% concatenated from Arithmetic-Bohr-1.tex
\documentclass[a4paper,11pt, reqno]{amsart}

\usepackage{amsmath}
\usepackage{amssymb}
\usepackage{mathrsfs}
\usepackage{ifthen}
\usepackage{graphicx}
\usepackage[T1]{fontenc} %skandit
\usepackage{lipsum}
\usepackage[a4paper,margin=0.75in]{geometry}

%\setlength{\topmargin}{-0.02in}
%\setlength{\topskip}{0.1in} % between header and text
%\setlength{\textheight}{9.3 in} % height of main text
%\setlength{\textwidth}{5.65in} % width of text
%\setlength{\oddsidemargin}{0.20in} % odd page left margin
%\setlength{\evensidemargin}{0.30in} % even page left margin
%\addtolength{\evensidemargin}{4cm} \addtolength{\oddsidemargin}{-0.9cm} \addtolength{\textwidth}{2cm}
\usepackage{cite}

%%%%%%%%%%%%%%%%%%%%%%%%%%%%%%%%%%%%%%%%%%%%%%%%%%%%%%%%%%%%%%%%%%


%%%%%%%%%%%%
%\usepackage{lineno}

%\linenumbers

%\usepackage[running]{lineno}
%\linenumbers

%\internallinenumbers
%%%%%%%%%%%%%

%%%%%%%%%%%%%%%%%%%%%%%%%%%%%%%%%%%%%%%%%%%%%%%%%%%%%


\def\switchlinenumbers{\@ifstar
	{\let\makeLineNumberOdd\makeLineNumberRight
		\let\makeLineNumberEven\makeLineNumberLeft}%
	{\let\makeLineNumberOdd\makeLineNumberLeft
		\let\makeLineNumberEven\makeLineNumberRight}%
}

\def\setmakelinenumbers#1{\@ifstar
	{\let\makeLineNumberRunning#1%
		\let\makeLineNumberOdd#1%
		\let\makeLineNumberEven#1}%
	{\ifx\c@linenumber\c@runninglinenumber
		\let\makeLineNumberRunning#1%
		\else
		\let\makeLineNumberOdd#1%
		\let\makeLineNumberEven#1%
		\fi}%
}

\def\leftlinenumbers{\setmakelinenumbers\makeLineNumberLeft}
\def\rightlinenumbers{\setmakelinenumbers\makeLineNumberRight}




%%%%%%%%%%%%%%%%%%%%%%%%%%%%%%%%%%%%%%%%%%%%%%%%%%%%%%

%\usepackage{geometry}
%\geometry{a4paper,textwidth=16cm,textheight=25cm,left=3cm}
\nonstopmode \numberwithin{equation}{section}
\setlength{\textwidth}{16.2cm} \setlength{\oddsidemargin}{0cm}
\setlength{\evensidemargin}{0cm} \setlength{\footskip}{30pt}
\pagestyle{plain}

\newtheorem*{theorem*}{Theorem}

%\theoremstyle{plain}
\newtheorem{thm}{Theorem}[section]
%\newtheorem{thm}[equation]{Theorem}
\newtheorem{cor}{Corollary}[section]
\newtheorem{lem}{Lemma}[section]
\newtheorem{prop}[equation]{Proposition}

\newtheorem{Cor}{Corollary}
\renewcommand{\theCor}{}
\newtheorem{claim}{Claim}
\newtheorem{conj}{Conjecture}

\theoremstyle{definition}
\newtheorem{defn}{Definition}[section]
\newtheorem{example}{Example}[section]
\newtheorem{qsn}{Question} [section]
\newtheorem{prob}[equation]{Problem}
\newtheorem{rem}{Remark}[section]

\newtheorem{innercustomthm}{Theorem}
\newenvironment{customthm}[1]
{\renewcommand\theinnercustomthm{#1}\innercustomthm}
{\endinnercustomthm}

%\newenvironment{rem}{%
	%\bigskip
	%\noindent \textsl{{\sl Remark. }}}{\bigskip}
%\newenvironment{rems}{%
	%\bigskip
	%\noindent \textsl{{\sl Remarks. }}}{\bigskip}


%%%%%%%%%%%% METHOD FOR HOUR AND MINUTE %%%%%%%%%%%%%
\newcounter{minutes}\setcounter{minutes}{\time}
\divide\time by 60
\newcounter{hours}\setcounter{hours}{\time}
\multiply\time by 60
\addtocounter{minutes}{-\time}
%%%%%%%%%%%%%%%%%%%%%%%%%%%%%%%%%%%%%%%%%%%%%%%%%%%%%

\newcounter {own}
\def\theown {\thesection       .\arabic{own}}


\newenvironment{nonsec}{\bf
	\setcounter{own}{\value{equation}}\addtocounter{equation}{1}
	\refstepcounter{own}
	\bigskip
	
	\indent \theown.$\ \,$}{$\,$.\ \ \ }


\newenvironment{pf}[1][]{%
	\vskip 3mm
	\noindent
	\ifthenelse{\equal{#1}{}}%
	{{\slshape {\bf Proof}. }}%
	{{\slshape #1.} }%
}%
{\qed\bigskip}


\DeclareMathOperator*{\Liminf}{\varliminf}
\DeclareMathOperator*{\Limsup}{\varlimsup}

\newcounter{alphabet}
\newcounter{tmp}
\newenvironment{Thm}[1][]{\refstepcounter{alphabet}%
	\bigskip%
	\noindent%
	{\bf Theorem \Alph{alphabet}}%
	\ifthenelse{\equal{#1}{}}{}{ (#1)}%
	{\bf .} \itshape}{\vskip 8pt}
%\newcommand{\Ref}[1]{\setcounter{tmp}{\ref{#1}}\Alph{tmp}}


%%%%%%%%%%%%%%%%%%%%%%%%%%%%%%%%%%%%%%%%%%%%%%%%%%%%%%%%%%%%%%%%%%%%%%%%%%%%%%%%%5


%\newcommand{\pad}[2]{\frac{\der #1}{\der #2}}
\def\be{\begin{equation}}
	\def\ee{\end{equation}}
\newcommand{\sep}{\itemsep -0.01in}
\newcommand{\seps}{\itemsep -0.02in}
\newcommand{\sepss}{\itemsep -0.03in}
\newcommand{\bee}{\begin{enumerate}}
	\newcommand{\eee}{\end{enumerate}}
\newcommand{\pays}{\!\!\!\!}
\newcommand{\pay}{\!\!\!}
\newcommand{\blem}{\begin{lem}}
	\newcommand{\elem}{\end{lem}}
\newcommand{\bthm}{\begin{thm}}
	\newcommand{\ethm}{\end{thm}}
\newcommand{\bcor}{\begin{cor}}
	\newcommand{\ecor}{\end{cor}}
\newcommand{\beg}{\begin{examp}}
	\newcommand{\eeg}{\end{examp}}
\newcommand{\begs}{\begin{examples}}
	\newcommand{\eegs}{\end{examples}}
\newcommand{\bdefe}{\begin{defin}}
	\newcommand{\edefe}{\end{defin}}
\newcommand{\bprob}{\begin{prob}}
	\newcommand{\eprob}{\end{prob}}
\newcommand{\bei}{\begin{itemize}}
	\newcommand{\eei}{\end{itemize}}
\newcommand{\real}{{\operatorname{Re}\,}}
\newcommand{\imaginary}{{\operatorname{Im}\,}}
\newcommand{\conjugate}{{\operatorname{^*}}}
\newcommand{\norm}[1]{\left\lVert#1\right\rVert}
\newcommand{\abs}[1]{\left\lvert#1\right\rvert}
\newcommand{\innpdct}[1]{\left\langle#1\right\rangle}
\newcommand{\tr}{{\operatorname{tr}\,}}

\newcommand{\B}{\mathbb{B}}
\newcommand{\C}{\mathbb{C}}
\newcommand{\D}{\mathbb{D}}
\newcommand{\E}{\mathbb{E}}
\newcommand{\G}{\mathbb{G}}
\newcommand{\LB}{\mathbb{L}}
\newcommand{\R}{\mathbb{R}}
\newcommand{\T}{\mathbb{T}}
\newcommand{\ov}{\overline}
\newcommand{\comment}[1]{}

%%%%%%%%%%%%%%%%%%%%%%%%%%%%%%%%%%%%%%%%%%%%%%%%%%%%%%%%%%%%%%%%%%%%%%%%%%%%%%%%%%%%%%%%%%%%%%%%%%%%%%%%%%%%%%%%%%%

\subjclass[{AMS} Subject Classification:]{Primary 32A05, 31C10, 46B07;  Secondary 32Q02, 46E40}
\keywords{Pluriharmonic functions, Bohr phenomenon; complete Reinhardt domain, Minkowski space, Banach sequence space, symmetric and convex Banach lattice, $p$-summing operator}

\begin{document}
	
	\title[]{Arithmetic Bohr radius and Local Banach space theory
		%Bohr and arithmetic Bohr radii for vector valued pluriharmonic functions via local Banach space theory
	}
	
	
	
	
	\author{Himadri Halder}
	\address{Himadri Halder,
		Department of Mathematics,
		Indian Institute of Science, Bangalore-560012, India}
	\email{himadrihalder119@gmail.com, himadrih@iisc.ac.in}
	
	%	\author{Sourav Pal}
	%	\address{Sourav Pal,
		%		Department of Mathematics,
		%		Indian Institute of Technology Bombay, Powai, Mumbai,
		%		Mumbai 400076, India.}
	%	\email{sourav@math.iitb.ac.in}
	
	
	
	
	
	%\def\thefootnote{}
	%\footnotetext{ {\tiny File:~\jobname.tex,
			%		printed: \number\year-\number\month-\number\day,
			%		\thehours.\ifnum\theminutes<10{0}\fi\theminutes }
		%} \makeatletter\def\thefootnote{\@arabic\c@footnote}\makeatother
	
	%	\pagestyle{abstract} 
	%	\thispagestyle{empty}
	
	\begin{abstract}
		This article introduces the notion of arithmetic Bohr radius for operator valued pluriharmonic functions on complete Reinhardt domains in $\mathbb{C}^n$. Using tools from local Banach space theory, we determine its asymptotic behavior in both finite and infinite dimensions. Asymptotic estimates for this constant are derived for both convex and non-convex complete Reinhardt domains. The framework developed in this article extends the classical Minkowski-space setting to a much broader class of sequence spaces, such as mixed Minkowski, Lorentz, and Orlicz spaces. Our results also apply to a wide class of Banach sequence spaces, including symmetric and convex Banach spaces. This generality allows for a unified and systematic investigation of Bohr’s theorem for both holomorphic and pluriharmonic functions. As an application of our results, we obtain several consequences extending known results in the scalar valued setting and in the existing literature.
	\end{abstract}
	
	
	
	\maketitle
	\pagestyle{myheadings}
	\markboth{Himadri Halder}{Bohr and arithmetic Bohr radii via local Banach space theory}
	
	\section{Introduction and the main results}\label{section-1}
	Over the past three decades, Bohr's theorem and its applications have become a central topic in the function theory of several complex variables, exhibiting deep connections with fundamental aspects of both the local and global theory of Banach spaces. These rich connections have stimulated extensive research on Bohr's theorem in various fields of mathematics. For instance, it has been investigated in the contexts of Banach algebras and uniform algebras (see \cite{paulsen-2002,paulsen-2004}), complex manifolds (see \cite{aizn-2000b,aizenberg-2001a}), ordinary and vector-valued Dirichlet series (see \cite{bala-2006,defant-2008}), elliptic equations (see \cite{aizenberg-2001b}), Faber–Green condensers (see \cite{lassere-2017}), free holomorphic functions (see \cite{popescu-2019}), vector-valued holomorphic functions (see \cite{defant-2011}), local Banach space theory (see \cite{defant-2003,defant-2025-TAMS}), domains of monomial convergence (see \cite{defant-2009}), harmonic and pluriharmonic mappings (see \cite{hamada-JFA-2022}), slice regular functions (see \cite{Rocchetta-Slice-Bohr-Math-Nachr,Xu-quaternion-AMPA,Xu-quaternion-PRSESA}), Hardy spaces (see \cite{bene-2004}), as well as various multidimensional settings (see \cite{aizn-2000a,boas-1997,boas-2000,defant-2004,defant-2006}). 
	\par 
	The investigation of Bohr's theorem for slice regular functions has recently attracted considerable attention in hypercomplex analysis. The pioneering work of Rocchetta, Gentili, and Sarfatti \cite{Rocchetta-Slice-Bohr-Math-Nachr} in $2012$ established the Bohr's theorem for slice regular functions, providing a fundamental extension of the classical Bohr phenomenon to the setting of noncommutative function theory. Their results opened new directions in the study of analytic properties of quaternionic and related function spaces. Subsequently, Xu \cite{Xu-quaternion-AMPA} in $2021$  made significant contributions by introducing the notion of the generalized Bohr radius for slice regular functions over quaternions, thereby refining and extending earlier results. In the same year, Xu \cite{Xu-quaternion-PRSESA} further extended the Bohr's theorem to slice regular functions over octonions, addressing additional challenges arising from nonassociativity. These developments have considerably enriched the theory and stimulated further research on Bohr's theorem in hypercomplex frameworks.
	\par In multidimensional settings, the Bohr radius problem becomes significantly more complicated. Many questions concerning the precise behavior of multidimensional Bohr radii remain open, largely due to the lack of certain essential analytical and geometric tools. Several of these problems are closely connected to fundamental topics in functional analysis, such as the geometry of Banach spaces \cite{Blasco-Collect-2017,defant-2012}, unconditional basis constants of spaces of polynomials \cite{aizenberg-2001a}, as well as Gordon–Lewis constants, projection constants, and the Banach–Mazur distance \cite{defant-2003}.
	 An especially intriguing direction arises when one considers this problem for vector valued functions on complete Reinhardt domains.
	\par	
	Our goal of this article is to build a local Banach space theoretic framework for the study of Bohr's theorem for vector valued holomorphic and pluriharmonic functions via the arithmetic Bohr radius. One of the core challenges in this study is to analyze the asymptotic behavior of both the Bohr radius and the arithmetic Bohr radius using tools and invariants from local Banach space theory. To proceed in more detail, we first recall the notion of the Bohr radius for vector valued holomorphic functions.
	%The framework and unified approach developed here provide a powerful tool for studying both vector valued holomorphic and pluriharmonic functions in broad classes of Banach sequence spaces, representing one of the most significant applications of our results in functional and complex analysis. 
	%Over the last three decades, Bohr's theorem and its applications have occupied a central position in geometric function theory, connecting ideas from several complex variables, functional analysis, and operator theory. 
	Let $\Omega \subset \mathbb{C}^n$ be a given complete Reinhardt domain and $n \in \mathbb{N}$. Let $U:X\rightarrow Y$ be a non-null bounded liner operator between two complex Banach spaces $X$ and $Y$, with $\norm{U} \leq \lambda$. For $1 \leq p < \infty$, the $\lambda$-powered Bohr radius of $U$, denoted by $K_{\lambda}(\Omega, p,U)$, is defined to be the supremum of all $r\geq 0$ such that for all holomorphic functions $f(z)=\sum_{\alpha}a_{\alpha}z^{\alpha}:\Omega \rightarrow X$ we have 
	\begin{equation} \label{e-1.2}
		\sup_{z \in r\Omega}\, \sum_{\alpha} \norm{U(a_{\alpha})z^\alpha}^p_{Y} \leq \lambda^p\,\norm{f}^p_{\Omega,X}, 
	\end{equation}
	where $\norm{f}_{\Omega,X}:=\sup_{z \in \Omega}\norm{f(z)}_{X}$. 
	\vspace{2mm}
	
	Let us fix some notations. We write $K(\Omega, p,U):=K_{1}(\Omega, p,U)$. When $U=I:X\rightarrow X$ is the identity operator, we set $K_{\lambda}(\Omega, p,X):=K_{\lambda}(\Omega, p,I)$ and $K(\Omega, p,X):=K_{1}(\Omega, p,X)$. In the scalar valued case, we denote $K_{\lambda}(\Omega, p):=K_{\lambda}(\Omega, p,\mathbb{C})$ and $K(\Omega,p):=K_1(\Omega,p)$.
	\vspace{1mm}
	  
	The above definition is motivated by Defant {\it et al.} \cite{defant-2012}, extending their result for the case $p=1$ and $\Omega=\mathbb{D}^n$ to a more general setting, thus making it possible to the study of Bohr phenomenon for broader classes of domains and functions. 
	%the Bohr radius $K_{n}(\Omega)$ of $\Omega$ is defined to be the supremum of all $r \in [0,1]$ such that the following inequality holds
%	\begin{equation} \label{e-hh-mult-BI}
%		\sup_{z \in r\Omega}\, \sum_{\alpha}|c_{\alpha}\,z^{\alpha}| \leq \sup_{z \in \Omega}\, \left|f(z)\right|
%	\end{equation}
%for each holomorphic function $f(z)=\sum_{\alpha}c_{\alpha}\,z^{\alpha}:\Omega\rightarrow \mathbb{C}$.
	With the above notations, a classical theorem of Harald Bohr concerning power series states that $K(\mathbb{D},1)=1/3$ (see \cite{Bohr-1914}). %One of the most striking aspects of Bohr's theorem arises in multidimensional settings, where the situation becomes significantly more complex and many problems remain open. 
	One of the most striking features of Bohr's theorem is that our understanding of the behavior of the constant $K(\Omega,1)$ remains rather limited. A central challenge is to determine the precise value of $K(\Omega,1)$, which is still unknown for $n>1$, even in the classical setting of the polydisc $\mathbb{D}^n$, where $\mathbb{D}^n:=\{z=(z_{1},\ldots,z_{n}) \in \mathbb{C}^n: |z_{j}|<1\, \mbox{for all}\, 1 \leq j \leq n\}$. To make progress in this direction, much of the existing research has focused on establishing bounds for $K(\Omega,1)$ and investigating its asymptotic behavior on general domains, and in particular well known settings such as the unit ball of Minkowski space and finite-dimensional complex Banach spaces. 
	\par
	Denote $B_{\ell^n _q}:=\left\{z=(z_{1},\ldots,z_{n}) \in \mathbb{C}^n:\norm{z}_{q}:=\left(\sum_{i=1}^{n}|z_{i}|^q\right)^{1/q}<1\right\}$, $1\leq q <\infty$ and $B_{\ell^n _\infty}:=\mathbb{D}^n$. For every $1\leq q \leq \infty$ and all $n\in \mathbb{N}$, together the results of Aizenberg, Boas, Dineen, Khavinson, Timoney, Defant, and Frerick from \cite{aizn-2000a,boas-1997,boas-2000,defant-2011-lpn} show that  there exist constants $C,D>0$ such that 
	\begin{equation*}
		\frac{1}{C}\, \left(\frac{\log\,n}{n}\right)^{1-\frac{1}{\min\{q,2\}}} \leq K(B_{\ell^n _q},1) \leq D \,\left(\frac{\log\,n}{n}\right)^{1-\frac{1}{\min\{q,2\}}}.
	\end{equation*}
Boas and Khavinson \cite{boas-1997} showed that the constant $K(\mathbb{D}^n,1)$ does not exceed $ 1/3$, and moreover that the sequence $\{K(\mathbb{D}^n,1)\}_{n=1}^{\infty}$ is decreasing with limit zero as $n \rightarrow \infty$. In the bidisc case, their work yields the bounds $0.23570226\leq K(\mathbb{D}^2,1) \leq 1/3=0.33333$. Considerable progress has been made only recently. Knese \cite{knese-2025} established in $2025$ that $K(\mathbb{D}^2,1)>0.3006$, while a subsequent result in $2026$ of Baran, Pikul, Woerdeman, and Wojtylak \cite{baran-2026} showed that $K(\mathbb{D}^2,1) < 0.3177$. These advances in particular ensure that $K(\mathbb{D}^n,1) <1/3$ for $n \geq 2$.
	\vspace{1mm}
	
    It is worth mentioning that the constant $K_{\lambda}(\Omega, p,U)$ is not necessarily nonzero for arbitrary Banach spaces $X,Y$, and the operator $U$. We refer to the occurrence of strictly positive constant $K_{\lambda}(\Omega, p,U)$ as the Bohr phenomenon. More precisely, we make the following definition. For given Banach spaces $X,Y$ and the operator $U$ as above, a class of holomorphic functions $f$, defined in a complete Reinhardt domain $\Omega$, is said to exhibit the Bohr phenomenon on $\Omega$ if there exists a universal constant $r=r_{0}\in (0,1]$, called the Bohr radius of $\Omega$ with respect to this class, such that inequality \eqref{e-1.2} is satisfied by every function in the class. Since the Bohr phenomenon does not occur for all classes of holomorphic functions (see \cite{aizn-2000b,Himadri-Vasu-PEMS,bene-2004,Blasco-Collect-2017}), an important problem is to identify and characterize those classes for which the Bohr phenomenon holds. Blasco \cite{Blasco-Collect-2017} initiated the systematic study of classes exhibiting the Bohr phenomenon by characterizing them in the classical one-dimensional setting $\Omega=\mathbb{D}$ with $U$ equal to the identity operator. Many thanks to Defant, Maestre, and Schwarting \cite{defant-2012}, who were the first to address this problem in multidimensional case $\Omega=\mathbb{D}^n$, $n>1$, with $p=1$ and for an arbitrary operator $U$. More recently, the problem for general bounded complete Reinhardt domains has been considered by the present author in \cite{Himadri-local-Banach-1}, where a framework based on local Banach space theory was developed.
    \vspace{2mm}
    
  %Among the diverse applications of the classical Bohr phenomenon, a particularly noteworthy development is its elegant connection with finite Blaschke products, as demonstrated by Garcia {\it et al.} \cite{Bohr-Blaschke-Book}.
		The constant $K_{\lambda}(\Omega, 1)$ was first introduced by Bombieri \cite{bombieri-1962}, who determined its exact value for $\lambda \in [1,\sqrt{2}]$. Its precise asymptotic behavior as $\lambda\rightarrow \infty$ was later investigated by Bombieri and Bourgain \cite{bombieri-2004}. Building upon these foundational results,  Defant {\it et al.} \cite{defant-2012} have introduced the more general constant $K_\lambda(\mathbb{D}^n,1,U)$ and obtained asymptotic estimates in both finite and infinite dimensional Banach spaces $X$. On the other hand, the constants $K(\mathbb{D},p)$ and $K(\mathbb{D}^n,p)$ were first studied by Djakov and Ramanujan \cite{Djakov & Ramanujan & J. Anal & 2000}, and subsequently developed by B\'{e}n\'{e}teau, Dahlner, Khavinson, Das \cite{bene-2004}. From the viewpoint of Banach space valued functions, Blasco \cite{Blasco-Collect-2017} has shown that $K(\mathbb{D},p,X)=0$ for $p \in [1,2)$, whereas  $K(\mathbb{D},p,X)>0$ precisely when $X$ is $p$-uniformly $\mathbb{C}$-convex for $2 \geq p < \infty$. 
	%Furthermore, in \cite{Blasco-OTAA-2010}, Blasco has proved that $K(\mathbb{D},1,B_{\ell^2 _q})=0$ for $1\leq q \leq \infty$. 
	These results show that $K(\mathbb{D},p,X)$ is not necessarily nonzero for all Banach spaces $X$. This observation inspired a refinement of the definition to guarantee non-vanishing behavior for arbitrary $X$, leading to the introduction of the parameter $\lambda$ in \cite[Definition 1.1]{defant-2012}, and hence in \eqref{e-1.2}.
	\vspace{2mm}
	
	One of the main aims of this note is to study the Bohr radius through the lens of the arithmetic Bohr radius within the framework of local Banach space theory, by establishing a connection between the asymptotic behaviors of these two radii. Among the various notions of Bohr radius, the arithmetic Bohr radius is more than a mere curious variant of the classical Bohr radius. It possesses rich structural properties; in particular, it plays an important role in describing the domain of existence of monomial expansions of bounded holomorphic functions on complete Reinhardt domains (see \cite{prengel-2005}). With this objective in mind, we first recall the following notion of the arithmetic Bohr radius, which is associated with the inequality \eqref{e-1.2}.
	\vspace{2mm}
	
	 Let $\Omega \subset \mathbb{C}^n$ be a given complete Reinhardt domain and $n \in \mathbb{N}$. Let $X$ be any complex Banach space. Let $\mathcal{F}(\Omega, X)$ be a set of $X$-valued holomorphic functions on $\Omega$. For any complex Banach spaces $X$ and $Y$, let $U:X\rightarrow Y$ be a bounded liner operator and $\norm{U} \leq \lambda$. For each $1\leq p<\infty$ and $\lambda\geq 1$, the $\lambda_{p}$-\textit{arithmetic Bohr radius} $A_{\lambda}(\mathcal{F}(\Omega, X), p, U)$ of $\Omega$ with respect to $\mathcal{F}(\Omega, X)$ is defined as \begin{equation*}  \sup \left\{\frac{1}{n}\sum_{i=1}^{n}r_i \,|\, r\in \mathbb{R}^{n}_{\geq 0},\, \mbox{for all}\, f\in \mathcal{F}(\Omega, X) : \sum_{\alpha} \norm{U(a_{\alpha})r^\alpha}^p_{Y} \leq \lambda^p\,\norm{f}^p_{\Omega,X}\right\}, 
	\end{equation*}
	where $\mathbb{R}^n_{\geq 0}=\{r=(r_1,.\,.\,.\,, r_n) \in \mathbb{R}^n: r_i\geq 0, 1\leq i\leq n\}.$ Let $H^{\infty}(\Omega,X)$ be the space of all bounded $X$-valued holomorphic functions on $\Omega$. We write $A_\lambda(\Omega,p,U)$ for $A_{\lambda}(H^{\infty}(\Omega,X),p,U)$ and $A(\Omega,p,U)$ for $A_1(\Omega,p,U)$. When $U=I:X\rightarrow X$ is the identity operator, we set $A_{\lambda}(\Omega, p,X):=A_{\lambda}(\Omega, p,I)$ and $A(\Omega, p,X):=A_{1}(\Omega, p,X)$. In the scalar valued case, we denote $A_{\lambda}(\Omega, p):=A_{\lambda}(\Omega, p,\mathbb{C})$ and $A(\Omega,p):=A_1(\Omega,p)$. 
	% and $A^k_\lambda(\Omega,p,U)$ for $A_\lambda(\mathcal{P}^k(\Omega,X),\lambda)$. 
The constant $A_{\lambda}(\Omega,1)$ was first considered by Defant {\it et al.} \cite{defant-2007} and is usually known as the classical arithmetic Bohr radius. Recently, the authors in \cite{	A-H-P-banach,Allu-Pal-arithnmetic} have studied the constant $A_\lambda(\Omega,p,U)$ and have investigated asymptotic estimates of it on the Minkowski spaces. 
\vspace{1mm}

In this paper, we investigate asymptotic estimates of the arithmetic Bohr radius on arbitrary bounded complete Reinhardt domains using techniques from local Banach space theory. In a recent work \cite{Himadri-local-Banach-1}, the author introduced a notion of Bohr radius for operator valued pluriharmonic functions defined on such domains and analyzed its asymptotic behavior using tools from local Banach space theory. Motivated by these developments, it is natural to ask whether the arithmetic Bohr radius can also be investigated in the setting of operator valued pluriharmonic functions. 

In order to address this problem, we begin by recalling the definitions of harmonic and pluriharmonic functions. A function $f:\Omega \rightarrow \mathbb{C}$ of class $C^2$ is called plurharmonic if its restriction to every complex line is harmonic. When $\Omega \subset\mathbb{C}^n$ is a simply connected domain in containing the origin, such a function is pluriharmonic precisely when it can be decomposed as $f=h+\overline{g}$, where $h,g$ are holomorphic on $\Omega$ with $g(0)=0$. This characterization admits a natural extension to the operator valued setting. In particular, if $\Omega$ is a simply connected complete Reinhardt domain, then a continuous function $f:\Omega \rightarrow \mathcal{B}(\mathcal{H})$ is pluriharmonic if and only it possesses a series representation of the form
\begin{equation} \label{e-1.3-a}
	f(z)=\sum_{m=0}^{\infty} \sum_{|\alpha|=m} a_{\alpha}\, z^{\alpha} + \sum_{m=1}^{\infty} \sum_{|\alpha|=m} b^{*}_{\alpha}\, \bar{z}^{\alpha},
\end{equation} 
where the associated functions $h(z)=\sum_{m=0}^{\infty} \sum_{\alpha} a_{\alpha}\, z^{\alpha}$ and $g(z)=\sum_{m=1}^{\infty} \sum_{\alpha} b_{\alpha}\, z^{\alpha}$ are $\mathcal{B}(\mathcal{H})$ valued holomorphic functions on $\Omega$, and $a_{\alpha}, b_{\alpha} \in \mathcal{B}(\mathcal{H})$. Here, $\mathcal{B}(\mathcal{H})$ denotes the Banach algebra of all bounded linear operators acting on a complex Hilbert space $\mathcal{H}$. For any $T \in \mathcal{B}(\mathcal{H})$, $\norm{T}$ will always denote the operator norm of $T$, and $T^{*}$ is the usual adjoint of $T$. Write $Re(T):=(T+T^*)/2$. Moreover, if $z=(z_1, \ldots, z_{n})$, then $\overline{z}$ denotes the componentwise complex conjugate $(\overline{z}_1, \ldots,\overline{z}_n)$. For such a domain $\Omega$ as above, we denote by $\mathcal{PH}(\Omega,X)$ the set of all bounded $X$-valued pluriharmonic functions on $\Omega$, where $X=\mathcal{B}(\mathcal{H})$. Boundedness is understood in the sense that $\norm{f}_{\Omega,X}:=\sup_{z \in \Omega}\,\norm{f(z)}_{X} < \infty$ for $f \in \mathcal{PH}(\Omega,X)$.
\vspace{2mm}

Let $Y$ be any complex Banach space. Let $\Omega\subset \mathbb{C}^n$ be a simply connected complete Reinhardt domain and $n\in \mathbb{N}$. Let $U:\mathcal{B}(\mathcal{H})\rightarrow Y$ be a bounded liner operator and $\norm{U} \leq \lambda$. For $1 \leq p < \infty$, the $\lambda$-powered Bohr radius of $U$, denoted by $R_{\lambda}(\Omega, p,U)$, is defined to be the supremum of all $r\geq 0$ such that for all $f \in \mathcal{PH}(\Omega,\mathcal{B}(\mathcal{H}))$ of the form \eqref{e-1.3-a} we have 
\begin{equation} \label{e-1.4-a}
	\sup_{z \in r\Omega}\,\sum_{m=0}^{\infty} \sum_{|\alpha|=m} (\norm{U(a_{\alpha})}^p_{Y} + \norm{U(b_{\alpha})}^p_{Y})|z^\alpha|^p \leq \lambda^p\,\norm{f}^p_{\Omega,\mathcal{B}(\mathcal{H})}. 
\end{equation}

\noindent Set $R(\Omega, p,U):=R_{1}(\Omega, p,U)$, $R_{\lambda}(\Omega, p,\mathcal{B}(\mathcal{H})):=R_{\lambda}(\Omega, p,U)$ whenever $U=I:\mathcal{B}(\mathcal{H})\rightarrow \mathcal{B}(\mathcal{H})$, $R(\Omega, p,\mathcal{B}(\mathcal{H})):=R_{1}(\Omega, p,\mathcal{B}(\mathcal{H}))$, $R_{\lambda}(\Omega, p):=R_{\lambda}(\Omega, p,\mathbb{C})$, and $R(\Omega,p):=R_1(\Omega,p)$. The Bohr phenomenon for harmonic functions in the scalar setting was first investigated in \cite{Abu-2010}. Subsequent developments extended this line of research to operator valued harmonic mappings (see \cite{bhowmik-2021}), and more recently to pluriharmonic functions taking values in Hilbert spaces (see \cite{hamada-Math-Nachr-2023}). The asymptotic behaviour of the constants $R(\Omega,1)$ and $R(\Omega,p)$ have been studied in \cite{hamada-JFA-2022,das-2024} whenever $\Omega=B_{\ell^n _q}$.
\vspace{2mm}

We now turn to the central theme of this article, namely the notion of the arithmetic Bohr radius for pluriharmonic functions. To the best of our knowledge, this concept has not been explored previously in the literature. This concept not been even studied in the simplest cases such as $U=I:\mathcal{B}(\mathcal{H})\rightarrow \mathcal{B}(\mathcal{H})$, $p=1$, $\lambda=1$, and for the classical domains such as the unit balls of Minkowski spaces, nor for the scalar valued harmonic and pluriharmonic functions. This gap motivates us to pose the following question. 
\begin{qsn} \label{qsn-1.2}
	Let $\Omega \subset \mathbb{C}^n$ be a simply connected complete Reinhardt domain. Can we study the arithmetic Bohr radius for operator valued pluriharmonic functions on $\Omega$?
\end{qsn}
We answer affirmatively to this question in this paper by introducing this notion in the following way.
Let $\mathcal{FP}(\Omega, \mathcal{B}(\mathcal{H}))$ be a set of $\mathcal{B}(\mathcal{H})$-valued pluriharmonic functions on $\Omega$. Let $U:\mathcal{B}(\mathcal{H})\rightarrow Y$ be a bounded liner operator and $\norm{U} \leq \lambda$.
	For each $1\leq p<\infty$ and $\lambda\geq 1$, the $\lambda_{p}$-\textit{arithmetic Bohr radius} $AP_{\lambda}(\mathcal{F}(\Omega, \mathcal{B}(\mathcal{H})), p, U)$ of $\Omega$ with respect to $\mathcal{FP}(\Omega, \mathcal{B}(\mathcal{H}))$ is defined as $\sup\left\{\frac{1}{n}\sum_{i=1}^{n}r_i:r\in \mathbb{R}^{n}_{\geq 0}\right\}$ such that
	\begin{equation*} 
\sum_{m=0}^{\infty} \sum_{|\alpha|=m} (\norm{U(a_{\alpha})}^p_{Y} + \norm{U(b_{\alpha})}^p_{Y})r^{p \alpha} \leq \lambda^p\norm{f}^p_{\Omega,\mathcal{B}(\mathcal{H})}
	\end{equation*}
for all $f\in \mathcal{FP}(\Omega, \mathcal{B}(\mathcal{H}))$ of the form \eqref{e-1.3-a}.
	Set $AP_\lambda(\Omega,p,U):=AP_\lambda(\mathcal{PH}(\Omega,\mathcal{B}(\mathcal{H})),p,U)$ and $AP(\Omega,P,U):=AP_1(\Omega,p,U)$. When $U=I:\mathcal{B}(\mathcal{H})\rightarrow \mathcal{B}(\mathcal{H})$ is the identity operator, we set $AP_{\lambda}(\Omega, p,\mathcal{B}(\mathcal{H})):=AP_{\lambda}(\Omega, p,I)$ and $AP(\Omega, p,X):=AP_{1}(\Omega, p,\mathcal{B}(\mathcal{H}))$. In the scalar valued case, we denote $AP_{\lambda}(\Omega, p):=AP_{\lambda}(\Omega, p,\mathbb{C})$ and $AP(\Omega,p):=AP_1(\Omega,p)$.  Let $AP^k_\lambda(\Omega,p,U):=AP_\lambda(\mathcal{PH}(^k\Omega,\mathcal{B}(\mathcal{H})),p,U)$. In the scalar valued case, we denote $AP^k_{\lambda}(\Omega, p):=AP^k_{\lambda}(\Omega, p,\mathbb{C})$ and $AP^k(\Omega,p):=A^k_1(\Omega,p)$. 
	
	
	
	Recently, the present author in \cite{Himadri-local-Banach-1} has shown that the constant $R_{\lambda}(\Omega, p,U)$ is not necessarily nonzero for arbitrary operator $U$. Consequently, in the same paper, the author has established the condition when this constant is non-zero. Moreover, the author investigated asymptotic behaviors of this constant using invariants from local Banach space theory. With these developments in mind, it is natural to ask the following question for the constant $AP_{\lambda}(\Omega, p,U)$.
	\begin{qsn} \label{qsn-1.1}
		For a given simply connected complete Reinhardt domain $\Omega \subset \mathbb{C}^n$, can the the constant $AP_{\lambda}(\Omega, p,U)$ be studied in detail? If so, does it always remain nonzero, and what can be said about its asymptotic behavior?
	\end{qsn}
	 %We now focus on complex-valued and operator-valued pluriharmonic functions defined on domains $\Omega$ in $\mathbb{C}^n$.  These advances naturally lead to the following question for operator-valued pluriharmonic functions.
	
%	In this paper, we provide an affirmative answer to this question, and to this end, we first introduce the notion of the powered Bohr radius for operator-valued pluriharmonic functions. 
	%For a simply connected complete Reinhardt domain $\Omega$ and $X=\mathcal{B}(\mathcal{H})$, let $\mathcal{PH}^{+}(\Omega,X)$ denote the set of all functions $f$ in $\mathcal{PH}(\Omega,X)$ such that $f(0) \geq 0$. 
	\comment{\begin{defn} \label{def-1.1}
			Let $X=\mathcal{B}(\mathcal{H})$ and $Y$ be any complex Banach space. Let $\Omega\subset \mathbb{C}^n$ be a simply connected complete Reinhardt domain and $n\in \mathbb{N}$. Let $U:X\rightarrow Y$ be a bounded liner operator and $\norm{U} \leq \lambda$. For $1 \leq p < \infty$, the $\lambda$-powered Bohr radius of $U$, denoted by $R_{\lambda}(\Omega, p,U)$, is defined to be the supremum of all $r\geq 0$ such that for all $f \in \mathcal{PH}(\Omega,X)$ of the form \eqref{e-1.3-a} we have 
			\begin{equation} \label{e-1.4-a}
				\sup_{z \in r\Omega}\,\sum_{m=0}^{\infty} \sum_{|\alpha|=m} (\norm{U(a_{\alpha})}^p_{Y} + \norm{U(b_{\alpha})}^p_{Y})|z^\alpha|^p \leq \lambda^p\,\norm{f}^p_{\Omega,X}. 
			\end{equation}
		\end{defn}
		\noindent Set $R(\Omega, p,U):=R_{1}(\Omega, p,U)$, $R_{\lambda}(\Omega, p,X):=R_{\lambda}(\Omega, p,U)$ whenever $U=I:X\rightarrow X$, $R(\Omega, p,X):=R_{1}(\Omega, p,X)$, $R_{\lambda}(\Omega, p):=R_{\lambda}(\Omega, p,\mathbb{C})$, and $R(\Omega,p):=R_1(\Omega,p)$. The asymptotic behaviour of the constants $R(\Omega,1)$ and $R(\Omega,p)$ have been studied in \cite{hamada-JFA-2022,das-2024} whenever $\Omega=B_{\ell^n _q}$.}
	
%	We now proceed to present our main results. To the best of our knowledge, no attempt has yet been made to study operator-valued analogues of Bohr’s theorem for pluriharmonic functions in the framework of Definition \ref{def-1.1}, 
%	\begin{qsn} \label{qsn-1.3}
%		For a given simply connected complete Reinhardt domain $\Omega$, can the the constant $R_{\lambda}(\Omega, p,U)$ be studied in detail? If so, does it always remain nonzero, and what can be said about its asymptotic behavior?
%	\end{qsn} 
	One of our main contributions to this paper is to affirmatively answer the aforesaid question for arbitrary bounded simply connected complete Reinhardt domain by establishing the relation between Bohr radius and arithmetic Bohr radius. The following theorem relates the Bohr radius constant $R_{\lambda}(\Omega, p,U)$ and the arithmetic Bohr radius $AP_{\lambda}(\Omega, p,U)$ on any bounded  simply connected complete Reinhardt domain, which as a consequence shows the strictly positivity of $AP_{\lambda}(\Omega, p,U)$. This theorem actually help to obtain the lower estimate of $AP_{\lambda}(\Omega, p,U)$. A part of the proof follows a similar argument to that of \cite[Lemma 4.3]{defant-2012}; however, in that work the authors restricted attention only to scalar valued holomorphic functions.
	Before stating our result, we introduce the following standard notation. For bounded Reinhardt domains $\Omega_{1}, \Omega_{2}\subset \mathbb{C}^n$, let $S(\Omega_1, \Omega_{2}):= \inf \left\{s>0 : \Omega_{1} \subset s \, \Omega_{2}\right\}$. Recall that if $Z$ and $W$ are Banach sequence spaces then $S(B_{Z},B_{W})=\norm{Id:Z \rightarrow W}$.
	\begin{thm} \label{thm-1.8-atm}
			Let $Z=(\mathbb{C}^n, ||.||)$ be a Banach space such that $\chi(\{e_k\}^n_{k=1})=1$. Then 
		\begin{equation} \label{e-1.1-added}
			AP_\lambda(B_{Z},p,U) \geq \frac{\norm{Id:Z \rightarrow \ell^n_1}}{n} \, R_\lambda(B_{Z},p,U).
		\end{equation}
	Moreover, if $\norm{U}<\lambda$, then $AP_\lambda(B_{Z},p,U) \geq D.\frac{\norm{Id:Z \rightarrow \ell^n_1}}{n} \,. \frac{1}{\sup_{z \in B_{Z}}\norm{z}_{p}}$, where 
	$$
	D=\begin{cases}
		\max \left\{\frac{1}{4 \lambda\,2^{\frac{1}{p}}}\left(\frac{\lambda^p - \norm{U}^p}{2\lambda^p - \norm{U}^p}\right)^{\frac{1}{p}}\, , \frac{1}{4 \lambda\,2^{\frac{1}{p}}}\left(\frac{\lambda^p - \norm{U}^p}{\lambda^p - \norm{U}^p +1}\right)^{\frac{1}{p}}\, \frac{1}{\norm{U}}\right\}\,  & \text{for $\norm{U}\geq \frac{1}{4 \lambda\,2^{\frac{1}{p}}}$},\\[3mm]
		\max \left\{\frac{1}{4 \lambda\,2^{\frac{1}{p}}}\left(\frac{\lambda^p - \norm{U}^p}{2\lambda^p - \norm{U}^p}\right)^{\frac{1}{p}}, \frac{1}{4 \lambda\,2^{\frac{1}{p}}}\left(\frac{\lambda^p - \norm{U}^p}{\lambda^p - \norm{U}^p +1}\right)^{\frac{1}{p}}\right\} & \text{for $0<\norm{U}< \frac{1}{4 \lambda\,2^{\frac{1}{p}}}$}.
	\end{cases}
	$$
	\end{thm}
Observing that $S(\Omega,B_{\ell^n_1})=\sup_{z \in \Omega} \norm{z}_{1}$, and following arguments similar to those used in the proof of Theorem~\ref{thm-1.8-atm}, it is straightforward to verify that the same conclusion holds for any simply connected complete Reinhardt domain. We therefore state the result and omit the details to avoid repetition.
\begin{cor} \label{cor-1.1-added}
Let $\Omega$ be a simply connected complete Reinhardt domain in $\mathbb{C}^n$. Then 
\begin{equation*}
	AP_\lambda(\Omega,p,U) \geq \frac{S(\Omega,B_{\ell^n_1})}{n} \, R_\lambda(\Omega,p,U).
\end{equation*}
Moreover, if $\norm{U}<\lambda$, then $AP_\lambda(\Omega,p,U) \geq D.\frac{S(\Omega,B_{\ell^n_1})}{n} \,. \frac{1}{\sup_{z \in \Omega}\norm{z}_{p}}$, where $D$ as in Theorem \ref{thm-1.8-atm}.
\end{cor}
Applying this corollary to the case $\Omega=B_{\ell^n_q}$, the Minkowski space, yields the following result. 
\begin{cor}
Let $Y$ be any complex Banach space and $U:\mathcal{B}(\mathcal{H})\rightarrow Y$ a bounded liner operator with $\norm{U} \leq \lambda$. For each $1\leq p<\infty$ and $1\leq q \leq\infty$,
\begin{equation*}
AP_\lambda(B_{\ell^n_q},p,U) \geq D\,
\begin{cases}	
	n^{- \frac{1}{q}} & \text{for $q\leq p$}, \\[2mm]
	 n^{ - \frac{1}{p}} & \text{for $q>p$},
\end{cases}
\end{equation*}
where $D$ as in Theorem \ref{thm-1.8-atm}.
\end{cor}
	In particular, we obtain the following lower bound of the $\lambda$-powered arithmetic Bohr radius for operator-valued pluriharmonic functions whenever $U$ is the identity operator on $X=\mathcal{B}(\mathcal{H})$.
	\begin{cor}
		Let $X=\mathcal{B}(\mathcal{H})$ and $1<\lambda$. Then for all $n\in \mathbb{N}$ and $1\leq p < \infty$, we have
		\begin{equation*}
			AP_{\lambda}(\Omega, p,X) \geq \frac{1}{2^{2+\frac{1}{p}}} \frac{\left(\lambda^p -1\right)^{\frac{1}{p}}}{\lambda^2}\, \frac{S(\Omega, B_{\ell^n_1})}{n}\, \frac{1}{\sup_{z \in \Omega}\norm{z}_{p}}.
		\end{equation*}
	\end{cor}
	\begin{rem}
	In view of Theorem \ref{thm-1.8-atm} and Corollary \ref{cor-1.1-added}, it is evident that the constant $AP_\lambda(\Omega,p,U)$ is non-zero whenever $\norm{U}<\lambda$. Moreover, a direct comparison of the definitions of arithmetic Bohr radius and Bohr radius shows that $R_\lambda(\Omega,p,U)\leq AP_\lambda(\Omega,p,U)$. As a consequence, the vanishing of $AP_\lambda(\Omega,p,U)$ necessarily implies the vanishing of $R_\lambda(\Omega,p,U)$. However, the converse implication fails in general, reflecting a genuine distinction between these two notions. In the recent work, the author in \cite{Himadri-local-Banach-1} has shown that $R_\lambda(\Omega,p,U)=0$ when $\norm{U}=\lambda$. This result naturally leads to the question of whether the arithmetic Bohr radius exhibits the same behavior, namely, whether $AP_\lambda(\Omega,p,U)$ also vanishes when $\norm{U}=\lambda$. The example presented below shows that this is not the case and the constant $AP_{\lambda}(\Omega, p,U)$ may be strictly positive. 
		For a bounded simply connected complete Reinhardt domain $\Omega \subseteq \mathbb{C}^n$ and $k \in \mathbb{N}$, consider the functions  
		$F_{k}: \Omega \rightarrow \mathcal{B}(\mathcal{H})$ by $$F_{k}(z)=(i\, \cos \frac{1}{k}) I_{\mathcal{H}} + (\frac{1}{2} \sin \frac{1}{k})  I_{\mathcal{H}}z_{1}+(\frac{1}{2} \sin \frac{1}{k})  I_{\mathcal{H}} \overline{z_{1}}$$
		for $z=(z_{1}, \ldots,z_{n}) \in \Omega$, where $I_{\mathcal{H}}$ is the identity on $\mathcal{H}$. Let $U=\lambda \, I$, with $I$ the identity on $\mathcal{B}(\mathcal{H})$. Then $F_{k}(0)=(i\, \cos \frac{1}{k})  I_{\mathcal{H}}$. Clearly, $F_{k} \in \mathcal{PH}(\Omega,\mathcal{B}(\mathcal{H}))$ and $\norm{F_{k}}_{\Omega,\mathcal{B}(\mathcal{H})} \leq 1$. Let if possible, there exists $r_{0}>0$ such that \eqref{e-1.4-a} holds for all $f \in \mathcal{PH}(\Omega,\mathcal{B}(\mathcal{H}))$ for all $z \in r_{0}\, \Omega$. Applying this to the functions $F_{k}$ yields, 
		\begin{equation} \label{e-example-inequality}
		\lambda^p \, |\cos \frac{1}{k}|^p + \lambda^p \, |\sin \frac{1}{k}|^p\, |z_{1}|^p \leq \lambda^p
		\end{equation}
		all $z \in r_{0}\, \Omega$ and for all $k \in \mathbb{N}$. However, this inequality fails for sufficiently large $k$, since $\cos (1/k) \rightarrow 1$ and $\sin(1/k) \rightarrow 0$ as $k \rightarrow\infty$, while $z$
		can be chosen in $r_{0}\, \Omega$ such that $z_{1} \neq 0$. This fact shows that there is no $r_{0}>0$ such that \eqref{e-example-inequality} holds, and hence $R_\lambda(\Omega,p,U)=0$. On the other hand, it is easy to check that \eqref{e-example-inequality} holds for all the the points $(0,\xi_2, \ldots, \xi_n) \in \Omega$ such that $\xi_j \geq 0$, $j=1,\ldots,n$. Thus, by the definition, $AP_\lambda(\Omega,p,U)>0$. Hence, we conclude that $R_\lambda(\Omega,p,U)=0$ while $AP_\lambda(\Omega,p,U)>0$ whenever $U=\lambda\, I$, thereby highlighting a fundamental difference between the arithmetic Bohr radius and the Bohr radius.
	\end{rem}
The following result links between the arithmetic Bohr radii of the spaces $\mathcal{PH}(\Omega,\mathcal{B}(\mathcal{H}))$ and $\cup^{\infty}_{m=1} \mathcal{PH}(^m \Omega,\mathcal{B}(\mathcal{H}))$. More precisely, we estimate $AP_{\lambda}(\Omega, p,U)$ in terms of $AP_{\lambda}\big(\cup^{\infty}_{m=1} \mathcal{PH}(^m \Omega,X), p,U\big)$. 
\begin{thm} \label{tthm-1.8-atm}
	Let $\Omega$ be a simply connected complete Reinhardt domain in $\mathbb{C}^n$. Let $X=\mathcal{B}(\mathcal{H})$ and $Y$ any complex Banach space, and $U: X \rightarrow Y$ be a non-null bounded linear operator such that $\norm{U}< \lambda$. Then, for all $p\in[1,\infty)$, $\lambda>1$, and $n \in \mathbb{N}$, we have\\
	$(1)$ $\left(\frac{\lambda^p - \norm{U}^p}{2\lambda^p - \norm{U}^p}\right)^{\frac{1}{p}}\, AP_{\lambda}\big(\cup^{\infty}_{m=1} \mathcal{PH}(^m \Omega,X), p,U\big) \leq AP_{\lambda}(\Omega, p,U) \leq AP_{\lambda}\big(\cup^{\infty}_{m=1} \mathcal{PH}(^m \Omega,X), p,U\big)$ \\
	$(2)$ $\left(\frac{\lambda^p - \norm{U}^p}{\lambda^p - \norm{U}^p +1}\right)^{\frac{1}{p}}\, AP\big(\cup^{\infty}_{m=1} \mathcal{PH}(^m \Omega,X), p,U\big) \leq AP_{\lambda}(\Omega, p,U) \leq \lambda \, \inf_{m\in \mathbb{N}} AP\big( \mathcal{PH}(^m \Omega,X), p,U\big)$.
\end{thm}

We now impose additional conditions on the initial coefficient, which lead to improve estimates for the arithmetic Bohr radius obtained in Theorem \ref{tthm-1.8-atm}. More precisely, we consider the following class of functions. Let $\mathcal{PH}^{*}(\Omega,\mathcal{B}(\mathcal{H})):=\{f \in \mathcal{PH}(\Omega,\mathcal{B}(\mathcal{H})): f(0)=\mu\, I, \, \mu \in [0, \norm{f}_{\Omega,\mathcal{B}(\mathcal{H})}]\}$. Denote $AP^{*}_{\lambda}(\Omega, p,U):= AP_{\lambda}(\mathcal{PH}^{*}(\Omega,\mathcal{B}(\mathcal{H})), p,U)$.
\begin{thm} \label{tthm-1.9-atm}
	Let $X=\mathcal{B}(\mathcal{H})$ and $Y$ any complex Banach space, and $U: X \rightarrow Y$ be a non-null bounded linear operator such that $\norm{U}< \lambda$. Then, for all $p\in[1,\infty)$, $\lambda>1$, and $n \in \mathbb{N}$, we have\\
	$(1)$ $C\, AP_{\lambda}\big(\cup^{\infty}_{m=1} \mathcal{PH}(^m \Omega,X), p,U\big) \leq AP^{*}_{\lambda}(\Omega, p,U) \leq AP_{\lambda}\big(\cup^{\infty}_{m=1} \mathcal{PH}(^m \Omega,X), p,U\big)$ \\
	$(2)$ $D\, AP\big(\cup^{\infty}_{m=1} \mathcal{PH}(^m \Omega,X), p,U\big) \leq AP^{*}_{\lambda}(\Omega, p,U) \leq \lambda \, \inf_{m\in \mathbb{N}} AP\big( \mathcal{PH}(^m \Omega,X), p,U\big)$.
	Here,
	$
	C=
	\max \left\{\frac{1}{(1+8^p)^{1/p}}, \left(\frac{\lambda^p - \norm{U}^p}{2\lambda^p - \norm{U}}\right)^{\frac{1}{p}}\right\}
	$
	and
	$D=
	\max \left\{\frac{\lambda}{(\lambda^p+8^p)^{1/p}}, \left(\frac{\lambda^p - \norm{U}^p}{\lambda^p - \norm{U}^p +1}\right)^{\frac{1}{p}}\right\}.
	$
\end{thm}
We now turn to the second part of Question \ref{qsn-1.1}, which is the central focus of this paper, namely the asymptotic estimate of the constant $AP_{\lambda}(\Omega, p,U)$ through invariants from local Banach space theory, such as unconditional bases and projection constants. 

In order to study this, we shall make use of the following correspondence.  Recall that a Schauder basis $\{w_{k}\}$ of a Banach space $W$ is said to be $1$-unconditional if $\chi(\{w_k\})=1$ (see more details in Section $3$). We usually denote the canonical basis vectors of $Z=(\mathbb{C}^n, ||.||)$ by $e_{k}$, $k=1, \ldots,n$. 
  Let $Z=(\mathbb{C}^n, ||.||)$ be an $n$-dimensional Banach space for which the canonical basis vectors $e_{k}$ form a normalized $1$-unconditional basis (or equivalently, the open unit ball $B_{Z}$ is a complete Reinhardt domain). There is a one-to-one correspondence between bounded convex complete Reinhardt domains $\Omega \subseteq\mathbb{C}^n$ and the open unit balls of norms on $\mathbb{C}^n$ for which the canonical basis vectors $e_{k}$ form a normalized $1$-unconditional basis. Indeed, the Minkowski functional $p_{\Omega}:\mathbb{C}^n \rightarrow \mathbb{R}_{+}$ of $\Omega$ defined by $$p_{\Omega}(z):= \inf \left\{t>0 : z/t \in \Omega\right\}$$
	is a norm on $\mathbb{C}^n$, and $\Omega$ coincides with the open unit ball of $(\mathbb{C}^n,p_{\Omega})$. Conversely, if $Z=(\mathbb{C}^n, ||.||)$ is a Banach space whose canonical basis is normalized $1$-unconditional, then its open unit ball $B_{	Z}$ is clearly a bounded convex complete Reinhardt domain in $\mathbb{C}^n$. Thus, the study of the constant $R_{\lambda}(\Omega, p,U)$ for the unit balls $\Omega=B_{Z}$ of finite dimensional complex Banach spaces $(\mathbb{C}^n, ||.||)$ with the normalized $1$-unconditional canonical bases, is equivalent to studying $R_{\lambda}(\Omega, p,U)$ over bounded convex complete Reinhardt domains $\Omega \subseteq\mathbb{C}^n$.
	
	Our investigation naturally separates into two cases according to whether $\mathcal{B}(\mathcal{H})$ is finite and infinite dimensional. In each case, taking into account the above correspondence, we estimate the asymptotic behavior of $AP_{\lambda}(\Omega, p,U)$ in the following steps. First, we establish the result for $\Omega=B_{Z}$, that is, for a convex complete Reinhardt domain. Next, by combining this estimate with the above correspondence and Lemma \ref{lem-3.6}, we derive the asymptotic behavior of $AP_{\lambda}(\Omega, p,U)$ for every bounded simply connected complete Reinhardt domain $\Omega \subseteq \mathbb{C}^n$, not necessarily convex.
	
	 We start with the finite-dimensional case. The following theorem determines the exact asymptotic behavior of the arithmetic Bohr radius $AP_\lambda(B_{Z},p,\mathcal{B}(\mathcal{H}))$ when $U$ is the identity operator on $\mathcal{B}(\mathcal{H})$.
	\begin{thm} \label{thm-1.2-atm}
		Let $Z=(\mathbb{C}^n, ||.||)$ be a Banach space such that $\chi(\{e_k\}^n_{k=1})=1$. Let $X=\mathcal{B}(\mathcal{H})$ be a finite dimensional complex Banach space and $\lambda>1$. Then 
		\comment{\begin{equation}
				R^{m}_{\lambda}(B_{Z}, p) \geq	\begin{cases}
					E_{1}(X)\, \lambda^{\frac{1}{m}} \, \max \left\{\frac{1}{\sqrt{n}\, \norm{Id: \ell^n_{2} \rightarrow Z}}, \, \frac{1}{\left(\frac{m^m}{m!}\right)^{1/m}\, \norm{Id: Z \rightarrow \ell^n_{1}}}\right\} & \text{for $p=1$}, \\[3mm]
					2^{-\frac{p+5}{pm}}\,  \lambda^{\frac{1}{m}}\, \norm{I:\ell^n _2 \rightarrow Z}^{-\frac{2}{p}} \,\norm{I-\real (a_{0})}^{-\frac{2}{pm}} & \text{for $p \geq 2$}, \\[2mm]
					
					E_{2}(X)\, \lambda^{\frac{1}{m}} \,\max \left\{\frac{1}{(\sqrt{n})^{1-\theta}\, \norm{Id: \ell^n_{2} \rightarrow Z}}, \, \frac{\norm{I:\ell^n _2 \rightarrow Z}^{-\theta}}{\left(\frac{m^m}{m!}\right)^{(1-\theta)/m}\, \norm{Id: Z \rightarrow \ell^n_{1}}^{1-\theta}}\right\} & \text{for $1<p< 2$}, 
				\end{cases}
			\end{equation}
			where $\theta=(2(p-1))/p$. Moreover,}
		\begin{equation*}
			AP_{\lambda}(B_{Z}, p,X) \geq	\begin{cases}
				E_{1}(X)\,  \max \left\{\frac{\norm{Id: Z \rightarrow \ell^n_{1}}}{n^{3/2}\, \norm{Id: \ell^n_{2} \rightarrow Z}}, \, \frac{1}{e\,n}\right\}\, \left(\frac{\lambda^p -1}{2\lambda^p - 1}\right)^{\frac{1}{p}} & \text{for $p=1$}, \\[2mm]
				E_{3} \frac{\norm{Id: Z \rightarrow \ell^n_{1}}}{n}\norm{I:\ell^n _2 \rightarrow Z}^{-\frac{2}{p}}\, \left(\frac{\lambda^p -1}{2\lambda^p - 1}\right)^{\frac{1}{p}}\, & \text{for $p \geq 2$}, \\[2mm]
				
				E_{2}(X)\,  \max \left\{\frac{\norm{Id: Z \rightarrow \ell^n_{1}}}{n^{\frac{3-\theta}{2}} \norm{Id: \ell^n_{2} \rightarrow Z}},  \frac{\norm{I:\ell^n _2 \rightarrow Z}^{-\theta}}{e^{(1-\theta)}n \norm{Id: Z \rightarrow \ell^n_{1}}^{-\theta}}\right\}\left(\frac{\lambda^p -1}{2\lambda^p - 1}\right)^{\frac{1}{p}} \hspace{-2mm} & \text{for $1<p< 2$}, 
			\end{cases}
		\end{equation*}
		%Here, $E_{1}(X)$, $E_{2}(X)$, and $E_{3}$ are positive constants depending respectively on $X$, on $X$ and $p$, and on $p$.
		where $\theta = (2(p-1))/p$. Here, $E_{1}(X)$, $E_{2}(X)$, and $E_{3}$ are positive constants: $E_{1}(X)$ depends only on $X$, $E_{2}(X)$ depends on both $X$ and $p$, while $E_{3}$ is independent of $X$ and depends only on $p$. Furthermore,
		for any $1\leq q \leq \infty$,
		\begin{equation*}
			AP_{\lambda}(B_{Z}, p,X) \leq d\,\, \lambda^{\frac{1}{\log\,n}}\, \norm{Id:Z \rightarrow \ell^n_q}\, (\log\,n)^{1-\frac{1}{\min\{2,q\}}}\,\,\left(\frac{1}{n}\right)^{\frac{1}{p}+\frac{1}{\max\{2,q\}}-\frac{1}{2}}
		\end{equation*}
		for some constant $d>0$.
	\end{thm}
In particular, for the Minkowski spaces $Z=\ell^n _q$, we obtain the following lower bound. The upper bound is obtained in the proof of Theorem \ref{thm-1.2-atm}, following the technique from \cite{defant-2007}.
\begin{cor} \label{cor-1.4}
		Let $X=\mathcal{B}(\mathcal{H})$ be finite dimensional and $\lambda>1$. Then, for $1 \leq q \leq \infty$,
	\begin{equation*}
		AP_{\lambda}(B_{\ell^n _q}, p,X) \geq	\begin{cases}
			E'_{3}(X)\,\left(\frac{\lambda^p -1}{2\lambda^p - 1}\right)^{\frac{1}{p}}\, n^{-\frac{1}{q}} \, \left(\frac{\log n}{n}\right)^{ \left(1- \frac{1}{\min\{q,2\}}\right)} & \text{for $p=1$}, \\[2mm]
			E'_{2} \left(\frac{\lambda^p -1}{2\lambda^p - 1}\right)^{\frac{1}{p}}\, n^{-\left(\frac{1}{p}+\frac{1}{q}\right)} & \text{for $p \geq q$}, \\[2mm]
			
			E_{4}(X)\,  \left(\frac{\lambda^p -1}{2\lambda^p - 1}\right)^{\frac{1}{p}} \,n^{-\frac{p(q-1)+q(p-1)}{pq(q-1)}}\, \left( \frac{\log n}{n}\right)^{\left(1- \frac{1}{\min\{q,2\}}\right)\, \frac{q-p}{p(q-1)}} & \text{for $1<p< q$}. 
		\end{cases}
	\end{equation*}
	Here, $E'_{3}(X)$, $E'_{2}$, and $E_{4}(X)$ are positive constants: $E'_{3}(X)$ depends only on $X$, $E_{4}(X)$ depends on both $X$ and $p$, while $E'_{2}$ is independent of $X$ and depends only on $p$
\end{cor}
Theorem \ref{thm-1.2-atm} leads to the following asymptotic estimates of the arithmetic Bohr radius for Banach sequence spaces, whose broad applications are described in Section $2$.
\begin{thm} \label{thm-1.3-a}
	Let $Z$ be a Banach sequence space, $X=\mathcal{B}(\mathcal{H})$ be finite dimensional, and $\lambda>1$. For $n\in \mathbb{N}$, let $Z_{n}$ be the linear span of $\{e_{k}, \, k=1,\ldots,n\}$. Then
	\begin{enumerate}
		\item if $Z$ is a subset of $\ell_2$ we have \begin{equation*}
			AP_{\lambda}(B_{Z}, p,X) \geq	\begin{cases}
				E_{1}(X)\,  \max \left\{\frac{\norm{\sum_{k=1}^{n}e^{*}_{k}}_{Z^{*}}}{n^{3/2}\, \norm{Id: \ell^n_{2} \rightarrow Z}}, \, \frac{1}{e\,n}\right\}\, \left(\frac{\lambda^p -1}{2\lambda^p - 1}\right)^{\frac{1}{p}} & \text{for $p=1$}, \\[2mm]
				E_{3} \frac{\norm{\sum_{k=1}^{n}e^{*}_{k}}_{Z^{*}}}{n}\norm{Id:\ell^n _2 \rightarrow Z}^{-\frac{2}{p}}\, \left(\frac{\lambda^p -1}{2\lambda^p - 1}\right)^{\frac{1}{p}}\, & \text{for $p \geq 2$}, \\[2mm]
				
				E_{2}(X)\,  \max \left\{\frac{\norm{\sum_{k=1}^{n}e^{*}_{k}}_{Z^{*}}}{n^{\frac{3-\theta}{2}} \norm{Id: \ell^n_{2} \rightarrow Z}},  \frac{\norm{Id:\ell^n _2 \rightarrow Z}^{-\theta}}{e^{(1-\theta)}n \norm{\sum_{k=1}^{n}e^{*}_{k}}_{Z^{*}}^{-\theta}}\right\}\left(\frac{\lambda^p -1}{2\lambda^p - 1}\right)^{\frac{1}{p}} \hspace{-2mm} & \text{for $1<p< 2$}, 
			\end{cases}
		\end{equation*}
		and 
		$$AP_{\lambda}(B_{Z_{n}}, p,X) \leq d\,  \lambda^{\frac{1}{\log\,n}}\,\, \sqrt{\log \, n} \, \left(\frac{1}{n}\right)^{\frac{1}{p}};$$
		\item if $Z$ is symmetric and $2$-convex, we have
		\begin{equation*}
			AP_{\lambda}(B_{Z}, p,X) \geq	\begin{cases}
				E_{1}(X)\,  \max \left\{\frac{\norm{\sum_{k=1}^{n}e^{*}_{k}}_{Z^{*}}}{n^{3/2}}, \, \frac{1}{e\,n}\right\}\, \left(\frac{\lambda^p -1}{2\lambda^p - 1}\right)^{\frac{1}{p}} & \text{for $p=1$}, \\[2mm]
				E_{3} \frac{\norm{\sum_{k=1}^{n}e^{*}_{k}}_{Z^{*}}}{n}\, \left(\frac{\lambda^p -1}{2\lambda^p - 1}\right)^{\frac{1}{p}}\, & \text{for $p \geq 2$}, \\[2mm]
				
				E_{2}(X)\,  \max \left\{\frac{\norm{\sum_{k=1}^{n}e^{*}_{k}}_{Z^{*}}}{n^{\frac{3-\theta}{2}}},  \frac{1}{e^{(1-\theta)}n \norm{\sum_{k=1}^{n}e^{*}_{k}}_{Z^{*}}^{-\theta}}\right\}\left(\frac{\lambda^p -1}{2\lambda^p - 1}\right)^{\frac{1}{p}} \hspace{-2mm} & \text{for $1<p< 2$}, 
			\end{cases}
		\end{equation*}
		and 
		$$AP_{\lambda}(B_{Z}, p,X) \leq d\,  \lambda^{\frac{1}{\log\,n}}\, \norm{\sum_{k=1}^{n}e^{*}_{k}}_{Z^{*}}\,\,\, \sqrt{\log \, n}\,\, \left(\frac{1}{n}\right)^{\frac{1}{p}+\frac{1}{2}}.
		$$
	\end{enumerate}
	The constants $E_{1}, E_{2}, E_{3}$, and $d$ are same as in Theorem \ref{thm-1.2-atm}.
\end{thm}
To determine the exact asymptotic behavior of the constant $AP_{\lambda}(B_{Z}, p,\mathcal{B}(\mathcal{H}))$ in the case where $\mathcal{B}(\mathcal{H})$ is infinite dimensional, we make essential use of Lemma \ref{lem-3.6}. This lemma allows us to reduce the problem to the study of this constant for such a domain in which the estimation of bounds becomes significantly more simpler. Among complete Reinhardt domains in $\mathbb{C}^n$, the most natural and widely studied examples are the unit balls $B_{\ell^n _q}$ of the Minkowski spaces $\ell^n _q$, $1\leq q \leq \infty$. Consequently, we begin by analyzing the behavior of $AP_{\lambda}(\Omega, p,\mathcal{B}(\mathcal{H}))$ when $\Omega=B_{\ell^n _q}$. 

In the finite-dimensional setting, Corollary \ref{cor-1.4} shows that the asymptotics of this constant necessarily involve a logarithmic factor. By contrast, the situation changes fundamentally in the infinite-dimensional case. As we show in the theorem below, the logarithmic term disappears entirely, and the asymptotic decay of $AP_{\lambda}(B_{\ell^n _q}, p,\mathcal{B}(\mathcal{H}))$ is governed instead by the geometric properties of the Banach space $X$, namely its optimal cotype $Cot(X)$(see Section $2$ for the definition), where $X=\mathcal{B}(\mathcal{H})$. 
	We now obtain the asymptotic estimates of the arithmetic Bohr radius for $B_{\ell^n _q}$ whenever $X=\mathcal{B}(\mathcal{H})$ is infinite dimensional.	
	\begin{thm} \label{thm-1.3-atm}
		Let $X=\mathcal{B}(\mathcal{H})$ be an infinite dimensional complex Banach space  and $\lambda>1$. Then $AP_{\lambda}(B_{\ell^n _q}, p,X) \geq \tilde{\Psi}_{1}(n,\lambda,p,q)$, where
		\begin{equation*}
			\tilde{\Psi}_{1}(n,\lambda,p,q):=	\begin{cases}
				%	\dfrac{\left(\lambda^p -1\right)^{\frac{1}{p}}}{\lambda}\, 2^{-\left(1+\frac{1}{p}\right)}\, e^{-\frac{1}{p}}\, 4^{-\frac{q}{p}} & \text{for $p > q$}, \\[2mm]
				
				E_{5}\,\dfrac{\left(\lambda^p -1\right)^{\frac{1}{p}}}{\lambda} \frac{1}{n^{1/q}} & \text{for $p>\min \{Cot(X),q\}$}, \\[2mm]
				E_{6}\, \dfrac{\left(\lambda^p -1\right)^{\frac{1}{p}}}{\lambda}\, n^{{\frac{1}{\min\{t,q\}}}-\frac{1}{p}-\frac{1}{q}} & \text{for $p\leq \min \{Cot(X),q\}$},
			\end{cases}
		\end{equation*}
		if $X$ has a finite cotype $t$. Here $E_{5}>0$ and $E_{6}>0$ are constant independent of $X$ and $n$. 
		On the other hand, we have
		\begin{equation*}
			AP_{\lambda}(B_{\ell^n _q}, p,X) \leq \tilde{\Psi}_{2}(n,\lambda,p,q):=	\begin{cases}	
				2^{1-\frac{1}{p}}\, \lambda\, n^{- \frac{1}{p}} & \text{for $q\leq Cot(X)$}, \\[2mm]
				2^{1-\frac{1}{p}}\,\lambda\, n^{\frac{1}{Cot(X)}\, - \,\frac{1}{p}-\frac{1}{q}} & \text{for $q>Cot(X)$}.
			\end{cases}
		\end{equation*}
	\end{thm}
	
	We are now in a stage to give the bounds of the arithmetic Bohr radius $AP_{\lambda}(B_{Z}, p,X)$ for the unit ball $B_{Z}$ of $Z=(\mathbb{C}^n,||.||)$ whenever $X=\mathcal{B}(\mathcal{H})$ is a infinite dimensional complex Banach space. In view of Theorem \ref{thm-1.3-atm} and Lemma \ref{lem-3.6}, we obtain the following result.
	\begin{thm} \label{thm-1.4-atm}
		Let $Z=(\mathbb{C}^n, ||.||)$ be a Banach space such that $\chi(\{e_k\}^n_{k=1})=1$. Let $\mathcal{B}(\mathcal{H})$ be infinite dimensional, and $\tilde{\Psi}_{1}(n,\lambda,p,q)$ and $\tilde{\Psi}_{2}(n,\lambda,p,q)$ be as in Theorem \ref{thm-1.3-atm}. For $\lambda>1$, 
		$$\frac{\tilde{\Psi}_{1}(n,\lambda,p,q)}{\norm{Id: \ell^n_{q} \rightarrow Z}} \, \leq AP_{\lambda}(B_{Z}, p,X) \leq \norm{Id: Z \rightarrow \ell^n_{q}} \tilde{\Psi}_{2}(n,\lambda,p,q).
		$$
	\end{thm}
The correspondence preceding Theorem \ref{thm-1.2-atm} allows Theorem \ref{thm-1.4-atm} to be extended to bounded simply connected convex complete Reinhardt domains. Moreover, combining Theorem \ref{thm-1.3-atm} with Lemma \ref{lem-3.6} yields the asymptotic estimates for all bounded simply connected complete Reinhardt domains, not necessarily convex. We record this result as a corollary.
\begin{cor} \label{cor-1.4-atm}
	Let $\Omega$ be a bounded simply connected complete Reinhardt domain. Let $\mathcal{B}(\mathcal{H})$ be infinite dimensional, and $\tilde{\Psi}_{1}(n,\lambda,p,q)$ and $\tilde{\Psi}_{2}(n,\lambda,p,q)$ be as in Theorem \ref{thm-1.3-atm}. For $\lambda>1$, 
	$$\frac{\tilde{\Psi}_{1}(n,\lambda,p,q)}{S(B_{\ell^n _q},\Omega)} \, \leq AP_{\lambda}(\Omega, p,X) \leq S(\Omega,B_{\ell^n _q}) \tilde{\Psi}_{2}(n,\lambda,p,q).
	$$
\end{cor}
An application of Theorem \ref{thm-1.4-atm}, together with the norm estimates for the identity operators mentioned in the proof of Theorem \ref{thm-1.3-a}, yields the following result.
	\begin{thm} \label{thm-1.6}
		Let $Z$ be a Banach sequence space and $X=\mathcal{B}(\mathcal{H})$ be infinite dimensional. For $n\in \mathbb{N}$, let $Z_{n}$ be the linear span of $\{e_{k}, \, k=1,\ldots,n\}$. Then for $\lambda>1$
		\begin{enumerate}
			\item if $Z$ is subset of $\ell_2$, we have 
			$$\frac{\tilde{\Psi}_{1}(n,\lambda,p,2)}{\norm{Id: \ell^n_{2} \rightarrow Z_n}} \, \leq AP_{\lambda}(B_{Z_n}, p,X) \leq  \tilde{\Psi}_{2}(n,\lambda,p,2);
			$$
			\item if $Z$ is symmetric and $2$-convex, we have $$\tilde{\Psi}_{1}(n,\lambda,p,2)\, \leq AP_{\lambda}(B_{Z}, p,X) \leq  \frac{\norm{\sum_{m=1}^{n}e^{*}_{m}}_{Z^{*}}}{\sqrt{n}}\, \tilde{\Psi}_{2}(n,\lambda,p,2).
			$$
		\end{enumerate}
	\end{thm}
	We have so far investigated the asymptotic behavior of the arithmetic Bohr radius associated with the identity operator $U:\mathcal{B}(\mathcal{H}) \rightarrow \mathcal{B}(\mathcal{H})$. We now turn to the study of this constant for a general operator $U$, which need not be the identity. To this end, we first analyze the Bohr radius constant corresponding to an arbitrary operator. Using this analysis, we then derive bounds for the arithmetic Bohr radius.
	\vspace{1mm}
	
	 With this objective in mind, we begin by establishing the following result, which may be viewed as an analogue of \cite[Theorem 5.3]{defant-2012}. While the authors there proved the result for the unit polydisc, our approach extends it to the unit ball of an arbitrary Banach space admitting a $1$-unconditional basis.
	\begin{thm} \label{thm-1.8}
	Let $Y$ be a $s$-concave Banach lattice, with $2\leq s < \infty$, and $U: X \rightarrow Y$ an $(t,1)$-summing operator with $1 \leq t \leq s$. Then there is a constant $C>0$ such that
	\begin{equation*}
	\left(\sum_{|\alpha|=m} \norm{U(a_{\alpha})z^{\alpha}}^{\beta}_{Y}\right)^{\frac{1}{\beta}} \leq C^m \, \norm{P}_{B_{Z},X}
	\end{equation*}
holds for every homogeneous polynomial $P(z)=\sum_{|\alpha|=m} a_{\alpha}\, z^{\alpha}$ on $Z$, where 
\begin{equation*}
\beta= \frac{mst}{s+ (m-1)t}.
\end{equation*}
	\end{thm}

We are now in a position to establish lower bounds for both the Bohr radius and the arithmetic Bohr radius for a general operator.
	\begin{thm} \label{thm-1.9}
	Let $1 \leq p < \infty$ and $1\leq q \leq \infty$.
	\begin{enumerate}
		\item Let $\mathcal{H}_1$ and $\mathcal{H}_2$ be two complex Hilbert spaces, and $U: \mathcal{B}(\mathcal{H}_1) \rightarrow \mathcal{B}(\mathcal{H}_2)$ be a non-null bounded linear operator with $\norm{U}<\lambda$. If $\mathcal{B}(\mathcal{H}_1)$ or $\mathcal{B}(\mathcal{H}_2)$ is of cotype $t$ with $2 \leq t \leq \infty$, then for every $n \in \mathbb{N}$,	
	\begin{align*}
	R_{\lambda}(B_{Z}, p,U) \geq \frac{\Phi_{1}(n,\lambda,p,q)}{\norm{Id:Z \rightarrow \ell^n _q} \norm{Id:\ell^n _q \rightarrow Z}}
	\end{align*} 
and
\begin{equation*}
	AP_{\lambda}(B_{Z}, p,U) \geq \frac{\Phi_{1}(n,\lambda,p,q)\, \norm{Id:Z \rightarrow \ell^n _1}}{n\norm{Id:Z \rightarrow \ell^n _q} \norm{Id:\ell^n _q \rightarrow Z}},
\end{equation*}
	where
	\begin{equation*}
		\Phi_{1}(n,\lambda,p,q):=	\begin{cases}
			%	\dfrac{\left(\lambda^p -1\right)^{\frac{1}{p}}}{\lambda}\, 2^{-\left(1+\frac{1}{p}\right)}\, e^{-\frac{1}{p}}\, 4^{-\frac{q}{p}} & \text{for $p > q$}, \\[2mm]
			
			E_{5}\,\dfrac{\left(\lambda^p -\norm{U}^p\right)^{\frac{1}{p}}}{\lambda} & \text{for $p>\min \{Cot(\mathcal{B}(\mathcal{H}_1)),Cot(\mathcal{B}(\mathcal{H}_2)),q\}$}, \\[2mm]
			E_{6}\, \dfrac{\left(\lambda^p -\norm{U}^p\right)^{\frac{1}{p}}}{\lambda}\,\, n^{{\frac{1}{\min\{t,q\}}}-\frac{1}{p}} & \text{for $p\leq \min \{Cot(\mathcal{B}(\mathcal{H}_1)),Cot(\mathcal{B}(\mathcal{H}_2)),q\}$}.
		\end{cases}
	\end{equation*}
	Here $E_{5}>0$ and $E_{6}>0$ are constants independent of $\mathcal{H}_1$, $\mathcal{H}_2$, and $n$.
	\item Let $Y$ be any complex Banach space and $U: \mathcal{B}(\mathcal{H}) \rightarrow Y$ be a non-null bounded linear operator with $\norm{U}<\lambda$. If $Y$ is a $s$-concave Banach lattice with $2 \leq s < \infty$ and there is a $1 \leq t <s$ such that $U$ is $(s,1)$-summing, then
	\begin{equation*}
		R_{\lambda}(B_Z,p,U) \geq D\, \left(\frac{\lambda^p - \norm{U}^p}{2\lambda^p - \norm{U}^p}\right)^{\frac{1}{p}} \, \left(\frac{\log \, n}{n}\right)^{1-\frac{1}{s}} 
	\end{equation*}
and \begin{equation*}
AP_{\lambda}(B_Z,p,U) \geq D\, \left(\frac{\lambda^p - \norm{U}^p}{2\lambda^p - \norm{U}^p}\right)^{\frac{1}{p}} \, \frac{\norm{Id:Z \rightarrow \ell^n _1}}{n}\,\left(\frac{\log \, n}{n}\right)^{1-\frac{1}{s}}
\end{equation*}
for some constant $D>0$ independent of $\mathcal{H}$ and $n$.
\end{enumerate}
	\end{thm}
\begin{rem}
\begin{enumerate}
	\item In particular, setting $q=1$ in Theorem \ref{thm-1.9} $(1)$ gives the following
		\begin{equation*}
		AP_{\lambda}(B_{Z}, p,U) \geq \frac{\Phi_{1}(n,\lambda,p,1)}{n \norm{Id:\ell^n _1 \rightarrow Z}}.
	\end{equation*}
	
	\item Since the lower bounds in $(1)$ of Theorem \ref{thm-1.9} are expressed in terms of the norms $\norm{Id:Z \rightarrow \ell^n _q}$ and $\norm{Id:\ell^n _q \rightarrow Z}$, varying $q$ over the interval $1 \leq q \leq \infty$ leads to an improvement of these estimates. The improved bounds are then given as follows:
	\begin{align*}
		R_{\lambda}(B_{Z}, p,U) \geq 
		\sup_{1\leq q \leq \infty} \left\{\frac{\Phi_{1}(n,\lambda,p,q)}{\norm{Id:Z \rightarrow \ell^n _q} \norm{Id:\ell^n _q \rightarrow Z}}\right\}
	\end{align*} 
	and
	\begin{equation*}
		AP_{\lambda}(B_{Z}, p,U) \geq 
		\sup_{1\leq q \leq \infty} \left\{\frac{\Phi_{1}(n,\lambda,p,q)\, \norm{Id:Z \rightarrow \ell^n _1}}{n\norm{Id:Z \rightarrow \ell^n _q} \norm{Id:\ell^n _q \rightarrow Z}}\right\}.
	\end{equation*}
\end{enumerate}
\end{rem}	
	 

Theorem \ref{thm-1.8-atm} provides a lower bound for the $n$-dimensional arithmetic Bohr radius $AP_\lambda(\Omega,p, U)$ in terms of the corresponding $n$-dimensional Bohr radius $R_\lambda(\Omega,p, U)$. For a given simply connected complete Reinhardt domain $\Omega \subset \mathbb{C}^n$, we now show that $AP_\lambda(\Omega,p, U)$ can be estimated in terms of the one-dimensional Bohr radius for $\mathbb{D}$, namely $R_\lambda(\mathbb{D},p,U)$.
\begin{thm}\label{tthm-1.10-atm}
Let $U:\mathcal{B}(\mathcal{H})\rightarrow Y$ be a non-null bounded liner operator and $\norm{U} < \lambda$.	Let $1\leq p<\infty$. Then for every $n$ we have 
	\begin{equation*}
		\frac{R_\lambda(\mathbb{D},p,U)}{n}\leq AP_\lambda(B_{Z},p, U)\leq \norm{Id: Z \rightarrow \ell^n_{1}}\, \frac{R_\lambda(\mathbb{D},p,U)}{n^{1/p}}.
	\end{equation*}
\end{thm}
By considering the function $\phi$ defined in \eqref{e-function} with $t=S(\Omega,B_{\ell^n_1})$ and proceeding along lines similar to those in the proof of Theorem \ref{tthm-1.10-atm}, one can verify that the result remains valid for any simply connected complete Reinhardt domain $\Omega$. We therefore state the result below and omit the details.
\begin{cor}\label{cor-1.10-atm}
	Let $U:\mathcal{B}(\mathcal{H})\rightarrow Y$ be a non-null bounded liner operator and $\norm{U} < \lambda$. Let $1\leq p<\infty$. Then for every $n$ we have 
	\begin{equation*}
		\frac{R_\lambda(\mathbb{D},p,U)}{n}\leq AP_\lambda(\Omega,p, U)\leq S(\Omega,B_{\ell^n_1})\, \frac{R_\lambda(\mathbb{D},p,U)}{n^{1/p}}.
	\end{equation*}
\end{cor}

In view of Theorem \ref{tthm-1.10-atm} and the fact $\norm{Id: B_{\ell^n_q}\rightarrow B_{\ell^n_1}} =n^{1-1/q}$ for $q \geq 1$, we deduce the following result for $Z=\ell^n_q$.
\begin{cor} \label{cor-1.5}
Let $U:\mathcal{B}(\mathcal{H})\rightarrow Y$ be a non-null bounded liner operator and $\norm{U} < \lambda$. Then for every $n$, $1\leq p<\infty$ and $1 \leq q \leq \infty$ we have 
\begin{equation*}
	\frac{R_\lambda(\mathbb{D},p,U)}{n}\leq AP_\lambda(B_{\ell^n_q},p, U)\leq  \frac{R_\lambda(\mathbb{D},p,U)}{n^{\frac{1}{p}+\frac{1}{q}-1}}.
\end{equation*}
\end{cor}
Unlike the Bohr radius constant $R_\lambda(\Omega,p, U)$, for which no exact value is currently known for any simply connected complete Reinhardt domain $\Omega \subset \mathbb{C}^n$ in dimension $n \geq 2$, the arithmetic Bohr radius can be explicitly computed in certain settings. In particular, for the unit ball of $\ell^n_1$ we are able to determine its precise value. This follows directly from Corollary \ref{cor-1.5} by taking $p=1$ and $q=1$.
\begin{cor} \label{cor-1.8}
Let $U:\mathcal{B}(\mathcal{H})\rightarrow Y$ be a non-null bounded liner operator and $\norm{U} < \lambda$.	Then for every $n\in \mathbb{N}$, we have 
	\begin{equation*}
		AP_\lambda(B_{\ell^n_1},1, U)=\frac{R_\lambda(\mathbb{D},1,U)}{n}.
	\end{equation*}
\end{cor}
It worth mentioning that Corollary \ref{cor-1.8} is actually pluriharmonic analogue of \cite[Theorem 3.1]{defant-2008}.
\vspace{1mm}

In view of the corollaries \ref{cor-1.10-atm}, \ref{cor-1.5}, and \ref{cor-1.8}, it is evident that to estimate $AP_\lambda(\Omega,p, U)$ it is enough to estimate $R_\lambda(\mathbb{D},p,U)$. To do this, we prove the following result in the case $U$ is the identity operator on $\mathbb{C}$ and $p=1$.
\begin{lem} \label{lem-1.1}
	
If $U=I:\mathbb{C} \rightarrow \mathbb{C}$ and $1 <\lambda$, then 
	\begin{equation*}
	R_\lambda(\mathbb{D},1,\mathbb{C}) \geq \frac{(\lambda-1)\pi}{4+(\lambda-1)\pi}.
	\end{equation*}
 \comment{If $U=I:\mathcal{B}(\mathcal{H}) \rightarrow \mathcal{B}(\mathcal{H})$ and $1 <\lambda$, then
 	\begin{equation*}
 		R_\lambda(\mathbb{D},1,\mathcal{B}(\mathcal{H})) \geq 	
 \end{equation*}}
\end{lem}
 Now, by making use of Lemma \ref{lem-1.1} and the corollaries \ref{cor-1.10-atm}, \ref{cor-1.5}, and \ref{cor-1.8}, we deduce the following important estimates, which we record as a corollary.
 \begin{cor}
 Let $\Omega \subset \mathbb{C}^n$ be a bounded simply connected complete Reinhardt domain. Let $U$ be the identity operator on $\mathbb{C}$ and $1 <\lambda$. Then for every $n \in \mathbb{N}$, we have
   
 	$$
 	\frac{(\lambda-1)\pi}{(4+(\lambda-1)\pi)\,n} \leq AP_\lambda(\Omega,1, \mathbb{C}).
 	$$
 \end{cor}
 This corollary actually provides a universal lower bound of the arithmetic Bohr radius for complex valued pluriharmonic functions on any bounded simply connected complete Reinhardt domain $\Omega \subset \mathbb{C}^n$.
 \vspace{1mm}
  
On the other hand, by applying the same arguments as in the pluriharmonic case in Theorem \ref{thm-1.8-atm}, Theorem \ref{tthm-1.10-atm}, and Corollary \ref{cor-1.10-atm}, we obtain the relation between the Bohr radius and the arithmetic Bohr radius for any Banach space $X$-valued holomorphic functions.
\begin{thm} \label{thm-1.8-atm-holo}
	Let $Z=(\mathbb{C}^n, ||.||)$ be a Banach space such that $\chi(\{e_k\}^n_{k=1})=1$. Then 
	\begin{equation*}
		A_\lambda(B_{Z},p,U) \geq \frac{\norm{Id: Z \rightarrow \ell^n_{1}}}{n} \, K_\lambda(B_{Z},p,U).
	\end{equation*}
\end{thm}

\comment{We now turn our attention to answering Question \ref{qsn-1.1}. To begin with, it is useful to point out some important differences between the study of Bohr’s theorem for vector-valued holomorphic functions and that for pluriharmonic functions. First, in the pluriharmonic setting we restrict ourselves to simply connected complete Reinhardt domains in $\mathbb{C}^n$. This restriction is necessary because the series expansion \eqref{e-1.3-a} is valid only on such domains. Second, in this framework we focus exclusively on $\mathcal{B}(\mathcal{H})$-valued pluriharmonic functions, and correspondingly restrict the operator $U$ to $\mathcal{B}(\mathcal{H})$. The reason is that the notion of complex conjugation is well defined in $\mathcal{B}(\mathcal{H})$, which is essential for our analysis. On the other hand, in the study of Bohr’s theorem for vector-valued holomorphic functions, the situation is somewhat simpler. Here, it suffices to consider complete Reinhardt domains (without the additional requirement of simply connectedness) because the domain of convergence of multivariable power series is complete Reinhardt domain and within its domain of convergence it defines a holomorphic function. Moreover, since complex conjugation does not play any role in this context, we are not constrained to $\mathcal{B}(\mathcal{H})$-valued functions. Instead, we can work with holomorphic functions taking values in any complex Banach space. Therefore, in addressing Question \ref{qsn-1.1}, we adopt the following framework: we consider bounded holomorphic functions with values in an arbitrary complex Banach space $X$, defined on complete Reinhardt domains in $\mathbb{C}^n$. In the following theorem, we first establish that the constant $K_{\lambda}(\Omega, p, U)$ is nonzero whenever $\Omega = B_{Z}$, which serves as the holomorphic analogue of Theorem \ref{thm-1.1}. By applying the same arguments as in the pluriharmonic case (see the discussion preceding Theorem \ref{thm-1.1}), it then follows that $K_{\lambda}(\Omega, p, U)$ is nonzero for every bounded complete Reinhardt domain in $\mathbb{C}^n$.
	\begin{thm} \label{thm-1.1-a}
		Let $X$ and $Y$ be any complex Banach spaces and $U: X \rightarrow Y$ be a non-null bounded linear operator such that $\norm{U}< \lambda$. Then, for $\lambda>1$ and $n \in \mathbb{N}$, we have
		%		\begin{equation} \label{e-homo-1}
			%			R^{m}_{\lambda}(B_{Z}, p,U) \geq \max \left\{\phi(\lambda,p,m)\,\frac{1}{\sup_{z \in B_{Z}}\norm{z_{p}}}, \, \phi(\lambda,p,m)\,\left(\frac{\lambda}{\norm{U}}\right)^{\frac{1}{m}}\, \frac{1}{\sup_{z \in B_{Z}}\norm{z_{p}}}\right\},
			%		\end{equation} 
		%		where $\phi(\lambda,p,m)=2^{-\frac{1}{pm}}\,  (4\lambda)^{-\frac{1}{m}}\, \norm{I-\real(a_0)}^{-\frac{1}{m}}$. Moreover,
		$$K_{\lambda}(B_{Z}, p,U) \geq  \frac{C}{\sup_{z \in B_{Z}}\norm{z}_{p}},$$
		where 
		$	C=\begin{cases}
			\max \left\{\left(\frac{\lambda^p - \norm{U}^p}{2\lambda^p - \norm{U}}\right)^{\frac{1}{p}}\, , \left(\frac{\lambda^p - \norm{U}^p}{\lambda^p - \norm{U}^p +1}\right)^{\frac{1}{p}}\, \frac{1}{\norm{U}}\right\}\,  & \text{for $\norm{U}\geq 1$},\\[3mm]
			\max \left\{\left(\frac{\lambda^p - \norm{U}^p}{2\lambda^p - \norm{U}}\right)^{\frac{1}{p}}, \left(\frac{\lambda^p - \norm{U}^p}{\lambda^p - \norm{U}^p +1}\right)^{\frac{1}{p}}\right\} & \text{for $0<\norm{U}< 1$}.
		\end{cases}
		$
\end{thm}}
\begin{thm}\label{tthm-1.10-atm-holo}
	Let $U:X\rightarrow Y$ be a non-null bounded liner operator and $\norm{U} < \lambda$.	Let $1\leq p<\infty$. Then for every $n$ we have 
	\begin{equation*}
		\frac{K_\lambda(\mathbb{D},p,U)}{n}\leq A_\lambda(B_{Z},p, U)\leq \norm{Id: Z \rightarrow \ell^n_{1}}\, \frac{K_\lambda(\mathbb{D},p,U)}{n^{1/p}}.
	\end{equation*}
\end{thm}
Theorems \ref{thm-1.8-atm-holo} and \ref{tthm-1.10-atm-holo} remain valid for any bounded complete Reinhardt domain, not only for $B_Z$. Moreover, we recover \cite[Lemma 4.3]{defant-2008} from Theorem \ref{thm-1.8-atm-holo} by choosing $U=I:\mathbb{C} \rightarrow \mathbb{C}$ and $p=1$.
\par In particular, Theorem \ref{tthm-1.10-atm-holo} yields the following result.
\begin{cor} \label{cor-1.8-a}
	Let $U:X\rightarrow Y$ be a non-null bounded liner operator and $\norm{U} < \lambda$.	Then for every $n \in \mathbb{N}$ we have 
	\begin{equation*}
		A_\lambda(B_{\ell^n_1},1, U)=\frac{K_\lambda(\mathbb{D},1,U)}{n}.
	\end{equation*}
\end{cor}
It is interesting to note that \cite[Theorem 3.1]{defant-2008} follows from Corollary \ref{cor-1.8-a} by choosing $U$ the identity operator on $\mathbb{C}$.
\begin{rem}
	For $X=\mathcal{B}(\mathcal{H})$, the bounds for $AP_{\lambda}(B_{Z}, p,X)$ established in Theorems \ref{thm-1.2-atm}–\ref{thm-1.6} extend to the constant $A_{\lambda}(B_{Z}, p,X)$ for any complex Banach space $X$, possibly with different constants. In particular, for $X=\mathcal{B}(\mathcal{H})$, this implies that the asymptotic behaviors of $AP_{\lambda}(B_{Z}, p,X)$ and $A_{\lambda}(B_{Z}, p,X)$ coincide . Since the proofs follow the same lines as those of the corresponding results for $AP_{\lambda}(B_{Z}, p,X)$, we omit the details. Moreover, an analogue of Theorem \ref{tthm-1.8-atm} holds for the constant $A_{\lambda}(\Omega, p,X)$ for any bounded complete Reinhardt domain.
\end{rem}
	\section{Application}
	This section is devoted to examples that demonstrate the scope of our results when applied to various classes of sequence spaces. Our first example is the family of mixed Minkowski spaces, which form a natural generalization of the classical Minkowski spaces. Given $m,n \in \mathbb{N}$ and $1 \leq s,t \leq \infty$, the mixed Minkowski space is defined as 
	$$\ell^m_s(\ell^n_t):=\{(z_{k})^m_{k=1}:z_{1}, \ldots,z_{m} \in \mathbb{C}^n\}$$ 
	with $\norm{(z_{k})^m_{k=1}}_{s,t}:=(\sum_{k=1}^{m}\norm{z_{k}}^s_t)^{1/s}$. In the special case $s=t$, this construction coincides with the classical Minkowski space $\ell^{mn}_s$. While earlier contributions such as \cite{Allu-Pal-arithnmetic,A-H-P-banach} dealt primarily with classical Minkowski spaces and the holomorphic case, the framework proposed in this article extends beyond that setting. It treats a significantly larger family of sequence spaces and enables a systematic study of holomorphic as well as pluriharmonic functions.
	\begin{example}[{\bf Mixed spaces}]
		Let $m,n \geq 2$. Then
		\begin{enumerate}
			\item Let $1 \leq s ,t \leq 2$. If $X=\mathcal{B}(\mathcal{H})$ is finite dimensional, then
			\begin{equation*}
				AP_{\lambda}(B_{\ell^m_s(\ell^n_t)}, p,X) \geq	\begin{cases}
					E_{1}(X)\,  \, \max \left\{\frac{1}{n^{2/t}\,m^{2/s}}, \, \frac{1}{e\,mn }\right\}\, \left(\frac{\lambda^p -1}{2\lambda^p - 1}\right)^{\frac{1}{p}} & \text{for $p=1$}, \\[2mm]
					E_{3}\, n^{\frac{1}{p}-\frac{2}{tp}-\frac{1}{t}}\,m^{\frac{1}{p}-\frac{2}{sp}-\frac{1}{s}} \left(\frac{\lambda^p -1}{2\lambda^p - 1}\right)^{\frac{1}{p}}\, & \text{for $p \geq 2$}, \\[2mm]
					
					E_{2}(X)  \max \left\{\frac{1}{ n^{\frac{2}{t}+\frac{1}{p}-1}m^{\frac{2}{s}+\frac{1}{p}-1}},  \frac{m^{2+\frac{4}{ps}-\frac{3}{p}-\frac{4}{s}}}{e^{(1-\theta)} n^{\frac{3}{p}+\frac{4}{t} -2-\frac{4}{pt}}}\right\}\left(\frac{\lambda^p -1}{2\lambda^p - 1}\right)^{\frac{1}{p}} \hspace{-1mm} & \text{for $1<p< 2$}, 
				\end{cases}
			\end{equation*}
			%Here, $E_{1}(X)$, $E_{2}(X)$, and $E_{3}$ are positive constants depending respectively on $X$, on $X$ and $p$, and on $p$.
			and 
			$$AP_{\lambda}(B_{\ell^m_s(\ell^n_t)}, p,X) \leq d\,  \lambda^{\frac{1}{\log\,(mn)}}\,\, \sqrt{\log \, (mn)} \, \left(\frac{1}{mn}\right)^{\frac{1}{p}}.
			$$
			If $X=\mathcal{B}(\mathcal{H})$ is infinite dimensional then
			$$\frac{\tilde{\Psi}_{1}(mn,\lambda,p,2)}{m^{1/s -1/2}\,n^{1/t-1/2}} \, \leq AP_{\lambda}(B_{Z}, p,X) \leq  \tilde{\Psi}_{2}(mn,\lambda,p,2).
			$$ 
			\item Let $2 \leq s,t \leq \infty$. If $X=\mathcal{B}(\mathcal{H})$ is finite dimensional then
			\begin{equation*}
				AP_{\lambda}(B_{\ell^m_s(\ell^n_t)}, p,X) \geq	\begin{cases}
					E_{1}(X)\,  \, \max \left\{\frac{1}{n^{\frac{1}{2}+\frac{1}{t}}\, m^{\frac{1}{2}+\frac{1}{s}}}, \, \frac{1}{e\, mn}\right\}\, \left(\frac{\lambda^p -1}{2\lambda^p - 1}\right)^{\frac{1}{p}} & \text{for $p=1$}, \\[2mm]
					E_{3}\, n^{-\frac{1}{t}}\, m^{-\frac{1}{s}}\left(\frac{\lambda^p -1}{2\lambda^p - 1}\right)^{\frac{1}{p}}\, & \text{for $p \geq 2$}, \\[2mm]
					
					E_{2}(X)  \max \left\{\frac{1}{n^{\frac{1-\theta}{2}+\frac{1}{t}}\,m^{\frac{1-\theta}{2}+\frac{1}{s}}}, \frac{m^{\theta -\frac{\theta}{s}-1}}{e^{(1-\theta)} n^{1-\theta +\frac{\theta}{t}}}\right\}\left(\frac{\lambda^p -1}{2\lambda^p - 1}\right)^{\frac{1}{p}}  \hspace{-2mm} & \text{for $1<p< 2$}, 
				\end{cases}
			\end{equation*}
			and
			$$AP_{\lambda}(B_{\ell^m_s(\ell^n_t)}, p,X) \leq d\,n^{1-\frac{1}{t}-\frac{1}{p}}m^{1-\frac{1}{s}-\frac{1}{p}}\, \lambda^{\frac{2}{\log \, n}}\, \sqrt{\log mn}.
			$$
			If $X=\mathcal{B}(\mathcal{H})$ is infinite dimensional then
			$$\tilde{\Psi}_{1}(mn,\lambda,p,2)\, \leq AP_{\lambda}(B_{Z}, p,X) \leq m^{\frac{1}{2}-\frac{1}{s}}\,n^{\frac{1}{2}-\frac{1}{t}} \, \tilde{\Psi}_{2}(mn,\lambda,p,2).
			$$
		\end{enumerate}
		Here, the constants $E_{1}(X)$, $E_{2}(X)$, $E_{3}$, and $d$ are as in Theorem \ref{thm-1.2-atm}.	
	\end{example}
	\begin{pf}
	For $1\leq s,t,l,r \leq \infty$, it is easy to check that the identity operator between mixed sequence spaces factorizes as 
		$$\norm{Id:\ell^m_s(\ell^n_t) \rightarrow \ell^m_l(\ell^n_r)}= \norm{Id:\ell^m_s \rightarrow \ell^m_l}\, \norm{Id:\ell^n_t \rightarrow \ell^n_r}$$
		 (see \cite[p. 188]{defant-2003}). Moreover, for finite-dimensional $\ell^n_p$ spaces, the norm of the identity operator satisfies: $\norm{Id:\ell^n_s \rightarrow \ell^n_t}=1$ for $s \leq t$ and $\norm{Id:\ell^n_s \rightarrow \ell^n_t}=n^{1/t - 1/s}$ for $s>t$. Applying these observations together with Theorems \ref{thm-1.2-atm}, \ref{thm-1.4-atm}, and \ref{thm-1.6}, we obtain the desired result. This completes the proof.
	\end{pf}

	The scope of our results further includes classical symmetric Banach sequence spaces such as Lorentz spaces $\ell_{s,t}$ as well as Orlicz spaces $\ell_{\psi}$; see \cite{Lindenstrauss-book} for their definitions. 
	\begin{example}[{\bf Lorentz spaces}]
		Let $1 \leq s,t \leq \infty$.
		\begin{enumerate}
			\item Let $1 \leq s <2$ and $1\leq t \leq \infty$, or $s=2$ and $1 \leq t \leq 2$. If $X=\mathcal{B}(\mathcal{H})$ is finite dimensional 
			\begin{equation*}
				 AP_{\lambda}(\ell^n_{s,t}, p,X) \geq	\begin{cases}
					E_{1}(X)\,  \max \left\{\frac{1}{n^{\frac{1}{2}+\frac{1}{s}}\, \norm{Id: \ell^n_{2} \rightarrow \ell^n_{s,t}}}, \, \frac{1}{e\,n}\right\}\, \left(\frac{\lambda^p -1}{2\lambda^p - 1}\right)^{\frac{1}{p}} & \text{for $p=1$}, \\[2mm]
					E_{3}\, \frac{1}{n^{\frac{1}{s}}}\norm{Id:\ell^n _2 \rightarrow \ell^n_{s,t}}^{-\frac{2}{p}}\, \left(\frac{\lambda^p -1}{2\lambda^p - 1}\right)^{\frac{1}{p}}\, & \text{for $p \geq 2$}, \\[2mm]
					
					E_{2}(X)\,  \max \left\{\frac{1}{n^{1-\frac{\theta}{2}+\frac{1}{s}} \norm{Id: \ell^n_{2} \rightarrow \ell^n_{s,t}}},  \frac{\norm{Id:\ell^n _2 \rightarrow \ell^n_{s,t}}^{-\theta}}{e^{(1-\theta)}\,n^{1-\theta+\frac{\theta}{s}}}\right\}\left(\frac{\lambda^p -1}{2\lambda^p - 1}\right)^{\frac{1}{p}} \hspace{-2mm} & \text{for $1<p< 2$}, 
				\end{cases}
			\end{equation*}
			and 
			$$AP_{\lambda}(B_{\ell^n_{s,t}}, p,X) \leq d\,  \lambda^{\frac{1}{\log\,n}}\,\, \sqrt{\log \, n} \, \left(\frac{1}{n}\right)^{\frac{1}{p}}.
			$$
			If $X=\mathcal{B}(\mathcal{H})$ is infinite dimensional, then
			$$\frac{\tilde{\Psi}_{1}(n,\lambda,p,2)}{\norm{I:\ell^n _2 \rightarrow \ell^n_{s,t}}} \, \leq AP_{\lambda}(B_{\ell^n_{s,t}}, p,X)  \leq  \tilde{\Psi}_{2}(n,\lambda,p,2).
			$$
			\item Let $2<s \leq \infty$ and $1 \leq t \leq \infty$. If $X=\mathcal{B}(\mathcal{H})$ is finite dimensional
			\begin{equation*} 
			AP_{\lambda}(\ell^n_{s,t}, p,X) \geq	\begin{cases}
				E_{1}(X)\,  \max \left\{\frac{1}{n^{\frac{1}{2}+\frac{1}{s}}}, \, \frac{1}{e\,n}\right\}\, \left(\frac{\lambda^p -1}{2\lambda^p - 1}\right)^{\frac{1}{p}} & \text{for $p=1$}, \\[2mm]
				E_{3}\, \frac{1}{n^{\frac{1}{s}}}\, \left(\frac{\lambda^p -1}{2\lambda^p - 1}\right)^{\frac{1}{p}}\, & \text{for $p \geq 2$}, \\[2mm]
				
				E_{2}(X)\,  \max \left\{\frac{1}{n^{1-\frac{\theta}{2}+\frac{1}{s}}},  \frac{1}{e^{(1-\theta)}\,n^{1-\theta+\frac{\theta}{s}}}\right\}\left(\frac{\lambda^p -1}{2\lambda^p - 1}\right)^{\frac{1}{p}} \hspace{-2mm} & \text{for $1<p< 2$}, 
			\end{cases}
			\end{equation*}
			and $$AP_{\lambda}(B_{\ell^n_{s,t}}, p,X) \leq d\,  \lambda^{\frac{1}{\log\,n}}\,\, \sqrt{\log \, n} \, \left(\frac{1}{n}\right)^{\frac{1}{p}+\frac{1}{s}-\frac{1}{2}}.
			$$ If $X=\mathcal{B}(\mathcal{H})$ is infinite dimensional, then $$\tilde{\Psi}_{1}(n,\lambda,p,2) \, \leq AP_{\lambda}(B_{\ell^n_{s,t}}, p,X)  \leq \left(\frac{1}{n}\right)^{\frac{1}{s}-\frac{1}{2}} \, \tilde{\Psi}_{2}(n,\lambda,p,2).
			$$
			%; if $2 < s \leq \infty$, $1 \leq t \leq 2$, the lower bound of $R_{\lambda}(B_{\ell^n_{s,t}}, p,X)$ is same as \eqref{e-1.32-a} and $R_{\lambda}(B_{\ell^n_{s,t}}, p,X) \leq d\,  \lambda^{2/\log \, n}\, n^{-\left(\frac{1}{p}+\frac{1}{s}\right)}\, \sqrt{\log \, n}$.
		\end{enumerate}
	\end{example}
	\begin{pf}
		We begin by observing that $\norm{\sum_{k=1}^{n}e^{*}_{k}}_{\ell_{s,t}^{*}}=n^{1-1/s}$. Moreover, Lorentz sequence spaces satisfy the following lexicographical inclusion property: $$\ell_{p,q} \subseteq \ell_{s,t}\,\, \mbox{if and only if}\,\,\, p<s, \,\, \mbox{or}\,\, p=s\,\, \mbox{and}\,\, q \leq t.
		$$
		 These facts, together with Theorems \ref{thm-1.3-a}(1), \ref{thm-1.4-atm}, and \ref{thm-1.6}, immediately yield the assertion in part $(1)$. 
		 \vspace{1mm}
		 
		 We now turn to part $(2)$. Suppose first that $2<s \leq \infty$ and $2 \leq t \leq \infty$. In this case, $\ell_{s,t}$ is $2$-convex (see \cite[p.189]{defant-2003}), and the conclusion follows directly from Theorems \ref{thm-1.3-a}$(2)$, \ref{thm-1.4-atm}, and \ref{thm-1.6}. It remains to consider the case $2 < s \leq \infty$, $1 \leq t \leq 2$. Here we apply Theorems \ref{thm-1.2-atm}, \ref{thm-1.4-atm}, and \ref{thm-1.6}, together with the estimates 
		 $$
		 \norm{Id: \ell^n_{2} \rightarrow \ell^n_{s,t}} \leq 1\,\,\,\mbox{and}\,\,\, \norm{Id: \ell^n_{s,t} \rightarrow \ell^n_{2}} \leq n^{1/2 - 1/s}
		 $$ 
		 (see \cite[p.189]{defant-2003}). This completes the proof.
	\end{pf}
	\begin{example} [{\bf Orlicz spaces}]
		Let $\psi$ be an Orlicz function which satisfies the $\Delta_{2}$-condition. 
		\begin{enumerate}
			\item Let $a^2 \leq T \psi(a)$ for all $a$ and some $T>0$.  If $X=\mathcal{B}(\mathcal{H})$ is finite dimensional, then
			\begin{equation*}
				AP_{\lambda}(B_{\ell^n_{\psi}}, p,X) \geq	\begin{cases}
					E_{1}(X)\,  \, \max \left\{\frac{\psi^{-1}(1/n)}{\sqrt{n}\, \norm{Id: \ell^n_{2} \rightarrow \ell^n_{\psi}}}, \, \frac{1}{e\,n}\right\}\, \left(\frac{\lambda^p -1}{2\lambda^p - 1}\right)^{\frac{1}{p}} & \text{for $p=1$}, \\[2mm]
					E_{3}\, \psi^{-1}(1/n)\,\norm{Id:\ell^n _2 \rightarrow \ell^n_{\psi}}^{-\frac{2}{p}}\, \left(\frac{\lambda^p -1}{2\lambda^p - 1}\right)^{\frac{1}{p}}\, & \text{for $p \geq 2$}, \\[2mm]
					
					E_{2}(X)  \max \left\{\frac{\psi^{-1}(1/n)}{(\sqrt{n})^{1-\theta} \norm{Id: \ell^n_{2} \rightarrow \ell^n_{\psi}}},  \frac{\norm{Id:\ell^n _2 \rightarrow \ell^n_{\psi}}^{-\theta}}{(en)^{(1-\theta)} (\psi^{-1}(1/n))^{-\theta}}\right\}\left(\frac{\lambda^p -1}{2\lambda^p - 1}\right)^{\frac{1}{p}} \hspace{-4mm}  & \text{for $1<p< 2$}, 
				\end{cases}
			\end{equation*}
			and $$AP_{\lambda}(B_{\ell^n_{\psi}}, p,X) \leq d\,  \lambda^{\frac{1}{\log\,n}}\,\, \sqrt{\log \, n} \, \left(\frac{1}{n}\right)^{\frac{1}{p}}.
			$$
			If $X=\mathcal{B}(\mathcal{H})$ is infinite dimensional, 
			$$\frac{\tilde{\Psi}_{1}(n,\lambda,p,2)}{\norm{Id:\ell^n _2 \rightarrow \ell^n_{\psi}}} \, \leq AP_{\lambda}(B_{\ell^n_{\psi}}, p,X) \leq  \tilde{\Psi}_{2}(n,\lambda,p,2).
			$$
			
			\item Let $\psi(\beta \,a) \leq T \beta^2\,\psi(a)$ for $0 \leq \beta,a \leq 1$ and some $T>0$. If $X=\mathcal{B}(\mathcal{H})$ is finite dimensional,
			\begin{equation*}
				AP_{\lambda}(B_{\ell^n_{\psi}}, p,X) \geq	\begin{cases}
					E_{1}(X)\,  \, \max \left\{\frac{\psi^{-1}(1/n)}{\sqrt{n}}, \, \frac{1}{e\,n}\right\}\, \left(\frac{\lambda^p -1}{2\lambda^p - 1}\right)^{\frac{1}{p}} & \text{for $p=1$}, \\[2mm]
					E_{3}\,\psi^{-1}(1/n)\, \left(\frac{\lambda^p -1}{2\lambda^p - 1}\right)^{\frac{1}{p}}\, & \text{for $p \geq 2$}, \\[2mm]
					
					E_{2}(X)\,  \max \left\{\frac{\psi^{-1}(1/n)}{(\sqrt{n})^{1-\theta}}, \, \frac{1}{(en)^{(1-\theta)}\, (\psi^{-1}(1/n))^{-\theta}}\right\}\left(\frac{\lambda^p -1}{2\lambda^p - 1}\right)^{\frac{1}{p}}  & \text{for $1<p< 2$}, 
				\end{cases}
			\end{equation*}
			and 
			$$AP_{\lambda}(B_{\ell^n_{\psi}}, p,X) \leq d\,  \lambda^{\frac{1}{\log\,n}}\, \,\psi^{-1}\left(\frac{1}{n}\right)\,\, \sqrt{\log \, n}\, \left(\frac{1}{n}\right)^{\frac{1}{p}-\frac{1}{2}}.$$
			If $X=\mathcal{B}(\mathcal{H})$ is infinite dimensional, 
			$$\tilde{\Psi}_{1}(n,\lambda,p,2) \, \leq AP_{\lambda}(B_{\ell^n_{\psi}}, p,X) \leq \sqrt{n}\, \psi^{-1}\left(\frac{1}{n}\right) \, \tilde{\Psi}_{2}(n,\lambda,p,2).
			$$
		\end{enumerate}
	\end{example}
	\begin{pf}
		It is important to note the facts $\norm{\sum_{k=1}^{n}e^{*}_{k}}_{\ell_{\psi}^{*}}\, \norm{\sum_{k=1}^{n}e_{k}}_{\ell_{\psi}}=n$ 
		and
		$$
		\norm{\sum_{k=1}^{n}e_{k}}_{\ell_{\psi}}=\frac{1}{\psi^{-1}\left(\frac{1}{n}\right)}
		$$
		(see \cite[p.192]{defant-2003}). The condition imposed on $\psi$ in part $(1)$ guarantees that $\ell_{\psi} \subseteq \ell_2$.  On the other hand, the condition in part $(2)$ ensures that $\ell_{\psi}$ is $2$-convex (see \cite[p.189]{defant-2003}). Taking these properties into account, the desired conclusion follows directly from Theorems \ref{thm-1.3-a} and \ref{thm-1.6}, together with the preceding facts. This completes the proof.
	\end{pf}	
	\section{Preliminaries}
	We adopt standard notation and terminology from Banach space theory. Throughout, all Banach spaces $W$ are assumed to be complex. Their topological duals are denoted by $W^{*}$, and their open unit balls by $B_{W}$. A Banach space $W$ is said to have cotype $t\in [0,\infty]$ if there exists a constant $C>0$ such that, for every finite family vectors $w_{1},\ldots,w_{n}\in W$, 
	$$\left(\sum_{k=1}^{n}\norm{w_{k}}^t\right)^{\frac{1}{t}} \leq C \left(\int_{0}^{1}\norm{\sum_{k=1}^{n}r_{k}(s)w_{k}}^2\, ds\right)^{\frac{1}{2}},$$
	where $r_{k}$ denotes the $k$-th Rademacher function on $[0,1]$. We define 
	$$Cot(W):=\inf\{2 \leq t \leq \infty: W\, \mbox{has cotype}\, t\}.
	$$
	It is well known that every Banach space has cotype $\infty$.  When $Cot(W)=\infty$,  we adopt the convention $1/Cot(W)=0$. A Schauder basis $\{w_{k}\}$ of a Banach space $W$ is called unconditional if there is a constant $c\geq 1$ such that $$\norm{\sum_{j=1}^{k}\varepsilon_{j}\mu_{j}w_{j}} \leq c \norm{\sum_{j=1}^{k}\mu_{j}w_{j}}
	$$
	for all $k\in\mathbb{N}$, all scalars $\mu_{j}\in\mathbb{C}$, and all $\varepsilon_{j}\in\mathbb{C}$ with $|\varepsilon_{j}|\le 1$. The smallest such constant is denoted by $\chi({w_{k}})$ and is called the unconditional basis constant of ${w_{k}}$. We say that ${w_{k}}$ is a $1$-unconditional basis if $\chi({w_{k}})=1$. The unconditional basis constant of the Banach space $W$ is defined by $\chi(W):=\inf \chi(\{w_{k}\})\in [1,\infty]$, the infimum taken over all unconditional basis $\{w_{k}\}$ of $W$. A Banach space $Z$ satisfying $\ell_1 \subset Z \subset c_{0}$ (with normal inclusions) is called a Banach sequence space if the canonical unit vectors ${e_{k}}$ form a $1$-unconditional basis of $Z$. A Banach lattice $Z$ is said to be $2$-convex if there exists a constant $C>0$ such that 
	$$\norm{\left(\sum_{k=1}^{n}|x_{k}|^2\right)^{\frac{1}{2}}} \leq C 	\left(\sum_{k=1}^{n}\norm{x_{k}}^2\right)^{\frac{1}{2}}
	$$
	for all finite families $x_{1},\ldots,x_{n}\in Z$. Let $Z=(\mathbb{C}^n, ||.||)$ be a Banach space, $Y$ any Banach space, and $m \in \mathbb{N}$. We denote by $\mathcal{P}(^m Z,Y)$ the space of all $m$-homogeneous polynomials $Q:Z \rightarrow Y$ of the form $Q(z)=\sum_{|\alpha|=m}c_{\alpha}z^{\alpha}$, together with the norm $\norm{Q}_{\mathcal{P}(^m Z,Y)}:=\sup_{z \in B_{Z}}||Q(z)||_{Y}$. The unconditional basis constant of the family of monomials $z^{\alpha}$, $\alpha \in (\mathbb{N} \cup \{0\})^n$ is denoted by $\chi_{M}(\mathcal{P}(^m Z,Y))$.
 We denote by $\mathcal{PH}(^m Z,\mathcal{B}(H))$ the space of all $m$-homogeneous pluriharmonic polynomials $P:Z \rightarrow \mathcal{B}(H)$ of the form 
	\begin{equation} \label{e-1.3-aa}
		P(z)= \sum_{|\alpha|=m} a_{\alpha}\, z^{\alpha} +  \sum_{|\alpha|=m} b^{*}_{\alpha}\, \bar{z}^{\alpha}
	\end{equation} 
	and set $\norm{P}_{\mathcal{PH}(^m Z,\mathcal{B}(H))}:=\sup_{z \in B_{Z}}|||P(z)||_{\mathcal{B}(H)}$. 
	\par A bounded linear operator $U:W\rightarrow Y$ between two Banach spaces $W$ and $Y$ is said to be $(r,t)$-summing, $r,t \in [1,\infty)$, if exists a constant $C>0$ such that for every finite collection $w_{1}, \ldots,w_{n} \in W$,
	$$\left(\sum_{k=1}^{n} \norm{U(w_{k})}^r\right)^{1/r} \leq C \, \sup_{\phi \in B_{X^{*}}} \left(\sum_{k=1}^{n} |\phi(w_{k})|^t\right)^{1/t}.
	$$
	The smallest constant $C$ for which this inequality holds is denoted by $\pi_{r,t}(U)$. In the special case $r=t$, the operator $U$ is called $r$- summing, and we write $\pi_{r}(U)$ instead of $\pi_{r,r}(U)$.
	\section{Auxiliary Lemmas} 
	%In the scalar case, a standard approach to deriving nontrivial estimates for Bohr radii of holomorphic functions is to first analyze $m$-homogeneous polynomials, as this reduction often isolates the essential difficulties while providing sharper insight into the general framework. 
	The following definition introduces the $m$-homogeneous analogue of the definition of Bohr radius for pluriharmonic functions, which plays a fundamental role in the analysis carried out in this article. Let $\Omega\subset \mathbb{C}^n$ be a simply connected complete Reinhardt domain and $n\in \mathbb{N}$. Let $U:X\rightarrow Y$ be a bounded liner operator and $\norm{U} \leq \lambda$. For $1 \leq p < \infty$, the $\lambda$-powered Bohr radius of $U$ with respect to $\mathcal{PH}(^m \Omega,X)$, denoted by $R^m_{\lambda}(\Omega, p,U)$, is defined to be the supremum of all $r\geq 0$ such that for every pluriharmonic polynomials $P:\Omega \rightarrow X$ of the form \eqref{e-1.3-aa} we have 
	\begin{equation} \label{e-1.4-aa}
		\sup_{z \in r\Omega}\, \sum_{|\alpha|=m} (\norm{U(a_{\alpha})}^p_{Y} + \norm{U(b_{\alpha})}^p_{Y})|z^\alpha|^p \leq \lambda^p\,\norm{f}^p_{\Omega,X}. 
	\end{equation}
	
	\noindent Here $\norm{f}_{\Omega,X}:=\sup_{z \in \Omega}\,\norm{f(z)}_{X}$. Set $R^m(\Omega, p,U):=R^m_{1}(\Omega, p,U)$, $R^m_{\lambda}(\Omega, p,X):=R^m_{\lambda}(\Omega, p,U)$ whenever $U=I:X\rightarrow X$, $R^m(\Omega, p,X):=R^m_{1}(\Omega, p,X)$, $R^m_{\lambda}(\Omega, p):=R^m_{\lambda}(\Omega, p,\mathbb{C})$, and $R^m(\Omega,p):=R^m_1(\Omega,p)$. It is straightforward that $R^m_{\lambda}(\Omega, p,U)=\lambda^{1/m}\,R^m(\Omega, p,U)$. Similarly as above, we can define the constant $K^m_{\lambda}(\Omega, p,U)$ for the class $\mathcal{P}^n_m(W)$ {\it i.e.}, for any Banach space $W$ valued $m$-homogeneous holomorphic polynomial $Q(z)=\sum_{|\alpha|=m}c_{\alpha}\, z^{\alpha}$. Clearly, we have $K^m_{\lambda}(\Omega, p,U)=\lambda^{1/m}\,K^m(\Omega, p,U)$. Using Bohnenblust-Hille inequality, Defant {\it et al.} \cite{defant-2011} have shown that the constant $K^m(\mathbb{D}^n,1)$ is hypercontractive.
	\vspace{2mm}
	
%	Schwarz–Pick type lemmas and distortion estimates for holomorphic and harmonic functions constitute an important direction of research in geometric function theory. Coefficient-type Schwarz–Pick estimates for complex-valued harmonic functions on the unit disc were obtained in \cite{Abu-2010, chen-2011}, and subsequently extended to complex-valued pluriharmonic functions on $B_{\ell^n_q}$ in \cite[Theorem 2.7]{hamada-JFA-2022}. 
	%Among the more recent developments in this area are the boundary Schwarz–Pick lemma and associated rigidity results developed by Bracci {\it et al.} \cite{Roth-schwarz-Adv-Math} and the geometric Schwarz-type estimates for spherically convex functions established by Kourou and Roth \cite{Roth-Scwarw-Canad}. 
%	In analogy with \cite[Theorem 2.7]{hamada-JFA-2022}, we now extend these estimates to the operator valued setting by establishing a coefficient-type Schwarz–Pick lemma for pluriharmonic functions defined on general complete Reinhardt domains. This lemma constitutes a fundamental ingredient in the proofs of several main results of this article.
The following lemma, proved in \cite{Himadri-local-Banach-1}, provides a coefficient-type Schwarz–Pick inequality for operator valued pluriharmonic functions defined on general complete Reinhardt domains. It plays an essential role in the proofs of several of our main results.
	\begin{lem} \label{lem-3.1} \cite{Himadri-local-Banach-1}
		Let $f \in \mathcal{PH}(Z,X)$ be of the form \eqref{e-1.3-a}. Then for all $|\alpha|=m\geq 1$, we have $\norm{\sum_{|\alpha|=m}(a_{\alpha} \pm b_{\alpha})\, z^{\alpha}}_{B_{Z},X} \leq 4 \,\norm{\norm{f}_{B_{Z},X}\,I-\real(a_{0})}$. Moreover, if $B_{Z}=B_{\ell^n_{q}}$ then $\norm{a_{\alpha}+b_{\alpha}}	\leq \frac{4}{\pi} \, \rho_{\alpha}\, \norm{\sum_{|\alpha|=m}(a_{\alpha}+b_{\alpha})\, z^{\alpha}}_{B_{Z},X}$,  $\norm{a_{\alpha}-b_{\alpha}}	\leq \frac{4}{\pi}\,\rho_{\alpha}\, \norm{\sum_{|\alpha|=m}(a_{\alpha}-b_{\alpha})\, z^{\alpha}}_{B_{Z},X}$, where $\rho_{\alpha}:=\left(\frac{|\alpha|^{|\alpha|}}{\alpha^{\alpha}}\right)^{1/q}$.
		\comment{\begin{enumerate}
				\item $\norm{\sum_{|\alpha|=m}(a_{\alpha}+b_{\alpha})\, z^{\alpha}}_{B_{Z},X} \leq 4 \,\gamma_{0} \norm{f}_{B_{Z},X}$,
				$\norm{\sum_{|\alpha|=m}(a_{\alpha}-b_{\alpha})\, z^{\alpha}}_{B_{Z},X} \leq 4 \, \gamma_{0}\, \norm{f}_{B_{Z},X}$;
				\item $\norm{a_{\alpha}+b_{\alpha}}	\leq \frac{4}{\pi} \, \rho_{\alpha}\, \norm{\sum_{|\alpha|=m}(a_{\alpha}+b_{\alpha})\, z^{\alpha}}_{B_{Z},X}$,  $\norm{a_{\alpha}-b_{\alpha}}	\leq \frac{4}{\pi}\,\rho_{\alpha}\, \norm{\sum_{|\alpha|=m}(a_{\alpha}-b_{\alpha})\, z^{\alpha}}_{B_{Z},X}$,
			\end{enumerate}
			where $\gamma_{0}:=\norm{I-\real(a_{0})}$ and $\rho_{\alpha}:=\left(\frac{|\alpha|^{|\alpha|}}{\alpha^{\alpha}}\right)^{\frac{1}{q}}$.}
	\end{lem}
	
	
	\comment{\begin{lem} \label{lem-3.2}
			Let $P(z)=\sum_{|\alpha|=m} a_{\alpha}\, z^{\alpha} + \sum_{|\beta|=m} b^{*}_{\alpha}\, \bar{z}^{\alpha} \in \mathcal{PH}^m_{n}$. Then for all $|\alpha|=m\geq 1$, we have
			\begin{equation} \label{e-homo-coeff-3.6}
				\norm{a_{\alpha}+b_{\alpha}}	\leq \frac{4}{\pi} \, \left(\frac{|\alpha|^{|\alpha|}}{\alpha^{\alpha}}\right)^{\frac{1}{q}}\, \norm{\sum_{|\alpha|=m}(a_{\alpha}+b_{\alpha})\, z^{\alpha}}_{\infty},  \norm{a_{\alpha}-b_{\alpha}}	\leq \frac{4}{\pi} \, \left(\frac{|\alpha|^{|\alpha|}}{\alpha^{\alpha}}\right)^{\frac{1}{q}}\, \norm{\sum_{|\alpha|=m}(a_{\alpha}-b_{\alpha})\, z^{\alpha}}_{\infty}
			\end{equation}
		\end{lem}
		\begin{pf}
			To prove the left inequality of \eqref{e-homo-coeff-3.6}, we consider the operator-valued holomorphic function $\tilde{h}(z)=\sum_{m=0}^{\infty}\sum_{|\alpha|=m} (a_{\alpha} + b_{\alpha})z^{\alpha}$ on $B_{\ell^n _q}$. For any $\phi \in X^{*}$ with $\norm{\phi} \leq 1$, $\tilde{g}=\phi \circ \tilde{h}$ is a complex-valued holomorphic function on $B_{\ell^n _q}$ with $\norm{\tilde{g}}_{\infty} \leq \norm{\tilde{h}}_{\infty}$. Moreover, $\tilde{g}(z)=\sum_{m=0}^{\infty}\sum_{|\alpha|=m} \phi (a_{\alpha} +b_{\alpha})z^{\alpha}$ for $z \in B_{\ell^n _q}$. Observe that $\tilde{g}$ is also a complex-valued pluriharmonic function, and hence from \cite[Theorem 2.6]{hamada-JFA-2022}, we have for $|\alpha|=m\geq 1$,
			\begin{equation} \label{e-3.5-a}
				\norm{a_{\alpha} +b_{\alpha}}=\sup _{\phi \in B_{X^{*}}} |\phi(a_{\alpha}+b_{\alpha})| \leq \frac{4}{\pi} \left(\frac{|\alpha|^{|\alpha|}}{\alpha^{\alpha}}\right)^{\frac{1}{q}} \norm{\tilde{g}}_{\infty} \leq \frac{4}{\pi} \left(\frac{|\alpha|^{|\alpha|}}{\alpha^{\alpha}}\right)^{\frac{1}{q}} \norm{\tilde{h}}_{\infty}.
			\end{equation}
			The right inequality in \eqref{e-homo-coeff-3.6} follows by applying similar arguments as above to the function $\tilde{P}=-iP$.
	\end{pf}}
	
	\comment{\begin{lem}
			Let $f(z)=\sum_{m=0}^{\infty}\sum_{|\alpha|=m} a_{\alpha}\, z^{\alpha} + \sum_{m=1}^{\infty}\sum_{|\beta|=m} b^{*}_{\alpha}\, \bar{z}^{\alpha} \in \mathcal{PH}_{n}$. Then for all $m\geq 1$, we have
			\begin{equation} \label{e-homo-coeff-33}
				\norm{\sum_{|\alpha|=m}(a_{\alpha}+b_{\alpha})\, z^{\alpha}} \leq 4 \norm{I-\real(a_{0})}\, \norm{P}_{\infty},
			\end{equation}
			and
			\begin{equation} \label{e-homo-coeff-4}
				\norm{\sum_{|\alpha|=m}(a_{\alpha}-b_{\alpha})\, z^{\alpha}} \leq 4 \norm{I-\real(a_{0})}\, \norm{P}_{\infty}.
			\end{equation}
		\end{lem}
	}
	The following lemma, established in \cite{Himadri-local-Banach-1}, provides a fundamental connection between the Bohr radius for pluriharmonic functions and that for homogeneous pluriharmonic functions. Conceptually, it may be viewed as a pluriharmonic counterpart of \cite[Lemma 3.2]{defant-2012}, where analogous ideas were developed for Minkowski spaces in the setting of the powered Bohr inequality. In contrast to that earlier result, the following lemma is formulated in the general framework of Banach spaces.
	\begin{lem} \cite{Himadri-local-Banach-1} \label{lem-3.3}
		Let $X=\mathcal{B}(\mathcal{H})$ and $Y$ any complex Banach space, and $U: X \rightarrow Y$ be a non-null bounded linear operator such that $\norm{U}< \lambda$. Then, for all $p\in[1,\infty)$, $\lambda>1$, and $n \in \mathbb{N}$, we have
		\begin{enumerate}
			\item $\left(\frac{\lambda^p - \norm{U}^p}{2\lambda^p - \norm{U}^p}\right)^{\frac{1}{p}}\, \inf_{m\in \mathbb{N}}\, R^{m}_{\lambda}(B_{Z}, p,U) \leq R_{\lambda}(B_{Z}, p,U) \leq \inf_{m\in \mathbb{N}}\, R^{m}_{\lambda}(B_{Z}, p,U)$
			
			\item $\left(\frac{\lambda^p - \norm{U}^p}{\lambda^p - \norm{U}^p +1}\right)^{\frac{1}{p}}\, \inf_{m\in \mathbb{N}}\, R^{m}(B_{Z}, p,U) \leq R_{\lambda}(B_{Z}, p,U) \leq \lambda \, \inf_{m\in \mathbb{N}}\, R^{m}(B_{Z}, p,U)$.
		\end{enumerate}
	\end{lem}
	%\begin{rem} \label{rem-4.1}
	It worth mentioning that analogous results of Lemma \ref{lem-3.3} hold for the constants $K^m_{\lambda}(B_{Z}, p,U)$ and $K_{\lambda}(B_{Z}, p,U)$.
	\vspace{1mm}
	
%	\end{rem}
	%\noindent  Its proof is similar to that of \cite[Lemma 2.5]{defant-2004} and is therefore omitted. 
	%This lemma plays a pivotal role in the development of several of our main results.
%	\begin{lem} \label{lem-3.5}
%		Let $\Omega_{1}$ and $\Omega_{2}$ be two bounded complete Reinhardt domains in $\mathbb{C}^n$. Then we have $$\frac{R_{\lambda}(\Omega_{2}, p,U)}{S(\Omega_{1},\Omega_{2})S(\Omega_{2},\Omega_{1})} \, \leq R_{\lambda}(\Omega_{1}, p,U) \leq S(\Omega_{1},\Omega_{2})S(\Omega_{2},\Omega_{1}) \, R_{\lambda}(\Omega_{2}, p,U).
%		$$
%	\end{lem}
%	\noindent As an immediate application, we observe that the Bohr radius $R_{\lambda}(\mathbb{D}^n, p,U)$ serves as a universal lower bound for the Bohr radius $R_{\lambda}(\Omega, p,U)$ of every complete Reinhardt domain.
%	\begin{lem} \label{lem-3.4}
%		For any complete Reinhardt domain $\Omega$, we have 
%		$
%		R_{\lambda}(\Omega, p,U) \geq R_{\lambda}(\mathbb{D}^n, p,U).
%		$
%	\end{lem}
The following lemma compares the arithmetic Bohr radius constants for two complete Reinhardt domains. This result serves as a key ingredient in the proofs of several of our main theorems.
\begin{lem} \label{lem-3.6}
	Let $\Omega_{1}$ and $\Omega_{2}$ be two bounded simply connected complete Reinhardt domains in $\mathbb{C}^n$. Then we have $$\frac{AP_{\lambda}(\Omega_{2}, p,U)}{S(\Omega_{2},\Omega_{1})} \, \leq AP_{\lambda}(\Omega_{1}, p,U) \leq S(\Omega_{1},\Omega_{2}) \, AP_{\lambda}(\Omega_{2}, p,U).
	$$
\end{lem}
\begin{pf}
Let $X=\mathcal{B}(\mathcal{H})$. We shall first prove the left hand inequality. Let $r\in \mathbb{R}^n_{\geq 0}$ be such that for all $f \in \mathcal{PH}(\Omega_2,X)$ of the form \eqref{e-1.3-a},
\begin{equation} \label{ee-1.8-atm-added}
	\sum_{m=0}^{\infty}	\sum_{|\alpha|=m} (\norm{U(a_{\alpha})}^p_{Y} + \norm{U(b_{\alpha})}^p_{Y})r^{p \alpha} \leq \lambda^{p} \norm{f}^p_{\Omega_2, X}.	
\end{equation}
Then, for every $\epsilon >0$, we obtain
\begin{align*}
&\sum_{m=0}^{\infty}	\sum_{|\alpha|=m} (\norm{U(a_{\alpha})}^p_{Y} + \norm{U(b_{\alpha})}^p_{Y}) \left(\frac{r}{S(\Omega_{2},\Omega_{1}) +\epsilon}\right)^{p \alpha} \\
& \leq \lambda^p \norm{\sum_{m=0}^{\infty} \sum_{|\alpha|=m} a_{\alpha}\, \left(\frac{z}{S(\Omega_{2},\Omega_{1}) +\epsilon}\right)^{\alpha} + \sum_{m=1}^{\infty} \sum_{|\alpha|=m} b^{*}_{\alpha}\, \bar{\left(\frac{z}{S(\Omega_{2},\Omega_{1}) +\epsilon}\right)}^{\alpha}}_{\Omega_2,X} \\
& \leq \norm{f}_{\Omega_{1},X}.
\end{align*} 
This implies that
\begin{equation*}
\frac{1}{S(\Omega_{2},\Omega_{1}) +\epsilon} \, \frac{\sum_{i=1}^{n} r_i}{n}=\frac{1}{n}\left(\frac{r_1}{S(\Omega_{2},\Omega_{1}) +\epsilon} +\cdots + \frac{r_n}{S(\Omega_{2},\Omega_{1}) +\epsilon}\right) \leq AP_{\lambda}(\Omega_{1}, p,U)
\end{equation*}
for all $\epsilon>0$.
Consequently, 
\begin{equation*}
\frac{1}{S(\Omega_{2},\Omega_{1})}.\, AP_{\lambda}(\Omega_{2}, p,U) \leq  AP_{\lambda}(\Omega_{1}, p,U). 
\end{equation*}
By interchanging $\Omega_{1}$ and $\Omega_{2}$ in the above inequality, we establish the right-hand inequality.
\end{pf}

	We recall the following remarkable result by Maurey and Pisier to prove the desired upper bound in the case of infintedimensional complex Banach space. 
	\begin{customthm}{A} \cite[Theorem 14.5]{diestel-abs-summing-1995} \label{thm-A}
		Given any infinite dimensional complex Banach space $X$, there exist $x_{1}, \ldots , x_{n} \in X$ for each $n \in \mathbb{N}$ such that $\norm{z}_{\infty}/2 \leq \norm{\sum_{j=1}^{n} x_{j}z_{j}} \leq \norm{z}_{Cot(X)}$  for every choice of $z=(z_{1}, \ldots,z_{n}) \in \mathbb{C}^n$. Clearly, setting $z=e_{j}$, gives $\norm{x_{j}}\geq 1/2$, where $e_{j}$ is the $j$-th canonical basis vector of $\mathbb{C}^n$.
	\end{customthm}
	\section{Proof of the main results}
	%We now begin the proofs of the main results. Let $f(z)=\sum_{|\alpha|=m} c_{\alpha}\, z^{\alpha}$ be an $X$ valued function. For any $p \geq 1$, we write $M_{f,p}:=\sum_{|\alpha|=m}\norm{c_{\alpha}\, z^{\alpha}}^p$, which will be used throughout the proofs.
	\begin{pf} [{\bf Proof of Theorem \ref{thm-1.8-atm}}]
		%We first want to prove that 
		%$$
		%AP_\lambda(B_{Z},p,U)\geq  \frac{\norm{Id: Z \rightarrow \ell^n_{1}}}{n} \, R_\lambda(B_{Z},p,U).
		%$$
		Since $\norm{Id: Z \rightarrow \ell^n_{1}}=\sup_{z \in B_{Z}} \norm{z}_{1}$, for any $0< \epsilon < R_\lambda(B_{Z},p,U)$, we can find an element $\tilde{z}\in B_{Z}$ such that 
		\begin{equation*}
			\norm{\tilde{z}}_{\ell^n_1}\geq \norm{Id: Z \rightarrow \ell^n_{1}}-\epsilon.
		\end{equation*}
		Let $s:=R_\lambda(B_{Z},p,U)-\epsilon$, $w:=s\tilde{z}$, and $r:=s|\tilde{z}|=|w|.$ Since $w\in s\, B_{Z}$ and $s<R_\lambda(B_{Z},p,U)$, for $f \in \mathcal{PH}(B_{z},\mathcal{B}(\mathcal{H}))$ be of the form \eqref{e-1.3-a}, we have
		\begin{align*}
			\sum_{m=0}^{\infty} \sum_{|\alpha|=m} (\norm{U(a_{\alpha})}^p_{Y} + \norm{U(b_{\alpha})}^p_{Y})r^{p \alpha} &=\sum_{m=0}^{\infty} \sum_{|\alpha|=m} (\norm{U(a_{\alpha})}^p_{Y} + \norm{U(b_{\alpha})}^p_{Y})|w^\alpha|^p \\
			&\leq \sup_{z \in s\, B_{Z}}\sum_{m=0}^{\infty} \sum_{|\alpha|=m} (\norm{U(a_{\alpha})}^p_{Y} + \norm{U(b_{\alpha})}^p_{Y})|z^\alpha|^p \\
			&\leq \lambda^{p} \norm{f}^p_{B_{Z}, \mathcal{B}(\mathcal{H})}.
		\end{align*}
		Therefore, we obtain 
		\begin{equation*}
			AP_\lambda(B_{Z},p,U) \geq \frac{1}{n}\sum_{i=1}^{n}r_i = \frac{R_\lambda(B_{Z},p,U)-\epsilon}{n}\norm{\tilde{z}}_{\ell^n_{1}}\geq \frac{R_\lambda(B_{Z},p,U)-\epsilon}{n}\left(\norm{Id: Z \rightarrow \ell^n_{1}}-\epsilon\right)
		\end{equation*}
		holds for all $\epsilon>0$, and so \eqref{e-1.1-added} immediately follows. 
		On the other hand, when $\norm{U}< \lambda$, the following lower bound of $R_\lambda(B_{Z},p,U)$ was established in \cite[Theorem 1.1]{Himadri-local-Banach-1}:
		\begin{equation*}
			R_{\lambda}(B_{Z}, p,U) \geq D. \frac{1}{\sup_{z \in B_{Z}}\norm{z}_{p}},
		\end{equation*}
	where $D$ as in the statement of Theorem \ref{thm-1.8-atm}.
		 Consequently, the second part of the theorem follows from \eqref{e-1.1-added} together with the preceding bound. This completes the proof.	
		\comment{We now only need to show that
			\begin{equation}\label{Pal-Vasu-P3-e-2.13}
				AP_\lambda(B_{Z},p,U) \leq \frac{\norm{Id: Z \rightarrow \ell^n_{1}}}{n} \, R_\lambda(B_{Z},p,U).
			\end{equation}
			Let $r=(r_1,.\,.\,., r_n)\in \mathbb{R}^n_{\geq 0}$ be such that 
			\begin{equation*}
				\sum_{m=0}^{\infty} \sum_{|\alpha|=m} (\norm{U(c_{\alpha})}^p_{Y} + \norm{U(d_{\alpha})}^p_{Y})r^{p \alpha} \leq \lambda^{p} \norm{g}^p_{B_{Z}, \mathcal{B}(\mathcal{H})}	
			\end{equation*}
			for all $g(z)=\sum_{m=0}^{\infty} \sum_{|\alpha|=m} c_{\alpha}\, z^{\alpha} + \sum_{m=1}^{\infty} \sum_{|\alpha|=m} d^{*}_{\alpha}\, \bar{z}^{\alpha}\in \mathcal{PH}(B_{z},\mathcal{B}(\mathcal{H}))$. Then, there exists $\tilde{z}=(z_1,\ldots,z_n)\in R_\lambda(B_{Z},p,U)\overline{B_{Z}}$ with $|z_j|=r_j$ for $1\leq j\leq n$ satisfying $\norm{\tilde{z}}=\norm{r}\leq R_\lambda(B_{Z},p,U)$ and
			\begin{equation*}
				\sum_{m=0}^{\infty} \sum_{|\alpha|=m} (\norm{U(c_{\alpha})}^p_{Y} + \norm{U(d_{\alpha})}^p_{Y})|\tilde{z}|^{p \alpha} \leq \lambda^{p} \norm{g}^p_{B_{Z}, \mathcal{B}(\mathcal{H})}.
			\end{equation*}
			To prove \eqref{Pal-Vasu-P3-e-2.13}, it suffices to show that $	\norm{r}_{1}/\norm{Id: Z \rightarrow \ell^n_{1}} \leq R_\lambda(B_{Z},p,U)$.
			Fix $f\in \mathcal{PH}(B_{z},\mathcal{B}(\mathcal{H}))$ be of the form \eqref{e-1.3-a}, then for $v\in \frac{\norm{r}_1}{\norm{Id: Z \rightarrow \ell^n_{1}}}\overline{B_{Z}}$, we have $\norm{v}\leq \norm{r}\leq R_\lambda(B_{Z},p,U)$.
			Hence, by the virtue of the definition of Bohr radius we obtain
			\begin{equation*}
				\sum_{m=0}^{\infty} \sum_{|\alpha|=m} (\norm{U(a_{\alpha})}^p_{Y} + \norm{U(b_{\alpha})}^p_{Y})|v|^{p \alpha} \leq \lambda^{p} \norm{g}^p_{B_{Z}, \mathcal{B}(\mathcal{H})}
			\end{equation*}
			for every $v\in \frac{\norm{r}_1}{\norm{Id: Z \rightarrow \ell^n_{1}}}\overline{B_{Z}}$. Therefore, it follows that
			$\norm{r}_1 \leq \norm{Id: Z \rightarrow \ell^n_{1}} \, R_\lambda(B_{Z},p,U)$, 
			which gives our conclusion.  This completes the proof.}
	\end{pf}
\begin{pf} [{\bf Proof of Theorem \ref{tthm-1.8-atm}}]
	We shall first prove the left-hand inequality in $(1)$. Let $r=(r_1,.\,.\,., r_n)\in \mathbb{R}^n_{\geq 0}$ be such that 
	\begin{equation*}
		\sum_{|\alpha|=m} (\norm{U(c_{\alpha})}^p_{Y} + \norm{U(d_{\alpha})}^p_{Y})r^{p \alpha} \leq \lambda^{p} \norm{g}^p_{\Omega, X}	
	\end{equation*}
	for all $g(z)=\sum_{|\alpha|=m} c_{\alpha}\, z^{\alpha} +  \sum_{|\alpha|=m} d^{*}_{\alpha}\, \bar{z}^{\alpha}\in \mathcal{PH}(^m\Omega,X)$. Fix $f\in \mathcal{PH}(\Omega,X)$ be of the form \eqref{e-1.3-a}. Let $R \in (0,1)$. Then 
	\begin{align} \label{ee-1.5-atm}
		&\sum_{m=0}^{\infty} \sum_{|\alpha|=m} (\norm{U(a_{\alpha})}^p_{Y} + \norm{U(b_{\alpha})}^p_{Y})(Rr)^{p \alpha}  \nonumber\\
		&= \norm{U(a_{0})}^p +  \sum_{m=1}^{\infty} R^{pm}\sum_{|\alpha|=m} (\norm{U(a_{\alpha})}^p_{Y} + \norm{U(b_{\alpha})}^p_{Y})r^{p \alpha} \nonumber\\ 
		& \leq \norm{U}^p \norm{a_{0}}^p +  \sum_{m=1}^{\infty} R^{pm} \lambda^{p} \norm{\sum_{|\alpha|=m} a_{\alpha}\, z^{\alpha} +  \sum_{|\alpha|=m} b^{*}_{\alpha}\, \bar{z}^{\alpha}}^p_{\Omega,X} \\
		&\leq \norm{U}^p \norm{f}^p _{\Omega,X} + \sum_{m=1}^{\infty} R^{pm} \lambda^{p} \norm{f}^p _{\Omega,X}= \left(\norm{U}^p + \lambda^p\, \frac{R^p}{1-R^p}\right) \norm{f}^p _{\Omega,X} \leq \lambda^p \norm{f}^p _{\Omega,X} ,\nonumber
	\end{align}
	provided 
	$$R \leq \left(\frac{\lambda^p - \norm{U}^p}{2\lambda^p - \norm{U}^p}\right)^{\frac{1}{p}}.$$ This yields
	\begin{equation*}
		\left(\frac{\lambda^p - \norm{U}^p}{2\lambda^p - \norm{U}^p}\right)^{\frac{1}{p}} \, \frac{1}{n} \, \sum_{i=1}^{n} r_{i} \leq AP_{\lambda}(B_{Z}, p,U),
	\end{equation*}
	and hence, the desired bound follows. The uper bound of the right-hand inequality of $(1)$ directly follows from the fact $\cup^{\infty}_{m=1} \mathcal{PH}(^m \Omega,X) \subseteq \mathcal{PH}(\Omega,X)$.
	\vspace{1mm}
	
	 We next prove the left-hand inequality in $(2)$. Let $r=(r_1,.\,.\,., r_n)\in \mathbb{R}^n_{\geq 0}$ be such that 
	\begin{equation*}
		\sum_{|\alpha|=m} (\norm{U(c_{\alpha})}^p_{Y} + \norm{U(d_{\alpha})}^p_{Y})r^{p \alpha} \leq \norm{g}^p_{\Omega, X}	
	\end{equation*}
	for all $g(z)=\sum_{|\alpha|=m} c_{\alpha}\, z^{\alpha} +  \sum_{|\alpha|=m} d^{*}_{\alpha}\, \bar{z}^{\alpha}\in \mathcal{PH}(^m\Omega,X)$. Fix $f\in \mathcal{PH}(\Omega,X)$ be of the form \eqref{e-1.3-a}. Let $R \in (0,1)$. Then 
	\begin{align} \label{ee-1.6-atm}
		&\sum_{m=0}^{\infty} \sum_{|\alpha|=m} (\norm{U(a_{\alpha})}^p_{Y} + \norm{U(b_{\alpha})}^p_{Y})(Rr)^{p \alpha}  \nonumber\\
		&= \norm{U(a_{0})}^p +  \sum_{m=1}^{\infty} R^{pm}\sum_{|\alpha|=m} (\norm{U(a_{\alpha})}^p_{Y} + \norm{U(b_{\alpha})}^p_{Y})r^{p \alpha} \nonumber\\ 
		& \leq \norm{U}^p \norm{a_{0}}^p +  \sum_{m=1}^{\infty} R^{pm}  \norm{\sum_{|\alpha|=m} a_{\alpha}\, z^{\alpha} +  \sum_{|\alpha|=m} b^{*}_{\alpha}\, \bar{z}^{\alpha}}^p_{\Omega,X} \\
		&\leq \norm{U}^p \norm{f}^p _{\Omega,X} + \sum_{m=1}^{\infty} R^{pm} \norm{f}^p _{\Omega,X}= \left(\norm{U}^p + \frac{R^p}{1-R^p}\right) \norm{f}^p _{\Omega,X} \leq  \lambda^p \norm{f}^p _{\Omega,X} ,\nonumber
	\end{align}
	provided 
	$$R \leq \left(\frac{\lambda^p - \norm{U}^p}{\lambda^p - \norm{U}^p +1}\right)^{\frac{1}{p}}.$$ This yields
	\begin{equation*}
		\left(\frac{\lambda^p - \norm{U}^p}{\lambda^p - \norm{U}^p+1}\right)^{\frac{1}{p}} \, \frac{1}{n} \, \sum_{i=1}^{n} r_{i} \leq AP_{\lambda}(B_{Z}, p,U),
	\end{equation*}
	and hence, the desired bound follows. The uper bound of the right-hand inequality in $(2)$ directly follows from the fact $\cup^{\infty}_{m=1} \mathcal{PH}(^m \Omega,X) \subseteq \mathcal{PH}(\Omega,X)$.
\end{pf}

\begin{pf} [{\bf Proof of Theorem \ref{tthm-1.9-atm}}]
	We first prove the inequality $(1)$.  By going the similar line of arguments as in the proof of Theorem \ref{tthm-1.8-atm}, from \eqref{ee-1.5-atm}, we obtain 
	\begin{equation*}
		\sum_{m=0}^{\infty} \sum_{|\alpha|=m} (\norm{U(a_{\alpha})}^p_{Y} + \norm{U(b_{\alpha})}^p_{Y})(Rr)^{p \alpha} \leq \norm{U}^p \norm{a_{0}}^p +  \sum_{m=1}^{\infty} R^{pm} \lambda^{p} \norm{\sum_{|\alpha|=m} a_{\alpha}\, z^{\alpha} +  \sum_{|\alpha|=m} b^{*}_{\alpha}\, \bar{z}^{\alpha}}^p_{\Omega,X}.
	\end{equation*} 
	Now, Lemma \ref{lem-3.1} gives
	\begin{align*}
		\norm{\sum_{|\alpha|=m} a_{\alpha} z^{\alpha}}
		& \leq \norm{\sum_{|\alpha|=m} \left(\frac{a_{\alpha} + b_{\alpha}}{2}\right)z^{\alpha}} + \norm{\sum_{|\alpha|=m} \left(\frac{a_{\alpha} - b_{\alpha}}{2}\right)z^{\alpha}} \nonumber \\
		& \leq 4 \, \norm{\norm{f}_{\Omega,X}\,I-\real(a_{0})}.
	\end{align*}
	The last inequality remains valid if the left hand quantity replaced by $\norm{\sum_{|\alpha|=m} b_{\alpha} z^{\alpha}}$. Thus, 
	\begin{equation} \label{ee-1.7-atm}
		\norm{\sum_{|\alpha|=m} a_{\alpha}\, z^{\alpha} +  \sum_{|\alpha|=m} b^{*}_{\alpha}\, \bar{z}^{\alpha}}^p \leq 8^p \, \norm{\norm{f}_{\Omega,X}\,I-\real(a_{0})}^p.
	\end{equation}
	Using \eqref{ee-1.7-atm} and the fact $f(0)=a_0=\mu I$, $0 \leq \mu \leq \norm{f}_{\Omega,X}$, a simple computation shows that 
	\begin{equation*}
		\sum_{m=0}^{\infty} \sum_{|\alpha|=m} (\norm{U(a_{\alpha})}^p_{Y} + \norm{U(b_{\alpha})}^p_{Y})(Rr)^{p \alpha} \leq \lambda^p \mu^p + \lambda^p 8^p \left(\norm{f}^p_{\Omega,X} - \mu^p \right)\frac{R^p}{1-R^p} \leq \lambda^p \norm{f}^p_{\Omega,X},
	\end{equation*}
	provided 
	$$R \leq \frac{1}{(1+8^p)^{1/p}}.$$ This yields
	$$ \frac{1}{(1+8^p)^{1/p}}\,\frac{1}{n} \, \sum_{i=1}^{n} r_{i} \leq AP^{*}_{\lambda}(B_{Z}, p,U),$$
	and hence, the desired bound follows. 
	\vspace{1mm}
	
	By going the similar line of arguments as in the proof of part $(2)$ of Theorem \ref{tthm-1.8-atm}, using  \eqref{ee-1.6-atm} and \eqref{ee-1.7-atm}, we obtain
	\begin{equation*}
		\sum_{m=0}^{\infty} \sum_{|\alpha|=m} (\norm{U(a_{\alpha})}^p_{Y} + \norm{U(b_{\alpha})}^p_{Y})(Rr)^{p \alpha} \leq \lambda^p \mu^p + 8^p \left(\norm{f}^p_{\Omega,X} - \mu^p \right)\frac{R^p}{1-R^p} \leq \lambda^p \norm{f}^p_{\Omega,X},
	\end{equation*}
	provided 
	$$R \leq \frac{\lambda}{(\lambda^p+8^p)^{1/p}}.$$ This yields
	$$ \frac{\lambda}{(\lambda^p+8^p)^{1/p}}\,\frac{1}{n} \, \sum_{i=1}^{n} r_{i} \leq AP^{*}_{\lambda}(B_{Z}, p,U),$$
	and hence, the desired bound follows. The proof is now complete.
\end{pf}

\begin{pf}[{\bf Proof of Theorem \ref{thm-1.2-atm}}]
Let $X=\mathcal{B}(\mathcal{H})$. We shall first prove the lower bound of $AP_{\lambda}(B_{Z}, p,X)$. The following lower bound of $R_{\lambda}(B_{Z}, p,X)$ was established in \cite[Theorem 1.2]{Himadri-local-Banach-1} in the case where $X$ is finite dimensional:
\begin{equation*}
	R_{\lambda}(B_{Z}, p,X) \geq	\begin{cases}
		E_{1}(X)\,  \max \left\{\frac{1}{\sqrt{n}\, \norm{Id: \ell^n_{2} \rightarrow Z}}, \, \frac{1}{e\, \norm{Id: Z \rightarrow \ell^n_{1}}}\right\}\, \left(\frac{\lambda^p -1}{2\lambda^p - 1}\right)^{\frac{1}{p}} & \text{for $p=1$}, \\[2mm]
		E_{3}\norm{I:\ell^n _2 \rightarrow Z}^{-\frac{2}{p}}\, \left(\frac{\lambda^p -1}{2\lambda^p - 1}\right)^{\frac{1}{p}}\, & \text{for $p \geq 2$}, \\[2mm]
		
		E_{2}(X)\,  \max \left\{\frac{1}{(\sqrt{n})^{1-\theta} \norm{Id: \ell^n_{2} \rightarrow Z}},  \frac{\norm{I:\ell^n _2 \rightarrow Z}^{-\theta}}{e^{(1-\theta)} \norm{Id: Z \rightarrow \ell^n_{1}}^{1-\theta}}\right\}\left(\frac{\lambda^p -1}{2\lambda^p - 1}\right)^{\frac{1}{p}} \hspace{-2mm} & \text{for $1<p< 2$}, 
	\end{cases}
\end{equation*}
where $\theta = (2(p-1))/p$. Here, $E_{1}(X)$, $E_{2}(X)$, and $E_{3}$ are positive constants: $E_{1}(X)$ depends only on $X$, $E_{2}(X)$ depends on both $X$ and $p$, while $E_{3}$ is independent of $X$ and depends only on $p$. Thus, by combining Theorem \ref{thm-1.8-atm} with the above estimate of $R_{\lambda}(B_{Z}, p,X)$, the desired lower bound of $AP_{\lambda}(B_{Z}, p,X)$ follows. 
\vspace{1mm}

We now proceed to estimate the upper bound. In this regard, we first obtain a lower bound of $AP_{\lambda}(B_{Z}, p,X)$ when $Z=\ell^n_q$, $1 \leq q \leq \infty$, and then the desired lower bound of $AP_{\lambda}(B_{Z}, p,X)$ for general Banach space $Z$ follows from Lemma \ref{lem-3.6}. Let $r=(r_1,.\,.\,., r_n)\in \mathbb{R}^n_{\geq 0}$ be such that 
\begin{equation*}
	\sum_{|\alpha|=m} (\norm{c_{\alpha}}^p + \norm{d_{\alpha}}^p)r^{p \alpha} \leq \lambda^p \, \norm{g}^p_{B_{\ell^n_q}, X}	
\end{equation*}
for all $g(z)=\sum_{|\alpha|=m} c_{\alpha}\, z^{\alpha} +  \sum_{|\alpha|=m} d^{*}_{\alpha}\, \bar{z}^{\alpha}\in \mathcal{PH}(^mB_{\ell^n_q},X)$. Then, for $\varepsilon_{\alpha} \in \{-1,1\}$ and $I$ the identity operator on $\mathcal{H}$, the $m$-homogeneous polynomial
\begin{equation*}
f(z)=\sum_{|\alpha|=m} \varepsilon_{\alpha}\, \frac{m!}{\alpha!}z^{\alpha}\, I +  \sum_{|\alpha|=m} \varepsilon_{\alpha}\,\frac{m!}{\alpha!}\, \bar{z}^{\alpha} I, \,\, z \in B_{\ell^n_q},
\end{equation*} 
satisfies
\begin{equation} \label{ee-5.4-added}
2\, \sum_{|\alpha|=m}  \left(\frac{m!}{\alpha!}\right)^p\, r^{p\alpha} \leq \lambda^p \,\norm{f}^p_{B_{\ell^n_q},X}.
\end{equation}
Using the facts $\norm{r}_1 \leq n^{1-\frac{1}{p}}\, \norm{r}_p$ and $m! /\alpha! \geq 1$, a simple computation shows that
\begin{align*}
\left(\sum_{i=1}^{n}r_i\right)^{pm} 
& \leq \left(n^{1-\frac{1}{p}}\right)^{pm}\, \left(\sum_{i=1}^{n}r^p_i\right)^m \nonumber \\
& = \left(n^{1-\frac{1}{p}}\right)^{pm}\, \sum_{|\alpha|=m} \frac{m!}{\alpha!}\, r^{p \alpha} \nonumber \\
& \leq \left(n^{1-\frac{1}{p}}\right)^{pm}\, \sum_{|\alpha|=m} \left(\frac{m!}{\alpha!}\right)^p\, r^{p \alpha}.
\end{align*}
In view of the preceding inequality, from \eqref{ee-5.4-added}, we deduce that
\begin{align} \label{ee-5.5-added}
2\, \left(\sum_{i=1}^{n}r_i\right)^{pm} 
& \leq \lambda^p\, \left(n^{1-\frac{1}{p}}\right)^{pm}\, \norm{\sum_{|\alpha|=m} \varepsilon_{\alpha}\, \frac{m!}{\alpha!}z^{\alpha}\, I +  \sum_{|\alpha|=m} \varepsilon_{\alpha}\,\frac{m!}{\alpha!}\, \bar{z}^{\alpha} I}^p_{B_{\ell^n_q},X} \nonumber\\
& \leq \lambda^p\, \left(n^{1-\frac{1}{p}}\right)^{pm}\, \left(\norm{\sum_{|\alpha|=m} \varepsilon_{\alpha}\, \frac{m!}{\alpha!}z^{\alpha}\, I}_{B_{\ell^n_q},X} + \norm{\sum_{|\alpha|=m} \varepsilon_{\alpha}\,\frac{m!}{\alpha!}\, \bar{z}^{\alpha} I}_{B_{\ell^n_q},X}\right)^p \nonumber \\
& = \lambda^p\, \left(n^{1-\frac{1}{p}}\right)^{pm}\, \left(\norm{\sum_{|\alpha|=m} \varepsilon_{\alpha}\, \frac{m!}{\alpha!}z^{\alpha}}_{B_{\ell^n_q},\mathbb{C}} + \norm{\sum_{|\alpha|=m} \varepsilon_{\alpha}\,\frac{m!}{\alpha!}\, \bar{z}^{\alpha}}_{B_{\ell^n_q},\mathbb{C}}\right)^p \nonumber\\
&= \lambda^p\, \left(n^{1-\frac{1}{p}}\right)^{pm}\, \left(\norm{\sum_{|\alpha|=m} \varepsilon_{\alpha}\, \frac{m!}{\alpha!}z^{\alpha}}_{B_{\ell^n_q},\mathbb{C}} + \norm{\sum_{|\alpha|=m} \varepsilon_{\alpha}\,\frac{m!}{\alpha!}\, z^{\alpha} }_{B_{\ell^n_q},\mathbb{C}}\right)^p \nonumber\\
&= 2^p\, \lambda^p\, \left(n^{1-\frac{1}{p}}\right)^{pm}\,\norm{\sum_{|\alpha|=m} \varepsilon_{\alpha}\,\frac{m!}{\alpha!}\, z^{\alpha} }^p_{B_{\ell^n_q},\mathbb{C}}.
\end{align}
We now recall the following probabilistic estimate due to Boas \cite[Theorem 4]{boas-2000}:
\begin{equation*}
\norm{\sum_{|\alpha|=m} \varepsilon_{\alpha}\,\frac{m!}{\alpha!}\, z^{\alpha} }_{B_{\ell^n_q},\mathbb{C}} \leq 
\gamma(m,n,q):= \sqrt{32m\, \log(6m)}
\begin{cases}
n^{\frac{1}{2}}\, (m!)^{1-\frac{1}{q}}\, & \mbox{if}\,\, q \leq 2 \\
n^{\frac{1}{2}+ \left(\frac{1}{2}-\frac{1}{q}\right)m}\, (m!)^{\frac{1}{2}}\, & \mbox{if}\,\, q \geq 2.
\end{cases}
\end{equation*}
This estimate, together with \eqref{ee-5.5-added}, yields
\begin{equation*}
	\frac{\sum_{i=1}^{n}r_i}{n} \leq \left(2^{1-\frac{1}{p}}\right)^m \lambda^{\frac{1}{m}}\, n^{-\frac{1}{p}}\, (\gamma(m,n,q))^{\frac{1}{m}}.
\end{equation*}
Consequently,
\begin{align} \label{ee-5.6-added}
AP^m_{\lambda}(B_{\ell^n_q}, p,X) 
& \leq \left(2^{1-\frac{1}{p}}\right)^m \lambda^{\frac{1}{m}}\, n^{-\frac{1}{p}}\, (\gamma(m,n,q))^{\frac{1}{m}} \nonumber\\
& \leq C_{1}\,\lambda^{\frac{1}{m}}\, \left(m\, \log\,m (m!)^{2 \left(1-\frac{1}{\min\{2,q\}}\right)}\,n\right)^{\frac{1}{2m}}\, \left(\frac{1}{n}\right)^{\frac{1}{p}+\frac{1}{\max\{2,q\}}-\frac{1}{2}} 
\end{align}
for $C_1>0$, a uniform constant. We already know that $AP_{\lambda}(B_{\ell^n_q}, p,X) \leq AP^m_{\lambda}(B_{\ell^n_q}, p,X)$, so for the desired upper bound of $AP_{\lambda}(B_{\ell^n_q}, p,X) $, it is enough to estimate the right side quantity of  \eqref{ee-5.6-added}. Write $t=1-1/\min\{2,q\}$ for short. By Stirling's formula, we have $m! \sim \sqrt{2\, \pi m} (m/e)^m$. Hence
\begin{equation*}
(m!)^{2t} = \left(\frac{m}{e}\right)^{2tm} (2 \pi m)^{t} \left(1+O\left(\frac{1}{m}\right)\right),
\end{equation*}
and so the right side quantity of \eqref{ee-5.6-added} becomes less than or equals to
\begin{equation} \label{ee-5.7-added}
C_1 \lambda^{\frac{1}{m}}\, (m\, \log\,m \, n)^{\frac{1}{2m}}\, \left(\frac{m}{e}\right)^{t} (2 \pi m)^{\frac{t}{2m}} \left(1+O\left(\frac{1}{m}\right)\right)^{\frac{1}{2m}}\, \, \left(\frac{1}{n}\right)^{\frac{1}{p}+\frac{1}{\max\{2,q\}}-\frac{1}{2}}.
\end{equation}
Since $(2 \pi m)^{\frac{t}{2m}} \rightarrow 1$ as $m \rightarrow \infty$, the sequence $\{(2 \pi m)^{\frac{t}{2m}}\}$ is bounded. On the other hand, $(1+O(1/m))^{\frac{1}{2m}}$ is also bounded by uniform constant. We now want to estimate the quantity $(m\, \log\,m \, n)^{\frac{1}{2m}}$. Take $m=[\log\,n]$. Then
\begin{equation*}
\log\, (m\, \log\,m \, n)=\log\, n + \log\, m +\log\log\,m= \log \,n + O(\log\log \,n).
\end{equation*}
Hence,
\begin{equation*}
(m\, \log\,m \, n)^{\frac{1}{2m}}= \exp \left(\frac{\log \,n + O(\log\log \,n)}{2m}\right).
\end{equation*}
Since $m \sim \log \, n$,
\begin{equation*}
\frac{\log \,n}{2m} \leq \frac{1}{2} + o(1)\,\, \mbox{and} \, \, O\left(\frac{\log\log \,n}{\log \, n}\right) \rightarrow 0.
\end{equation*}
Therefore,
\begin{equation*}
(m\, \log\,m \, n)^{\frac{1}{2m}} \leq C_2	
\end{equation*}
for some constant $C_2>0$. Combining these estimates, the quantity in \eqref{ee-5.7-added} becomes less than or equals to
\begin{equation*}
C_{3}\, \lambda^{\frac{1}{m}}\, m^t\,	\left(\frac{1}{n}\right)^{\frac{1}{p}+\frac{1}{\max\{2,q\}}-\frac{1}{2}}, 
\end{equation*}
which finally less than or equals to
\begin{equation*}
C\, \lambda^{\frac{1}{\log\,n}}\,\, (\log\,n)^{1-\frac{1}{\min\{2,q\}}}\,\,\left(\frac{1}{n}\right)^{\frac{1}{p}+\frac{1}{\max\{2,q\}}-\frac{1}{2}}
\end{equation*}
for some universal constant $C>0$.
Thus,
\begin{equation} \label{ee-5.8-added}
AP_{\lambda}(B_{\ell^n_q}, p,X) \leq C\, \lambda^{\frac{1}{\log\,n}}\,\, (\log\,n)^{1-\frac{1}{\min\{2,q\}}}\,\,\left(\frac{1}{n}\right)^{\frac{1}{p}+\frac{1}{\max\{2,q\}}-\frac{1}{2}}.
\end{equation}
For the upper bound of $AP_{\lambda}(B_{Z}, p,X)$, by making use of Lemma \ref{lem-3.6} and \eqref{ee-5.8-added}, we establish
\begin{equation*}
AP_{\lambda}(B_{Z}, p,X) \leq C\, \lambda^{\frac{1}{\log\,n}}\, \norm{Id:Z \rightarrow \ell^n_q}\, (\log\,n)^{1-\frac{1}{\min\{2,q\}}}\,\,\left(\frac{1}{n}\right)^{\frac{1}{p}+\frac{1}{\max\{2,q\}}-\frac{1}{2}}.
\end{equation*}
This completes the proof.
\end{pf}
\begin{pf} [{\bf Proof of Corollary \ref{cor-1.4}}]
Let $X=\mathcal{B}(\mathcal{H})$. The following lower bound of $R_{\lambda}(B_{\ell^n_q}, p,X)$ was established in \cite[Corollary 1.2]{Himadri-local-Banach-1} in the case where $X$ is finite dimensional:
\begin{equation*}
	R_{\lambda}(B_{\ell^n _q}, p,X) \geq	\begin{cases}
		E'_{3}(X)\,\left(\frac{\lambda^p -1}{2\lambda^p - 1}\right)^{\frac{1}{p}} \, \left(\frac{\log n}{n}\right)^{ \left(1- \frac{1}{\min\{q,2\}}\right)} & \text{for $p=1$}, \\[2mm]
		E'_{2} \left(\frac{\lambda^p -1}{2\lambda^p - 1}\right)^{\frac{1}{p}}\, n^{-\frac{1}{p}} & \text{for $p \geq q$}, \\[2mm]
		
		E_{4}(X)\,  \left(\frac{\lambda^p -1}{2\lambda^p - 1}\right)^{\frac{1}{p}} \,n^{-\frac{p-1}{p(q-1)}}\, \left( \frac{\log n}{n}\right)^{\left(1- \frac{1}{\min\{q,2\}}\right)\, \frac{q-p}{p(q-1)}} & \text{for $1<p< q$}. 
	\end{cases}
\end{equation*}
Here, $E'_{3}(X)$, $E'_{2}$, and $E_{4}(X)$ are positive constants: $E'_{3}(X)$ depends only on $X$, $E_{4}(X)$ depends on both $X$ and $p$, while $E'_{2}$ is independent of $X$ and depends only on $p$. Thus, by combining Theorem \ref{thm-1.8-atm} with the preceding estimate of $R_{\lambda}(B_{\ell^n_q}, p,X)$, the desired lower bound of $AP_{\lambda}(B_{\ell^n_q}, p,X)$ follows. The proof is now complete.
\end{pf}
\begin{pf} [{\bf Proof of Theorem \ref{thm-1.3-a}}]
	We start by recalling two classical facts from the theory of finite-dimensional Banach spaces with unconditional bases. Let $W=(\mathbb{C}^n, ||.||)$ be a Banach space whose canonical unit vectors $e_{k}$ form a normalized $1$-unconditional basis. Then the operator norm of the identity embedding from $W$ into $\ell^n_2$ is given by $\norm{Id: W \rightarrow \ell_2^n}=\sup_{z \in B_{W}}(\sum_{m=1}^{n}|z_{m}|^2)^{1/2}$. On the other hand, the identity embedding from $Y$ into $\ell^n_1$ satisfies  $\norm{Id: W \rightarrow \ell_1^n}=\sup_{z \in B_{W}} \sum_{m=1}^{n}|z_{m}|=\norm{\sum_{m=1}^{n}e^{*}_{m}}_{W^{*}}$ for all $z \in W$. We are now in a position to establish the assertions of the theorem.
	\vspace{1mm}
	
	Since $Z \subset \ell_2$, we may assume that the inclusion map $Z \subset \ell_2$ has norm $1$. Under this normalization, the claim $(1)$ follows directly from Theorem \ref{thm-1.2-atm} together with the above identities for the identity maps.
	\vspace{1mm}
	
	We shall now prove part $(2)$. By a result of \cite[Theorem 1.d.5]{Lindenstrauss-book}, we may assume without loss of generality that the constant $M^{(2)}(X)=1$. This implies that $\norm{Id: \ell_2^n \rightarrow Z_{n}}=1$. We now recall the following important fact from  \cite[(5.3), p. 191]{defant-2003}. If $W=(\mathbb{C}^n, ||.||)$ be a symmetric Banach space such that $M^{(2)}(W)=1$, then 
	\begin{equation} \label{e-5.4-added}
	\norm{Id: W \rightarrow \ell^n_2}= \frac{\norm{\sum_{m=1}^{n}e^{*}_{m}}_{W^{*}}}{\sqrt{n}}.
	\end{equation}
	 The conclusion then follows by combining Theorem \ref{thm-1.2-atm}, the above facts for identity embeddings, and the inequality \eqref{e-5.4-added}. This concludes the proof.	
\end{pf}
\begin{pf} [{\bf Proof of Theorem \ref{thm-1.3-atm}}]
	Let $X=\mathcal{B}(\mathcal{H})$. The lower bound of $AP_{\lambda}(B_{\ell^n _q}, p,X)$ follows from Theorem \ref{thm-1.8-atm} by observing the fact 
	$$\norm{Id: B_{\ell^n_q} \rightarrow B_{\ell^n_1}}= n^{1-\frac{1}{q}}$$
	for any $1 \leq q \leq \infty$.
	Now, we want to find the upper bound of $AP_{\lambda}(B_{\ell^n _q}, p,X)$. Let $r=(r_1,.\,.\,., r_n)\in \mathbb{R}^n_{\geq 0}$ be such that 
	\begin{equation} \label{e-1.5-atm}
		\sum_{k=1}^{n} (\norm{c_{k}}^p + \norm{d_{k}}^p)r^{p} \leq  \norm{g}^p_{B_{Z}, \mathcal{B}(\mathcal{H})}	
	\end{equation}
	for all $g(z)=\sum_{k=1}^{n}  c_{k}\, z_{k} + \sum_{k=1}^{n} d^{*}_{k}\, \bar{z}_{k}\in \mathcal{PH}(^1B_{z},\mathcal{B}(\mathcal{H}))$. We shall now make use of Theorem \ref{thm-A}. By considering $z=e_{j}$ in Theorem \ref{thm-A}, for each $n \in \mathbb{N}$ there exist $a_{1}, \ldots, a_{n} \in X$ such that $1/2 \leq \norm{a_{j}} $ for each $1\leq j \leq n$. This inequality together with  H\"{o}lder's inequality, by considering the $1$-homogeneous polynomial $F \in \mathcal{PH}(^1B_{Z},X)$ of the form $F(z)=\sum_{j=1}^{n} z_{j}a_{j}+ \sum_{j=1}^{n} a^{*}_{j}\bar{z_{j}}$, yields
	\begin{align*} \label{e-1.39-atm}
		\sum_{k=1}^{n} r_{k}  \leq \sum_{k=1}^{n}(\norm{a_{j}}+\norm{a_{j}}) r_{k} \leq \left(\sum_{k=1}^{n} (\norm{a_{k}}+\norm{b_{k}})^p r^p_{k}\right)^{\frac{1}{p}} (\sum_{k=1}^{n} 1)^{1-\frac{1}{p}}  = 2^{1-\frac{1}{p}} n^{1 -\frac{1}{p}} \norm{F}_{B_{\ell^n_q},X}.  
	\end{align*} 
	Using Theorem \ref{thm-A}, from the last inequality we obtain 
	\begin{equation} \label{e-1.6-atm}
		\frac{1}{n} \, \sum_{k=1}^{n} r_{k} \leq 2^{1-\frac{1}{p}} n^{ -\frac{1}{p}} \sup_{z \in B_{\ell^n _q}} \norm{z}_{Cot(X)}.
	\end{equation}
	For $q \leq Cot(X)$, we have $\sup_{z \in B_{\ell^n _q}} \norm{z}_{Cot(X)} \leq 1$, which together with \eqref{e-1.6-atm} gives 
	$$AP^1(B_{\ell^n _q}, p,X)\leq 2^{1-\frac{1}{p}} n^{ -\frac{1}{p}}.$$ On the other hand, for $q>Cot(X)$, $$\sup_{z \in B_{\ell^n _q}} \norm{z}_{Cot(X)} \leq n^{1/Cot(X)\, - \,1/q},$$ and so from \eqref{e-1.6-atm} we obtain 
	$$AP^1(B_{\ell^n _q}, p,X) \leq 2^{1-\frac{1}{p}}\, n^{1/Cot(X)\, - \,1/q -1/p}.$$ It is known that $AP_{\lambda}(B_{\ell^n _q}, p,X)\leq  AP^1_{\lambda}(B_{\ell^n _q}, p,X)=\lambda\,AP^1(B_{\ell^n _q}, p,X)$, and hence the upper bound of $AP_{\lambda}(B_{\ell^n _q}, p,X)$ immediately follows. This completes the proof.
\end{pf} 
		
	\begin{pf} [{\bf Proof of Theorem \ref{thm-1.8}}]
	We shall make use of the following result due to \cite[Theorem 5.3]{defant-2012}. If $Y$ is a $s$-concave Banach lattice, with $2\leq s < \infty$, and $U: X \rightarrow Y$ an $(t,1)$-summing operator with $1 \leq t \leq s$. Then there is a constant $C>0$ such that
	\begin{equation} \label{e-5.7-added}
		\left(\sum_{|\alpha|=m} \norm{U(c_{\alpha})}^{\beta}_{Y}\right)^{\frac{1}{\beta}} \leq C^m \, \norm{Q}_{\mathbb{D}^n,X}
	\end{equation}
	holds for every homogeneous polynomial $Q: \mathbb{D}^n \rightarrow X$ with $Q(z)=\sum_{|\alpha|=m} c_{\alpha}\, z^{\alpha}$ on $\mathbb{D}^n$, where 
	\begin{equation*}
		\beta= \frac{mst}{s+ (m-1)t}.
	\end{equation*}	
The proof of the theorem relies on the definition of complete Reinhardt domains and identifying $B_z$ as such a domain.  Recall that a domain $\Omega \subset \mathbb{C}^n$ is called complete Reinhardt if $(z_{1}, \ldots,z_{n}) \in \Omega$ then the new tuple $(\xi_1 z_1, \ldots , \xi _nz_n) \in \Omega$ for each $\xi_j \in \overline{\mathbb{D}}$, $j=1, \ldots,n$.
Here, $Z=(\mathbb{C}^n, ||.||)$ is an $n$-dimensional Banach space for which the canonical basis vectors $e_{k}$ form a normalized $1$-unconditional basis. That is, equivalent to say the open unit ball $B_{Z}$ is a complete Reinhardt domain.
 
Fix $z \in B_{Z}$. For every $m$-homogeneous polynomial $P: B_{Z} \rightarrow X$ with $P(z)=\sum_{|\alpha|=m} a_{\alpha}\, z^{\alpha}$ on $B_{Z}$, we consider the following $m$-homogeneous polynomial $\tilde{Q}: \mathbb{D}^n \rightarrow X$ by 
\begin{equation*}
	\tilde{Q}(w_1, \ldots,w_n):=P(z_1w_1, \ldots , z_nw_n)= \sum_{|\alpha|=m} (a_{\alpha}\, z^{\alpha}) w^{\alpha}
\end{equation*}
for $w=(w_1, \ldots ,w_n) \in \mathbb{D}^n$. The map $\tilde{Q}$ is well defined, since $B_Z$ is a complete Reinhardt domain and hence the point $(z_1w_1, \ldots , z_nw_n)$ belongs to $B_Z$. Clearly, $\norm{\tilde{Q}}_{\mathbb{D}^n,X} \leq \norm{P}_{B_{Z},X}$. Now, by virtue of \eqref{e-5.7-added}, we obtain
\begin{equation*} 
	\left(\sum_{|\alpha|=m} \norm{U(a_{\alpha})z^{\alpha}}^{\beta}_{Y}\right)^{\frac{1}{\beta}} \leq C^m \, \norm{\tilde{Q}}_{\mathbb{D}^n,X} \leq C^m \, \norm{P}_{B_{Z},X}.
\end{equation*}
This completes the proof.
	\end{pf}
	
\begin{pf} [{\bf Proof of Theorem \ref{thm-1.9}}]
We first want to prove the part $(1)$. Let $X=\mathcal{B}(\mathcal{H}_1)$ and $Y=\mathcal{B}(\mathcal{H}_2)$. By the definition of Bohr radius, it not difficult to see that
\begin{equation} \label{ee-5.8}
R_{\lambda}(B_{Z}, p,U) \geq \max\left\{R_{\frac{\lambda}{\norm{U}}}(B_{Z}, p,X), \, \, R_{\frac{\lambda}{\norm{U}}}(B_{Z}, p,Y)\right\}.
\end{equation}
In view of \cite[Theorem 1.5]{Himadri-local-Banach-1} and \eqref{ee-5.8}, for $1\leq p<\infty$ and $1 \leq q \leq \infty $, we deduce that
\begin{align} \label{e-5.9-added}
R_{\lambda}(B_{Z}, p,U) \geq \frac{\Phi_{1}(n,\lambda,p,q)}{S(B_{Z},B_{\ell^n _q})S(B_{\ell^n _q},B_{Z})},
\end{align} 
where
\begin{equation*}
\Phi_{1}(n,\lambda,p,q):=	\begin{cases}
	%	\dfrac{\left(\lambda^p -1\right)^{\frac{1}{p}}}{\lambda}\, 2^{-\left(1+\frac{1}{p}\right)}\, e^{-\frac{1}{p}}\, 4^{-\frac{q}{p}} & \text{for $p > q$}, \\[2mm]
	
	E_{5}\,\dfrac{\left(\lambda^p -\norm{U}^p\right)^{\frac{1}{p}}}{\lambda} & \text{for $p>\min \{Cot(X),Cot(Y),q\}$}, \\[2mm]
	E_{6}\, \dfrac{\left(\lambda^p -\norm{U}^p\right)^{\frac{1}{p}}}{\lambda}\,\, n^{{\frac{1}{\min\{t,q\}}}-\frac{1}{p}} & \text{for $p\leq \min \{Cot(X),Cot(Y),q\}$}.
\end{cases}
\end{equation*}
 Here $E_{5}>0$ and $E_{6}>0$ are constants independent of $X$ and $n$. The lower bound of $AP_{\lambda}(B_{Z}, p,U)$ follows from Theorem \ref{thm-1.8-atm} and \eqref{e-5.9-added}.
 
 We shall now prove the part $(2)$. Now let $X=\mathcal{B}(\mathcal{H})$.
Let 
$$P(z)=\sum_{|\alpha|=m} a_{\alpha}\, z^{\alpha} + \sum_{|\beta|=m} b^{*}_{\alpha}\, \bar{z}^{\alpha} \in \mathcal{PH}(^m Z,X).$$ Consider the holomorphic polynomial $$H(z):=\sum_{|\alpha|=m} \left(\frac{a_{\alpha} + b_{\alpha}}{2}\right)z^{\alpha}$$ on $B_{Z}$. In view of H\"{o}lder's inequality and Theorem \ref{thm-1.8}, we obtain
\begin{align*}
&\sum_{|\alpha|=m} \norm{U\left(\frac{a_{\alpha} + b_{\alpha}}{2}\right)z^{\alpha}}^p \nonumber\\
&\leq  \left(\sum_{|\alpha|=m} \norm{U\left(\frac{a_{\alpha} + b_{\alpha}}{2}\right)z^{\alpha}}\right)^p \nonumber \\
& \leq \left[\bigg(\sum_{|\alpha|=m}1\bigg)^{\frac{(s-1)mt-s+t}{mst}} \times \left(\sum_{|\alpha|=m} \norm{U\left(\frac{a_{\alpha} + b_{\alpha}}{2}\right)z^{\alpha}}^{\frac{mst}{s+(m-1)t}}\right)^{\frac{s+(m-1)t}{mst}}\right]^p \nonumber \\
& \leq \bigg(\sum_{|\alpha|=m}1\bigg)^{\frac{((s-1)mt-s+t)p}{mst}} c_{1}^{mp}\, \norm{H}^p_{\Omega,X} \leq \bigg(\sum_{|\alpha|=m}1\bigg)^{\frac{((s-1)mt-s+t)p}{mst}} c_{1}^{mp}\, 2^p \, \norm{P}^p_{B_{Z},X}.
\end{align*}
Similarly, by considering the holomorphic polynomial $G(z):=\sum_{|\alpha|=m} \left(\frac{a_{\alpha} - b_{\alpha}}{2}\right)z^{\alpha}$ on $B_{Z}$, we deduce that
\begin{equation*}
\sum_{|\alpha|=m} \norm{U\left(\frac{a_{\alpha} - b_{\alpha}}{2}\right)z^{\alpha}}^p \leq \bigg(\sum_{|\alpha|=m}1\bigg)^{\frac{((s-1)mt-s+t)p}{mst}} c_{2}^{mp}\, 2^p \, \norm{P}^p_{B_{Z},X}.
\end{equation*} 
A simple computation from last two inequalities shows that
\begin{align} \label{e-5.8}
&\left(\sum_{|\alpha|=m} \norm{U\left(\frac{a_{\alpha} + b_{\alpha}}{2}\right)z^{\alpha}}^p\right)^{1/p}  \nonumber \\
& \leq \left(\sum_{|\alpha|=m} \norm{U\left(\frac{a_{\alpha} + b_{\alpha}}{2}\right)z^{\alpha}}^p\right)^{1/p} + \left(\sum_{|\alpha|=m} \norm{U\left(\frac{a_{\alpha} - b_{\alpha}}{2}\right)z^{\alpha}}^p\right)^{1/p} \nonumber \\
& \leq 2(c^m_{1}+c^m_2) \bigg(\sum_{|\alpha|=m}1\bigg)^{\frac{(s-1)mt-s+t}{mst}} \norm{P}_{B_{Z},X}
\end{align} 
and 
\begin{equation} \label{e-5.9}
\left(\sum_{|\alpha|=m} \norm{U\left(\frac{a_{\alpha} - b_{\alpha}}{2}\right)z^{\alpha}}^p\right)^{1/p} \leq 2(c^m_{1}+c^m_2) \bigg(\sum_{|\alpha|=m}1\bigg)^{\frac{(s-1)mt-s+t}{mst}} \norm{P}_{B_{Z},X}.
\end{equation}
In view of \eqref{e-5.8} and \eqref{e-5.9}, we obtain 
\begin{equation} \label{e-5.10}
\sum_{|\alpha|=m} \left(\norm{U(a_{\alpha})}^p + \norm{U(b_{\alpha})}^p\right) |z^{\alpha}|^p \leq C^{pm} \bigg(\sum_{|\alpha|=m}1\bigg)^{\frac{((s-1)mt-s+t)p}{mst}} \norm{P}^p_{B_{Z},X}
\end{equation}
for some constant $C>0$. Using the estimate 
$$\sum_{|\alpha|=m} 1 =\binom{n+m-1}{m} \leq e^m \left(1+\frac{n}{m}\right)^m,
$$
from \eqref{e-5.10}, we deduce that 
\begin{equation} \label{e-5.11}
R^m(B_Z,p,U) \geq \tilde{D} \left(1+\frac{n}{m}\right)^{-\frac{(s-1)mt-s+t}{mst}}
\end{equation}
for some constant $\tilde{D}>0$. Since $R^m_{\lambda}(B_Z,p,U)=\lambda^{1/m} \, R^m(B_Z,p,U)$, \eqref{e-5.11} gives
\begin{equation} \label{e-5.12}
R^m_{\lambda}(B_Z,p,U) \geq\lambda^{\frac{1}{m}} \, \tilde{D} \left(1+\frac{n}{m}\right)^{-\frac{(s-1)mt-s+t}{mst}}.
\end{equation}
We now observe the elementary fact if $a,b>0$ and $0<p<1$, then $(a+b)^p \leq 2^p\, \max\{a^p, b^p\}$. Clearly, $$0<\frac{(s-1)mt-s+t}{mst}<1.$$
Thus, we have
\begin{equation*}
\left(1+\frac{n}{m}\right)^{\frac{(s-1)mt-s+t}{mst}} \leq 2^{\frac{(s-1)mt-s+t}{mst}} \max \left\{1, \left(\frac{n}{m}\right)^{\frac{(s-1)mt-s+t}{mst}}\right\} \leq 2\, \max \left\{1, \left(\frac{n}{m}\right)^{\frac{(s-1)m-(\frac{s}{t}-1)}{ms}}\right\}.
\end{equation*}
Consequently,
\begin{equation} \label{e-5.13}
	\left(1+\frac{n}{m}\right)^{-\frac{(s-1)mt-s+t}{mst}} \geq \frac{1}{2}\, \min \left\{1, \left(\frac{n}{m}\right)^{-\frac{(s-1)m-(\frac{s}{t}-1)}{ms}}\right\}=\frac{1}{2}\, \min \left\{1, \frac{m^{\frac{s-1}{s}}n^{\frac{(\frac{s}{t}-1)}{sm}}}{n^{\frac{s-1}{s}}m^{\frac{(\frac{s}{t}-1)}{sm}}}\right\}.
\end{equation}
It is not difficult  to see that the function $\psi: (0,\infty) \rightarrow \mathbb{R}$ defined by $$\psi(x)=x^{\frac{s-1}{s}} n^{\frac{(\frac{s}{t}-1)}{sm}},$$ attains a strict minimum at $x=\log \, n$ with the minimum value $$e^{\frac{(\frac{s}{t}-1)}{s}}\, (\log \, n)^{\frac{s-1}{s}}.$$ Finally, this fact together with Lemma \ref{lem-3.3} (1), \eqref{e-5.12}, and \eqref{e-5.13}, yields
\begin{equation} \label{e-5.15-added}
R_{\lambda}(B_Z,p,U) \geq D\,  \left(\frac{\lambda^p - \norm{U}^p}{2\lambda^p - \norm{U}^p}\right)^{\frac{1}{p}} \, \left(\frac{\log \, n}{n}\right)^{1-\frac{1}{s}}
\end{equation}
for some constant $D>0$. The lower bound of $AP_{\lambda}(B_{Z}, p,U)$ is a direct consequence Theorem \ref{thm-1.8-atm} and \eqref{e-5.15-added}.This completes the proof.
\comment{Observe that $1<t \leq s$, and so $(s-t)/mst \geq 0$. Using this fact, from \eqref{e-5.12}, we obtain 
	\begin{equation*}
		R^m_{\lambda}(B_Z,p,U) \geq\lambda^{\frac{1}{m}} \, \tilde{D} \left(1+\frac{n}{m}\right)^{-\frac{(s-1)}{s}}\, \left(1+\frac{n}{m}\right)^{\frac{(s-t)}{mst}} \geq \lambda^{\frac{1}{m}} \, \tilde{D}\, \left(1+\frac{n}{m}\right)^{-\frac{(s-1)}{s}}.
\end{equation*}}
\end{pf}






\begin{pf} [{\bf Proof of Theorem \ref{tthm-1.10-atm}}]
	Let $r=R_\lambda(\mathbb{D},p,U)$ and $f \in \mathcal{PH}(B_{z},X)$ be of the form \eqref{e-1.3-a}, where $X=\mathcal{B}(\mathcal{H})$.  Consider the function $\tilde{f}(z)=f(z, 0, \ldots,0)$ for $z \in \mathbb{D}$. Clearly, $\tilde{f}\in \mathcal{PH}(\mathbb{D},X)$ with
	the following expansion
	\begin{equation*}
		\tilde{f}(z)= \sum_{m=0}^{\infty} a_{(m,0,\ldots,0)} z^m + \sum_{m=1}^{\infty} b^{*}_{(m,0,\ldots,0)} \overline{z}^m, \,\,z \in \mathbb{D},
	\end{equation*}
	and $\norm{\tilde{f}}_{\mathbb{D},X}=\norm{f}_{B_{Z},X}$.
	Therefore, we have 
	\begin{align*}
		\sum_{m=0}^{\infty} \sum_{|\alpha|=m} (\norm{U(a_{\alpha})}^p_{Y} + \norm{U(b_{\alpha})}^p_{Y}) (r,0,\,.\,.\,.,0)^{p\alpha}
		&=\sum_{m=0}^{\infty}(\norm{U(a_{(m,0,\ldots,0)})}^p_{Y} + \norm{U(b_{(m,0,\ldots,0)})}^p_{Y}) r^{p\alpha}\nonumber,
	\end{align*}
	which is less than or equals to $\lambda^p \norm{\tilde{f}}^p_{\mathbb{D},X}=\lambda^p \norm{f}^p_{B_{Z},X}$.
	This shows that 
	$$\frac{r}{n}\leq AP_\lambda(B_{Z},p, U),$$
	 and hence we obtain $$\frac{R_\lambda(\mathbb{D},p,U)}{n} \leq AP_\lambda(B_{Z},p, U).
	 $$  Conversely, we prove that 
	\begin{equation*}
		AP_\lambda(B_{Z},p, U)\leq \norm{Id: Z \rightarrow \ell^n_{1}}\, \frac{R_\lambda(\mathbb{D},p,U)}{n^{1/p}}.
	\end{equation*}
	Suppose $r\in \mathbb{R}^n_{\geq 0}$ such that for all $f \in \mathcal{PH}(B_{z},X)$ be of the form \eqref{e-1.3-a},
	\begin{equation} \label{ee-1.8-atm}
		\sum_{m=0}^{\infty}	\sum_{|\alpha|=m} (\norm{U(a_{\alpha})}^p_{Y} + \norm{U(b_{\alpha})}^p_{Y})r^{p \alpha} \leq \lambda^{p} \norm{f}^p_{\Omega, X}.	
	\end{equation}
	We want to prove that 
	\begin{equation*}
		\frac{\sum_{j=1}^{n}r_j}{n}\leq \norm{Id: Z \rightarrow \ell^n_{1}}\, \frac{R_\lambda(\mathbb{D},p,U)}{n^{1/p}}.
	\end{equation*}
	Let $\tilde{f}:\mathbb{D}\rightarrow X \in \mathcal{PH}(\mathbb{D},X)$ be a bounded holomorphic map of the form $$\tilde{f}(\xi)=\sum_{m=0}^{\infty} c_{m} \xi^m + \sum_{m=1}^{\infty} d^{*}_m \overline{\xi}^m
	$$
	 for $\xi \in \mathbb{D}$. We now consider the function $\phi: B_{Z}\rightarrow \mathbb{C}$ defined by 
	\begin{equation} \label{e-function}
		\phi(z)=\frac{z_1+\cdots+z_n}{t},\quad z\in B_{Z},
	\end{equation}
	where $t= \norm{Id: Z \rightarrow \ell^n_{1}}$. Clearly, $\phi$ is holomorphic on $B_{Z}$ and $|\phi(z)|<1$ for all $z \in B_{Z}$.
	Now if we set $F=\tilde{f} \circ \phi $, then $F \in \mathcal{PH}(\Omega,X)$ with the following expansion 
	\begin{equation*}
		F(z)=\sum_{m=0}^{\infty} \sum_{|\alpha|=m} \left(\frac{c_{m}}{t^m} \frac{m!}{\alpha!}\right) z^{\alpha} +  \sum_{m=1}^{\infty} \sum_{|\alpha|=m} \left(\frac{d^{*}_{m}}{t^m} \frac{m!}{\alpha!}\right) \overline{z}^{\alpha},\,\, \,\, z \in B_{Z}
	\end{equation*}
	and $\norm{F}_{B_{Z},X}=\norm{\tilde{f}}_{\mathbb{D},X}$.
	By observing the fact $m!/\alpha! \geq 1$, for all $z\in B_{Z}$, we have
	\begin{align*}
		&\sum_{m=0}^{\infty} \sum_{|\alpha|=m} \left(\norm{U\left(\frac{c_{m}}{t^m} \frac{m!}{\alpha!}\right)}^p_{Y} + \norm{U\left(\frac{d_{m}}{t^m} \frac{m!}{\alpha!}\right)}^p_{Y}\right)|z|^{p \alpha} \nonumber \\
		& \geq \sum_{m=0}^{\infty} \sum_{|\alpha|=m} \left(\frac{m!}{\alpha!}\norm{U\left(\frac{c_{m}}{t^m} \right)}^p_{Y} + \frac{m!}{\alpha!}\norm{U\left(\frac{d_{m}}{t^m} \right)}^p_{Y}\right)|z|^{p \alpha} \nonumber \\
		& = \sum_{m=0}^{\infty} (\norm{U(c_{m})}^p_{Y} + \norm{U(d_{m})}^p_{Y})\, \left(\frac{\norm{z}_{p}}{t}\right)^{p m}.
	\end{align*}
	An observation, using preceding inequality and \eqref{ee-1.8-atm}, shows that
	\begin{align*}
		\sum_{m=0}^{\infty} (\norm{U(c_{m})}^p_{Y} + \norm{U(d_{m})}^p_{Y})\, \left(\frac{\norm{r}_{p}}{t}\right)^{p m}
		&\leq \sum_{m=0}^{\infty} \sum_{|\alpha|=m} \left(\norm{U\left(\frac{c_{m}}{t^m} \frac{m!}{\alpha!}\right)}^p_{Y} + \norm{U\left(\frac{d_{m}}{t^m} \frac{m!}{\alpha!}\right)}^p_{Y}\right)r^{p \alpha} \nonumber \\
		&\leq \lambda^p  \norm{F}^p_{B_{Z},X}=\lambda^p \norm{\tilde{f}}^p_{\mathbb{D},X},
	\end{align*}
	which implies that 
	$$
	\frac{\norm{r}_{p}}{t} \leq R_\lambda(\mathbb{D},p,U).
	$$ Again we have $\norm{r}_{1}\leq n^{(1-1/p)}\norm{r}_{p}$. Hence we obtain $$
	\frac{\norm{r}_{1}}{n}\leq \frac{t R_\lambda(\mathbb{D},p,U)}{n^{1/p}}.
	$$ 
	This completes the proof.
\end{pf}

\begin{pf} [{\bf Proof of Lemma \ref{lem-1.1}}]
Since $U$ is the identity operator on $\mathbb{C}$, we consider the complex valued harmonic functions on $\mathbb{D}$. Let $f:\mathbb{D} \rightarrow \mathbb{C}$ be a bounded harmonic function of the form
\begin{equation*}
	f(z)= \sum_{n=0}^{\infty} a_n\, z^n + \overline{\sum_{n=1}^{\infty} b_n\, z^n}, \,\,\, z \in \mathbb{D}.
\end{equation*}
Then, by \cite[Lemma 4]{Abu-2010}, we have 
\begin{equation*}
	|a_n|+ |b_n| \leq \frac{4}{\pi}\,\norm{f}_{\mathbb{D},\mathbb{C}}\,\,\,\,\, \mbox{for}\,\,\, n \geq 1.
\end{equation*}
By virtue of this estimate, for $|z|=r<1$, we obtain
\begin{align*}
|a_0|+ \sum_{n=1}^{\infty} (|a_n|+|b_n|)r^n 
& \leq |a_0|+ \frac{4}{\pi}\, \frac{r}{1-r}\,\norm{f}_{\mathbb{D},\mathbb{C}} \\
& \leq \norm{f}_{\mathbb{D},\mathbb{C}} + \frac{4}{\pi}\, \frac{r}{1-r}\,\norm{f}_{\mathbb{D},\mathbb{C}} \\
&= \left(1 + \frac{4}{\pi}\, \frac{r}{1-r} \right)\, \norm{f}_{\mathbb{D},\mathbb{C}} \leq \lambda \, \norm{f}_{\mathbb{D},\mathbb{C}},
\end{align*}
provided 
$$
r \leq \frac{(\lambda-1)\pi}{4+(\lambda-1)\pi}.
$$
Hence, the desired conclusion follows. The proof is now complete.
\comment{\\
	We now establish the part $(2)$. In this case, $U$ is the identity operator on $\mathcal{B}(\mathcal{H})$ and we consider the $\mathcal{B}(\mathcal{H})$ valued harmonic functions on $\mathbb{D}$. Let $f:\mathbb{D} \rightarrow \mathcal{B}(\mathcal{H})$ be a bounded harmonic function of the form
	\begin{equation*}
		f(z)= \sum_{n=0}^{\infty} a_n\, z^n + \sum_{n=1}^{\infty} b^{*}_n\, \overline{z}^n, \,\,\, z \in \mathbb{D}.
	\end{equation*}
	Then, by \cite[Lemma 4]{bhowmik-2021}, we have 
	\begin{equation*}
		||a_n||+ ||b_n|| \leq \frac{4}{\pi}\,\norm{f}_{\mathbb{D},\mathcal{B}(\mathcal{H})}\,\,\,\,\, \mbox{for}\,\,\, n \geq 1.
\end{equation*}}
\end{pf}



\noindent{\bf Statements and Declarations:}\\

\noindent{\bf Acknowledgment.} The author gratefully acknowledges the support received during the initial stages of this work while serving as a Senior Project Associate at IIT Bombay, through the \textit{Core Research Grant (CRG)}, sanctioned to Professor Sourav Pal by the Science and Engineering Research Board (SERB), Government of India. The author also sincerely acknowledges Professor B. V. Rajarama Bhat, ISI Bangalore, for support through the J. C. Bose Fellowship during the final stages of this manuscript. The author is currently supported by an NBHM Postdoctoral Fellowship from the Department of Atomic Energy (DAE), Government of India.\\

\noindent{\bf Conflict of interest.} The author declares that there is no conflict of interest regarding the publication of this paper.\\

\noindent{\bf Data availability statement.} Data sharing not applicable to this article as no datasets were generated or analysed during the current study.\\

\noindent{\bf Competing Interests.} The author declares none.\\
%\noindent\textbf{Acknowledgment:} 

\noindent{\bf ORCID Id.} 0000-0001-9466-
178X
%	The author gratefully acknowledges the support received during the initial stages of this work while serving as a Senior Project Associate at IIT Bombay, through the \textit{Core Research Grant (CRG)}, sanctioned to Professor Sourav Pal by the Science and Engineering Research Board (SERB), Government of India. The author is currently supported by the J. C. Bose Fellowship of Professor B. V. Rajarama Bhat, ISI Bangalore, whose support is also sincerely acknowledged.


\begin{thebibliography}{99}
	
	%	\bibitem{abouhajar-2007} {\sc A. Abouhajar, M. White,} and {\sc N. J. Young}, A Schwarz lemma for a domain related to $\mu$-synthesis, {\it J. Geom. Anal.} {\bf 17} (2007), 717--750.
	
	\bibitem{Abu-2010} {\sc Y. Abu-Muhanna},  Bohr's phenomenon in subordination and bounded harmonic classes, {\it Complex Var. Elliptic Equ.} {\bf  55} (2010), 1071--1078.
	
	
	%	\bibitem{abu-2011} {\sc Y. Abu-Muhanna} and  {\sc R. M. Ali}, Bohr's phenomenon for analytic functions into the exterior of a compact convex body, {\it J. Math. Anal. Appl.} {\bf  379}  (2011), 512--517.
	
	
	%	\bibitem{abu-2013} {\sc Y.  Abu Muhanna} and {\sc R. M.  Ali}, Bohr's phenomenon for analytic functions and the hyperbolic metric, {\it Math. Nachr.}  {\bf 286}  (2013), 1059--1065.
	
	
	%	\bibitem{abu-2014} {\sc Y. Abu Muhanna, R. M. Ali, Z. C. Ng} and  {\sc S. F. M Hasni}, Bohr radius for subordinating families of analytic functions and bounded harmonic mappings, 
	%	{\it J. Math. Anal. Appl.} {\bf 420} (2014), 124--136.
	
	%	\bibitem{Abu - Ali - Pon - 2017} {\sc Y.  Abu Muhanna}, {\sc R. M.  Ali} and {\sc S. Ponnusamy}, On the Bohr inequality, In ``Progress in Aporoximation Theory and Applicable Complex Analysis" (Edited by N.K. Govil et al.), \textit{Springer Optimization and its Applications}, \textbf{117} (2016), 265-295.
	
	%	\bibitem{Agler-Young-2004-JGEA} {\sc J. Agler} and {\sc N. J. Young}, The hyperbolic geometry of the symmetrized bidisc, {\it J. Geom. Anal.} {\bf 14} (2004), 375--403.
	
	%	\bibitem{Agler-Lykova-Young-2013-PLMS} {\sc J. Agler, Z. A. Lykova}, and {\sc N. J. Young}, Extremal holomorphic maps and the symmetrized bidisc, {\it Proc. London Math. Soc.} {\bf 106} (2013), 781--818.
	
	%	\bibitem{Agler-Lykova-Young-2015-JMMA} {\sc J. Agler, Z. A. Lykova}, and {\sc N. J. Young}, The complex geometry of a domain related to $\mu$-synthesis, {\it J. Math. Anal. Appl.} {\bf 422} (2015), 508--543.
	
	%	\bibitem{Ahamed-Allu-Halder-P3-2020} {\sc M. B. Ahamed, V. Allu} and {\sc H. Halder}, The Bohr Phenomenon for analytic functions on simply connected domains, {\it  Ann. Acad. Sci. Fenn. Math.} (2021), To appear.
	
	%	\bibitem{Ahamed-Allu-Halder-P3-2020} {\sc M. B. Ahamed, V. Allu} and {\sc H. Halder}, Bohr radius for certain classes of close-to-convex harmonic mappings, {\it Anal. Math. Phys.} {\bf 11} (111) (2021).
	
	%	\bibitem{aizn-1960} {\sc L. Aizenberg}, The spaces of functions analytic in $(p,q)$-circular regions, \textit{Soviet Math. Dokl.} {\bf 2} (1960), 79--82.
	
	\bibitem{aizn-2000a} {\sc L. Aizenberg}, Multidimensional analogues of Bohr's theorem on power series, \textit{Proc. Amer. Math. Soc.} {\bf 128} (2000), 1147--1155.
	
	\bibitem{aizn-2000b} {\sc L. Aizenberg, A. Aytuna}, and {\sc P. Djakov}, An abstract approach to Bohr's phenomenon, {\it Proc. Amer. Math. Soc.} {\bf 128} (2000), 2611--2619.
	
	
			\bibitem{aizenberg-2001a} {\sc L. Aizenberg, A. Aytuna},  and {\sc P. Djakov}, Generalization of theorem on Bohr for bases in spaces of holomorphic functions of several complex variables, 		{\it J. Math. Anal. Appl.} {\bf  258} (2001), 429--447.
	
			\bibitem{aizenberg-2001b} {\sc L. Aizenberg}  and {\sc N. Tarkhanov}, A Bohr Phenomenon for elliptic equations, {\it Proc. Lond. Math. Soc.} {\bf  82} (2001), 385--401.
	
	%		\bibitem{aizenberg-2003} {\sc L. Aizenberg, I. B. Grossman}  and {\sc Yu. F. Korobeinik}, Some remarks on the Bohr radius	for power series (Russian), {\it Izv. Vyssh. Uchebn. Zaved. Mat.}, 2002, no. 10, 3–10; {\it translation	in Russian Math. (Iz. VUZ)}, 46 (2002), no. 10, 1–8 (2003).
	
	%	\bibitem{aizn-2005} {\sc L. Aizenberg}, Generalization of Carath\'{e}odory’s inequality and the Bohr radius for multidimensional power series, Selected Topics in Complex Analysis. \textit{Operator Theory: Advances and Applications}, {\bf 158} (2005), Birkhäuser Basel, 87--94.
	
	
	%	\bibitem{aizn-2007} {\sc L. Aizenberg}, Generalization of results about the Bohr radius for power series, {\it Stud. Math.}  {\bf 180}  (2007), 161--168.  
	
	%	\bibitem{aizenberg-2012} {\sc L. Aizenberg}, Remarks on the Bohr and Rogosinski phenomena for power series, {\it Anal. Math. Phys.} {\bf 2} (2012), 69--78.
	
	%	\bibitem{albanese-2018} {\sc A. A. Albanese, J. Bonet}, and {\sc J. W. Ricker}, The Ces\'{a}ro operator on power series spaces, {\it Stud. Math.} {\bf 240} (2018), 47--68.
	
	%	\bibitem{Al-Kay-Pon-PAMS-2019} {\sc S. A. Alkhaleefah}, {\sc I. R. Kayumov} and {\sc S. Ponnusamy}, On the Bohr inequality with a fixed zero coefficients, {\it Proc. Amer. Math. Soc.} \textbf{147} (12) (2019), 5263--5274.
	
	%	\bibitem{Ali & Ng & CVEE & 2016} {\sc R. M. Ali} and {\sc Z. C. Ng}, The Bohr inequality in the hyperbolic plane, Complex Var. Elliptic Equ. \textbf{63}(11)(2018), 1539--1557.
	
	%\bibitem{Ali 1995} {\sc R. M .Ali}, {\sc S. Ponnusamy} and {\sc V. Singh}, Starlikeness of functions satisfying a differential inequality,  {\it Ann. Polon. Math.} {\bf  61} (1995), 135--140.
	
	%	\bibitem{Ali & Abdul & Ng & CVEE & 2016} {\sc R. M. Ali}, {\sc Z. Abdulhadi} and {\sc Z. C. Ng}, The Bohr radius for starlike logharmonic mappings, Complex Var. Elliptic Equ. \textbf{61}(1)(2016), 1--14.
	
	%	\bibitem{alshehri-2023} {\sc N. M. Alshehri} and {\sc Z. A. Lykova}, A Schwarz lemma for the Pentablock, {\it J. Geom. Anal.} {\bf 33} (2023), 44 pp.
	
	%	\bibitem{Ali-2017} {\sc R. M. Ali, R.W. Barnard} and {\sc  A.Yu. Solynin}, A note on Bohr's phenomenon for power series, {\it J. Math. Anal.Appl.} {\bf 449} (2017), 154-167.
	
	%	\bibitem{Ali-2019} {\sc R. M. Ali, N. K. Jain} and {\sc V. Ravichandran}, Bohr radius for classes of analytic functions, {\it Results Math.} {\bf 74} (2019) 179.
	
	%	\bibitem{alkhaleefah-2019} {\sc S. A. Alkhaleefah, I.R. Kayumov} and {\sc S. Ponnusamy}, On the Bohr inequality with a fixed zero coefficient, {\it Proc. Amer. Math. Soc.} {\bf 147} (2019), 5263--5274.
	
	%	\bibitem{Himadri-Vasu-P1} {\sc V. Allu} and {\sc H. Halder},  Bhor phenomenon for certain subclasses of Harmonic Mappings, {\it Bull. Sci. Math.} {\bf 173} (2021), 103053.
	
	%	\bibitem{Himadri-Vasu-P2} {\sc V. Allu} and {\sc H. Halder}, Bohr radius for certain classes of starlike and convex univalent functions, {\it J. Math. Anal. Appl.} {\bf 493}(1) (2021), 124519.
	
	%	\bibitem{Himadri-Vasu-P3} {\sc V. Allu} and {\sc H. Halder}, Bohr phenomenon for certain close-to-convex analytic functions, {\it Comput. Methods Funct. Theory} (2021),
	%	https://doi.org/10.1007/s40315-021-00412-6.
	
	%	\bibitem{Himadri-Vasu-P4} {\sc Vasudevarao Allu} and {\sc Himadri Halder}, Bohr inequality for certain harmonic mappings, see https://arxiv.org/pdf/2009.08683.pdf. 
	
	%	\bibitem{Himadri-Vasu-P6} {\sc V. Allu} and {\sc H. Halder}, Bohr phenomenon for operator valued functions on simply connected domains. 
	
	
	%	\bibitem{Himadri-Vasu-canad-2022} {\sc V. Allu } and {\sc H. Halder}, Operator valued analogues of multidimensional Bohr’s inequality, {\it Canad. Math. Bull.} {\bf 65} (2022), 1020--1035.
	
	%	\bibitem{Subhadip-Vasu-P1} {\sc V. Allu, H. Halder,} and {\sc S. Pal}, Bohr and Rogosinski inequalities for operator valued holomorphic functions, {\it Bull. Sci. Math.} {\bf 182} (2023), 103214, 18 pp.
	
	%	\bibitem{Himadri-Vasu-canad-2023} {\sc V. Allu } and {\sc H. Halder}, Bohr operator on operator-valued polyanalytic functions on simply connected domains {\it Canad. Math. Bull.}, (2023), 12 pp, doi:10.4153/S0008439523000541.
	
	\bibitem{Himadri-Vasu-PEMS} {\sc V. Allu } and {\sc H. Halder}, Bohr radius for banach spaces on simply connected domains, {\it Proc. Edinb. Math. Soc.} {\bf 67} (2024), 113--141.
	
	\bibitem{A-H-P-banach} {\sc V. Allu, H. Halder }, and {\sc S. Pal}, Multidimensional Bohr radii for holomorphic functions with values in complex Banach spaces, arXiv:2308.07825 (2023), 15 pp.
	
	%\bibitem{Allu-Pal-finite-banch} {\sc V. Allu} and {\sc S. Pal}, On multidimensional Bohr radius of finite dimensional Banach spaces, arXiv:2506.23540  (2025), 15 pp.
	
	\bibitem{Allu-Pal-arithnmetic} {\sc V. Allu} and {\sc S. Pal}, Arithmetic Bohr radius of bounded linear operators, (2025), 16 pp., arXiv:2410.18620v2.  
	
	%	\bibitem{Himadri-Vasu-CVEE-2025} {\sc V. Allu } and {\sc H. Halder}, Landau's theorem for harmonic Bloch functions on simply connected domains, {\it Complex Var. Elliptic Equ.} {\bf 70} (2025), 213--227.
	
	%	\bibitem{anderson-2006} {\sc J. M. Anderson} and {\sc J. Rovnyak}, On Generalized Schwarz-Pick Estimates, {\it Mathematika} {\bf 53} (2006), 161--168.
	
	
	%	\bibitem{bene-2004} {\sc C. B{\rm $\acute{E}$}n$\acute{E}$teau, A. Dahlner} and {\sc D. Khavinson}, Remarks on the Bohr phenomenon, {\it  Comput. Methods Funct. Theory}  {\bf 4}  (2004), 1--19.
	
	%   \bibitem{bharandar-2014} {\sc SV Bharanedhar} and {\sc S. Ponnusamy}, Coefficient conditions for harmonic univelent mappings and hypergeometric mappings, {\it Rocky Mountain J. Math.}	{\bf 44} (2014) 753--777.
	
	%	\bibitem{Ayt & Dja & BLMS & 2013} {\sc A. Aytuna} and {\sc P. Djakov}, Bohr property of bases in the space of entire functions and its generalizations, Bull. London Math. Soc. \textbf{45}(2)(2013), 411--420.
	\bibitem{baran-2026} {\sc R. Baran, P. Pikul, H. J. Woerdeman}, and {\sc M. Wojtylak}, Contractive realization theory for the annulus and other intersections of discs on the Riemann sphere, {\it J. Funct. Anal.} {\bf 290} (2026), 111346.
	
	\bibitem{bayart-advance-2014} {\sc F. Bayart, D. Pellegrino}, and {\sc J. B. Seoane-Sep$\rm \acute{U}$lveda}, The Bohr radius of the $n$-dimensional polydisk is equivalent to $\sqrt{(log \, n)/n}$, {\it Adv. Math.} {\bf 264} (2014), 726--746.
	
	%	\bibitem{beardon-minda-hyperbolic-density} {\sc A. F. Beardon} and {\sc D. Minda}, The hyperbolic metric and geometric function theory, {\it Proceedings of the International Workshop on Quasiconformal Mappings and their Applications (IWQCMA05)}.
	
		\bibitem{bala-2006} {\sc R. Balasubramanian}, {\sc B. Calado}, and {\sc H. Queffélec}, The Bohr inequality for ordinary Dirichlet series, {\it Studia Math.} {\bf 175} (2006), 285--304.
	
	\bibitem{bene-2004} {\sc C. B{\'e}n{\'e}teau}, {\sc A. Dahlner}, and {\sc D. Khavinson}, Remarks on the Bohr phenomenon, {\it Comput. Methods Funct. Theory} \textbf{4}(1) (2004), 1-19.
	
	%	\bibitem{Bermudez-2020} {\sc T. Berm{\' u}dez, A. Bonilla, V. M\"{u}ller}, and {\sc A. Peris}, Ces{\'a}ro bounded operators on Banach spaces, {\it J. Anal. Math.}
	%{\bf 140} (2020), 187--206.
	
	%	\bibitem{bernal-RACSAM-2021} {\sc L. Bernal-Gonz{\'a}lez, H. J. Cabana, D. Garc{\'\i}a, M. Maestre, G. A. Mu{\~n}oz-Fern{\'a}ndez}, and {\sc J. B. Seoane-Sep{\'u}lveda}, A new approach towards estimating the n-dimensional Bohr radius, {\it Rev. R. Acad. Cienc. Exactas Fís. Nat. Ser. A Mat. RACSAM} \textbf{115} (20021), 1-10.
	
	%	\bibitem{Bhattacharya-2012} {\sc T. Bhattacharyya, S. Pal}, and {\sc S. S. Roy}, Dilations of $\Gamma$-contractions by solving operator equations, {\it Adv. Math.} {\bf 230} (2012), 577--606.
	
	%	\bibitem{Bhattacharya-2014} {\sc T. Bhattacharyya} and {\sc S. Pal}, A functional model for pure $\Gamma$-contractions, {\it J. Operator Theory} {\bf 71} (2014), 327--339.
	
	\bibitem{bhowmik-2021} {\sc B. Bhowmik} and {\sc N. Das}, Bohr phenomenon for operator-valued functions, {\it Proc. Edinb. Math. Soc.} {\bf 64} (2021), 72--86.
	
	%	\bibitem{bhowmik-2018} {\sc B. Bhowmik} and {\sc N. Das},  Bohr phenomenon for subordinating families of certain univalent functions,  {\it J. Math. Anal. Appl.}  {\bf 462} (2018), 1087--1098.
	
	%	\bibitem{Blasco-OTAA-2010} {\sc O. Blasco}, The Bohr radius of a Banach space, {\it In Vector measures, integration and related topics, 5964, Oper. Theory Adv. Appl., 201, Birkh{\"a}user Verlag, Basel, 2010.}
	
	\bibitem{Blasco-Collect-2017} {\sc O. Blasco}, The $p$-Bohr radius of a Banach space, {\it Collect. Math.} {\bf 68} (2017), 87--100.
	
	
	\bibitem{boas-1997} {\sc H. P. Boas} and {\sc D. Khavinson}, Bohr's power series theorem in several variables, {\it Proc. Amer. Math. Soc.}  {\bf 125} (1997), 2975--2979.
	
	\bibitem{boas-2000} {\sc H. P. Boas}, Majorant Series,  {\it J. Korean Math. Soc.}  {\bf 37} (2000), 321--337.
	
	\bibitem{Bohr-1914} {\sc H. Bohr}, A theorem concerning power series,  {\it Proc. Lond. Math. Soc.} s2-13 (1914), 1--5.
	
	\bibitem{bombieri-1962} {\sc E. Bombieri}, Sopra un teorema di H. Bohr e G. Ricci sulle funzioni maggioranti delle serie di potenze, {\it Boll. Un. Mat. Ital.} {\bf 17} (1962), 276--282.
	
	\bibitem{bombieri-2004} {\sc E. Bombieri} and {\sc J. Bourgain}, A remark on Bohr's inequality, {\it Internat. Math. Res. Notices} {\bf 80} (2004), 4307--4330.
	
	%	\bibitem{Roth-schwarz-Adv-Math} {\sc F. Bracci, D. Kraus}, and {\sc O. Roth},  A new Schwarz-Pick lemma at the boundary and rigidity of holomorphic maps, {\it Adv. Math.} {\bf 432} (2023), Paper No. 109262, 41 pp.
	
	%	\bibitem{bshoty-2011} {\sc D. Bshouty} and {\sc A. Lyzzaik}, Close-to-convexity criteria for planar harmonic mappings, {\it Complex
		%		Appl. Oper. Theory} {\bf 5} (2011) 767--774.
	
	\bibitem{carando-2020} {\sc D. Carando, F. Marceca}, and {\sc P. Sevilla-Peris}, Hausdorff–Young-type inequalities for vector-valued Dirichlet series, {\it Trans. Amer. Math. Soc.} {\bf 373} (2020), 5627--5652.
	
%	\bibitem{chen-2011} {\sc S. Chen, S. Ponnusamy}, and {\sc X. Wang},  Bloch constant and Landau's theorems for planar $p$-harmonic mappings , {\it J. Math. Anal. Appl.} {\bf 373} (2011), 102-110.
	
	\bibitem{hamada-JFA-2022} {\sc S. Chen} and {\sc H. Hamada},  Some sharp Schwarz-Pick type estimates and their applications of harmonic and pluriharmonic functions, {\it J. Funct. Anal.} {\bf 282} (2022), Paper No. 109254, 42 pp.
	
		
	%	\bibitem{cho-2011} {\sc N. E. Cho}, {\sc O. S. Kwon} and {\sc V. Ravichandran}, Coefficient, distortion and growth inequalities for certain close-to-convex functions, {\it J. Inequal. Appl.} {\bf 2011:100} (2011), 7pp.
	
	%	\bibitem{chen-2011} {\sc S. Chen, S. Ponnusamy} and {\sc X. Wang}, Landau's theorem and Marden constant for harmonic $\nu$-Bloch mappings, {\it Bull. Aust. Math. Soc.} {\bf 84} (2011), 19--32.
	
	%	\bibitem{colona-1989} {\sc F. Colonna}, The Bloch constant of bounded harmonic mappings, {\it Indiana Univ. Math. J.} {\bf 38} (1989), 829--840.
	
	%	\bibitem{conway-functional-book-1990} {\sc J. B. Conway}, A course in functional analysis, Second edition. Graduate Texts in Mathematics, {\bf 96}. {\it Springer-Verlag, New York}, 1990.
	
	%	\bibitem{hamada-2022-JFA} {\sc S. Chen} and {\sc  H. Hamada}, Some sharp Schwarz-Pick type estimates and their applications of harmonic and pluriharmonic functions, {\it J. Funct. Anal.} {\bf 282} (2022), 109254, 42 pp.
	
	%		\bibitem{das-2023} {\sc N. Das}, Estimates for generalized Bohr radii in one and higher dimensions, {\it Canad. Math. Bull.} {\bf 66} (2023), 682--699.
	
	\bibitem{das-2024} {\sc N. Das}, The $p$-Bohr radius for vector-valued holomorphic and pluriharmonic functions, {\it Forum Math.} {\bf 36} (2024), 765--782.
	
	\bibitem{defant-2003} {\sc A. Defant, D. Garc\'{i}a}, and {\sc  M. Maestre}, Bohr power series theorem and local Banach space theory, {\it J. Reine Angew. Math.} {\bf 557} (2003), 173--197.
	
	\bibitem{defant-2004} {\sc A. Defant, D. Garc\'{i}a}, and {\sc  M. Maestre}, Estimates for the first and second Bohr radii of Reinhardt domains, {\it J. Appr. Theory} {\bf 128} (2004), 53--68.
	
	\bibitem{defant-2006} {\sc A. Defant} and {\sc L. Frerick}, A logarithmic lower bound for multi-dimenional Bohr radii, {\it Israel J. Math.} {\bf 152} (2006), 17--28.
	
	\bibitem{defant-2007} {\sc A. Defant,  M. Maestre}, and {\sc  C. Prengel}, The arithmetic Bohr radius, {\it Q. J. Math.} {\bf 59} (2008), 189--205.
	
			\bibitem{defant-2008} {\sc A. Defant, D. Garc\'{i}a, M. Maestre}, and {\sc D. P\'{e}rez-Garc\'{i}a}, Bohr's strip for vector-valued Dirichlet series, {\it Math. Ann.} {\bf 342} (2008), 533--555.
	
		\bibitem{defant-2009} {\sc A. Defant,  M. Maestre}, and {\sc  C. Prengel}, Domains of convergence for monomial expansions of holomorphic functions in infinitely many variables, {\it J. Reine Angew. Math.} {\bf 634} (2009), 13--49.
	
	\bibitem{defant-2011-lpn} {\sc A. Defant} and {\sc L. Frerick}, The Bohr radius of the unit ball of $\ell^n _p$, {\it J. Reine Angew. Math.} {\bf 660} (2011), 131--147.
	
	\bibitem{defant-2011} {\sc A. Defant, L. Frerick, J. Ortega-Cerd${\rm \grave{A}}$, M. Ouna{\"i}es}, and {\sc K. Seip}, The Bohnenblust-Hille inequality for homogeneous polynomils in hypercontractive, {\it Ann. of Math.} {\bf 174} (2011), 512--517.
	
	
	%			\bibitem{defant-2011} {\sc A. Defant, L. Frerick, J. Ortega-Cerd${\rm \grave{A}}$, M. Ouna{\"i}es}, and {\sc K. Seip}, The Bohnenblust-Hille inequality for homogeneous polynomils in hypercontractive, {\it Ann. of Math.} {\bf 174} (2011), 512--517.
	
	%		\bibitem{defant-book} {\sc A. Defant, D. García,  M. Maestre}, and {\sc P. Sevilla-Peris}, \textit{Dirichlet Series and Holomorphic Functions in High Dimensions} (New Mathematical Monographs) Cambridge: Cambridge University Press, (2019) doi:10.1017/9781108691611.
	
	
	\bibitem{defant-2012} {\sc A. Defant, M. Maestre}, and {\sc  U. Schwarting}, Bohr radii of vector valued holomorphic functions, {\it Adv. Math.} {\bf 231} (2012), 2837--2857.
	
	\bibitem{defant-2025-TAMS} {\sc A. Defant, D. Galicer, M. Mansilla, M. Mastyło}, and {\sc  S. Muro}, Local constants and Bohr’s phenomenon for Banach spaces of analytic polynomials, {\it Trans. Amer. Math. Soc. }(2025), DOI: https://doi.org/10.1090/tran/9643.
	
	%		\bibitem{Dixon & BLMS & 1995} {\sc P. G. Dixon}, Banach algebras satisfying the non-unital von Neumann inequality, \textit{Bull. Lond. Math. Soc.} \textbf{27} (4) (1995), 359--362.
	
	\bibitem{diestel-abs-summing-1995} {\sc J. Diestel, H. Jarchow}, and {\sc A. Tonge}, {\it Absolutely summing operators}, Cambridge Stud. Adv. Math. 43, Cambridge University, Cambridge, 1995.
	
	
	\bibitem{Djakov & Ramanujan & J. Anal & 2000} {\sc P. B. Djakov} and {\sc M. S. Ramanujan}, A remark on Bohr's theorem and its generalizations, \textit{J. Anal.} \textbf{8} (2000), 65--77.
	
	\bibitem{Dineen-Timoney-1989} {\sc S. Dineen} and {\sc R. M. Timoney}, Absolute bases, tensor products and a theorem of Bohr, \textit{Studia Math.} \textbf{94} (1989), 227--234.
	
	%	\bibitem{duren-1983} {\sc P. L. Duren}, {\it Univalent Functions} (Grundlehren der mathematischen Wisseenschaften 259, New York, Berlin, Heidelberg, Tokyo) Springer-Verlag, 1983.
	
	
	%	\bibitem{efraimidis-2017} {\sc I. Efraimidis, J. Gaona, R. Hern${\rm \acute{A}}$ndez} and {\sc O. Venegas}, On harmonic Bloch-type mappings, {\it Complex Var. Elliptic Equ.} {\bf 62} (2017), 1081--1092.
	
	%	\bibitem{edigarian-2005} {\sc A. Edigarian} and {\sc W. Zwonek}, Geometry of the symmetrized polydisc, {\it Arch. Math. (Basel)} {\bf 84} (2005), 364--374.
	
	%	\bibitem{evdoridis-2019} {\sc S. Evdoridis} and {\sc S. Ponnusamy}, Improved Bohr's inequality for locally univalent harmonic mappings, {\it  Indag. Math. (N.S.) } {\bf 30} (2019), 201--213.
	
	
	%	\bibitem{Evdoridis-Ponnusamy-2018} {\sc S. Evdoridis}, {\sc S. Ponnusamy} and {\sc A. Rasila}, Improved Bohr's inequality for locally univalent harmonic mappings, {\it Indag. Math. (N.S.)} {\bf 30} (2019), no. 1, 201–-213.
	
	%	\bibitem{Evd-Ponn-Rasi-2020} {\sc S. Evdoridis}, {\sc S. Ponnusamy}, and {A. Rasila}, Improved Bohr's inequality for shifted disks, {\it Results Math.} {\bf 76:14} (2021), 15 pages.
	
	%	\bibitem{Bohr-Blaschke-Book} {\sc S. R. Garcia, J. Mashreghi}, and W. T. Ross, {\it Finite Blaschke products and their connections}, Springer, Cham, 2018.
	
	%	\bibitem{goodman-1991-a} {\sc A. W. Goodman}, On uniformly convex functions, {\it Annl. Polon. Math.}, {\bf 56} (1991), 87--92.
	
	%	\bibitem{goodman-1991-b} {\sc A. W. Goodman}, On uniformly starlike functions, {\it J. Math. Anal. Appl.}, {\bf 155} (1991), 364--370.
	
	%	\bibitem{Four-Rusc-2010} {\sc R. Fournier} and {\sc St. Ruscheweyh}, On the Bohr radius for simply connected domains, \textit{Centre de Recherches Math$ \acute{e} $matiques CRM Proceedings and Lecture Notes}, Vol. \textbf{51} (2010), 165--171.
	
	%	\bibitem{gao-2005} {\sc C. Gao} and {\sc S. Zhou}, On a class of analytic functions related to the starlike functions, {\it Kyungpook Math. J.} {\bf 45} (2005), 123--130.
	
	%	\bibitem{gargi-2023} {\sc G. Ghosh} and {\sc W. Zwonek}, $2$-proper holomorphic images of classical Cartan domains, {\it Indiana Univ. Math. J.} (2024), to appear.
	
	%\bibitem{graham-1996} {\sc I. Graham} and {\sc D. Varolin}, Bloch constants in one and several variables, {\it Pacific J. Math.} {\bf 174} (1996), 347--357.	
	
	%	\bibitem{graham-2003} {\sc I. Graham} and {\sc G. Kohr}, Geometric Function Theory in One and Higher Dimensions, Monographs and Textbooks in Pure and Applied Mathematics, {\bf 255}, {\it Marcel Dekker, Inc., New York}, 2003. 
	
	%	\bibitem{nirupam cvee} {\sc Nirupam  Ghosh} and {\sc A. Vasudevarao}, Some basic properties of certain subclass of harmonic univalent functions, {\it Complex Var. Elliptic Equ.}  {\bf 63} (2018), 1687--1703.
	
	%\bibitem{nirupam-2019} {\sc Nirupam Ghosh} and {\sc A.Vasudevarao}, On a subclass of harmonic close-to-convex mappings, {\it Monatsh. Math.} {\bf 188} (2019), 247-267.
	
	
	%	\bibitem{nirupam-bull aus math} {\sc Nirupam Ghosh} and {\sc  A. Vasudevarao}, On some subclasses of harmonic mappings, {\it Bull. Aust. Math. Soc.} {\bf 101} (2020), 130--140.
	
	%	\bibitem{Hamada-ISJM-2009} {\sc H. Hamada, T. Honda}, and {\sc G. Kohr}, Bohr's theorem for holomorphic mappings with values in homogeneous balls, {\it Israel J. Math.} {\bf 173} (2009), 177--187.
	
	\bibitem{hamada-Math-Nachr-2023} {\sc H. Hamada},  Bohr's inequality for holomorphic and pluriharmonic mappings with values in complex Hilbert spaces, {\it Math. Nachr.} {\bf 296} (2023), 2795-2808.
	
	%	\bibitem{Hardy-1931} {\sc G. H. Hardy} and {\sc J. E. Littlewood}, Some properties of fractional integrals II, {\it Math. Z.} {\bf 34} (1931), 403--439.
	
	%   \bibitem{harnandez-2015} {\sc R. Hern${\rm \acute{A}}$ndez} and {\sc M. J. Mart${\rm \acute{I}}$n}, Pre-Schwarzian and Schwarzian derivatives of harmonic mappings, {\it J. Geom. Anal.} {\bf 25} (2015), 64--91.
	
	\bibitem{Himadri-local-Banach-1} {\sc H. Halder}, Local Banach space theoretic approach to Bohr's theorem for vector valued holomorphic and pluriharmonic functions, (2025), 19pp., arXiv:2512.09091.
	
	%	\bibitem{Huang-Liu-Ponnu} {\sc Y. Huang, M-S. Liu}, and {\sc S. Ponnusamy}, Refined Bohr-type inequalities with area measure for bounded analytic functions, {\it Anal. Math. Phys.} {\bf 10} (2020), 50. 
	
	%	\bibitem{Ismagilov-2020} {\sc A. Ismagilov, I. R. Kayumov} and {\sc S. Ponnusamy}, Sharp Bohr type inequality, 
	%	{\it J. Math. Anal. Appl.}  {\bf 489} (2020), 124147.
	
	%	\bibitem{Jarnici-Pflug-book} {\sc M. Jarnicki} and {\sc P. Pflug}, Invariant Distances and Metrics in Complex Analysis, 2nd ed., de Gruyter	Expos. in Math., {\bf 9}, Walter de Gruyter GmbH, Berlin, 2013. 
	
	%	\bibitem{kargar-2019} {\sc R. Kargar}, {\sc A. Ebadian} and {\sc J. Sokol}, On Both lemniscate and starlike functions, {\it Anal. Math. Phys.} {\bf 9} (2019), 143--154. 
	
	%	\bibitem{kanas-2019} {\sc S. Kanas}, {\sc V. S. Masih} and {\sc A. Ebadian}, Relations of a planar domains bounded by hyperbolas with families of holomorphic functions, {\it J. Inequal. Appl.} {\bf 246} (2019), 1--14.
	
	%	\bibitem{kayumov-2017} {\sc I.R. Kayumov} and {\sc S. Ponnusamy}, Bohr-Rogosinski radius for analytic functions, 
	%	preprint, see https://arxiv.org/abs/1708.05585.
	
	
	%	\bibitem{kayumov-2018-a} {\sc I.R. Kayumov, S. Ponnusamy} and {\sc N. Shakirov}, Bohr radius for localy unovalent harmonic mappings, {\it  Math. Nachr.} {\bf  291} (2018), 1757--1768.
	
	%	\bibitem{Kayumov-Ponnusamy-2018-b} {\sc I.R. Kayumov} and {\sc S. Ponnusamy}, Bohr's inequalities for the analytic functions with lacunary series and harmonic functions, {\it J. Math. Anal. Appl.}  {\bf 465} (2018), 857--871.
	
	%	\bibitem{Kay & Pon & AASFM & 2019} {\sc I.R. Kayumov} and {\sc S. Ponnusamy}, On a powered Bohr inequality, \textit{Ann. Acad. Sci. Fenn. Ser. A}, \textbf{44}(2019), 301--310.
	
	%	\bibitem{kayumov-2018-c} {\sc I. R. Kayumov} and {\sc S. Ponnusamy}, Improved version of Bohr's inequalities, 
	%	{\it C. R. Math. Acad. Sci. Paris} {\bf 358} (5) (2020), 615--620.
	
	%	\bibitem{kayumov-2019} {\sc I.R. Kayumov} and {\sc S. Ponnusamy}, On a powered Bohr inequality, {\it  Ann. Acad. Sci. Fenn. Math.} {\bf 44} (2019), 301--310.
	
	%	\bibitem{kayumova-2018} {\sc A. Kayumova, I. R. Kayumov} and {\sc S. Ponnusamy}, Bohr's inequality for harmonic mappings and beyond,
	%	{\it Mathematics and Computing}, 245--256, Commun. Comput. Inf. Sci., 834, Springer, Singapore, 2018.	
	
	%	\bibitem{Kayumov-Ponnuswamy-MN-2018} {\sc I. R. Kayumov}, {\sc S. Ponnusamy} and {\sc N. Shakirov}, Bohr radius for locally univalent harmonic mappings, {\it Math. Nachr} {\bf 291} (2018), 1757--1768.
	
	%	\bibitem{Kayumov-Ponnuswamy-Cesaro-CRM-2020} {\sc I. R. Kayumov}, {\sc D. M. Khammatova } and {\sc S. Ponnusamy}, On the Bohr inequality for the Ces{\'a}aro operator, {\it C. R. Math. Acad. Sci. Paris} {\bf 358} (2020), 615--620.
	
	%	\bibitem{Kayumov-Ponnusamy-Cesaro-average-arxiv-2021} {\sc I. R. Kayumov}, {\sc D. M. Khammatova } and {\sc S. Ponnusamy}, The Bohr inequality for the generalized Ces{\'a}aro averaging operators, https://arxiv.org/abs/2104.01550.
	
	
	%	\bibitem{khatter-2019} {\sc K.Khatter}, {\sc V. Ravichandran} and {\sc S. Sivaprasad Kumar}, Starlike functions associated with exponential function and the lemniscate of Bernoulli, 
	%	{\it Rev. R. Acad. Cienc. Exactas F$\acute{i}$s. Nat. Ser. A Mat. RACSAM} {\bf 113} (2019), 233--253.
	
	%	\bibitem{Roth-Scwarw-Canad} {\sc M. Kourou} and {\sc O. Roth }, Geometric versions of Schwarz's lemma for spherically convex functions, {\it Canad. J. Math.} {\bf 75} (2023), 1780–1799.
	
	%	\bibitem{kowalczyk-2010} {\sc J. Kowalczyk} and {\sc E. Les-Bomba}, On a subclass of close-to-convex functions, {\it Appl. Math. Lett.}  {\bf 23} (2010), 1147--1151.
	
	%	\bibitem{kumar-2016} {\sc S. Kumar} and {\sc V. Ravichandran}, A subclass of starlike functions associated with a rational function, {\it Southeast Asian Bull. Math.} {\bf 40} (2016), 199--212.
	
	%	\bibitem{Kosinski-2011} {\sc \L. Kosi\'{n}ski}, Geometry of quasi-circular domains and applications to tetrablock, {\it Proc. Amer. Math. Soc.} {\bf 139} (2011), 559-569. 
	
	%	\bibitem{Kosinski-2016} {\sc \L. Kosi\'{n}ski} and {\sc W. Zwonek}, Nevanlinna-Pick problem and uniqueness of left inverses in convex domains,	symmetrized bidisc and tetrablock, {\it J. Geom. Anal.} {\bf 26} (2016), 1863--1890. 
	
	%		\bibitem{kumar-2023} {\sc S. Kumar}, On the multidimensional Bohr radius, {\it Proc. Amer. Math. Soc.} {\bf 151} (2023), 2001--2009.
	\bibitem{knese-2025} {\sc G. Knese}, Three radii associated to Schur functions on the polydisk, {\it Proc. Amer. Math. Soc. Ser. B} {\bf 12} (2025), 48--63.
	
	\bibitem{Lindenstrauss-book} {\sc J. Lindenstrauss} and {\sc L. Tzafriri}, Classical Banach spaces $I$ and $II$, Springer-Verlag, 1977, 1979.
	
			\bibitem{lassere-2017} {\sc P. Lassère} and {E. Mazzilli}, Estimates for the Bohr Radius of a Faber–Green Condenser in the Complex Plane, {\it Constr. Approx.} {\bf 45} (2017), 409--426.
	
	%	\bibitem{landau-1986} {\sc E. Landau} and {\sc D. Gaier}, 
	%	\bibitem{sakaguchi-1959} {\sc K. Sakaguchi}, On a certain univalent mapping, {\it J. Math. Soc. Japan} {\bf 11} (1959), 72--75.
	
	%	\bibitem{sivaprasadkumar-2020} {\sc S. Sivaprasad Kumar} and {\sc S. Banga}, On a special type of Non-Ma-Minda function, arXiv:2006.02111v1, 2020.
	
	%	\bibitem{Y. Sun} {\sc Y. Sun}, {\sc Y-P Jiyang} and {\sc A. Rasila}, On a certain subclass of close-to-convex harmonic mappings, {\it Complex Var. Elliptic Equ.} {\bf 61} (2016) 1627-1643.
	
	%	\bibitem{Injectivity section} {\sc L. Li} and {\sc S. Ponnusamy}, Injectivity of sections of univalent harmonic mappings, {\it Nonlinear Analysis} {\bf  89} (2013), 276--283.
	
	%	\bibitem{Liu-Indag-Math-2018} {\sc G. Liu} and {\sc S. Ponnusamy}, Uniformly locally univalent harmonic mappings associated with the pre-Schwarzian norm, {\it Indag. Math. (N. S.)} {\bf 29} (2018), 752--778.
	
	%	\bibitem{Liu-2019} {\sc Z. Liu} and {\sc S. Ponnusamy}, Bohr radius for subordination and $k$-quasiconformal harmonic mappings, {\it  Bull. Malays. Math. Sci. Soc.} {\bf 42} (2019), 2151--2168.
	
	%	\bibitem{Liu-2020} {\sc G. Liu}, {\sc Z. Liu} and {\sc S. Ponnusamy}, Refined Bohr inequality for bounded analytic functions, priprint, see https://arxiv.org/pdf/2006.08930.
	
	%	\bibitem{Liu-Ponnusamy-BMMS-2019} {\sc Z. H. Liu} and {\sc S. Ponnusamy}, Bohr radius for subordination and $ K $-quasiconformal harmonic mappings, Bull. Malys. Math. Sci. Soc. \textbf{42} (2019), 2151--2168.
	
	%	\bibitem{Liu-Ponnusamy-Wang-2020} {\sc M. S. Liu, S. Ponnusamy} and {\sc J. Wang}, Bohr's phenomenon for the classes of Quasi-subordination and $ K $-quasiregular harmonic mappings, Rev. R. Acad. Cienc. Exactas Fs. Nat. Ser. A Mat. RACSAM  (2020) 114: 115, 15 pp.
	
	%	\bibitem{Liu-Pon-PAMS-2020} {\sc M. S. Liu} and {\sc S. Ponnusamy}, Multidimensional analogues of refined Bohr's inequality, {\it Proc. Amer. Math. Soc.} {\bf 149} (2021), 2133--2146.
	
	%	\bibitem{lyzzaik-1984} {\sc A. Lyzzail}, On a conjecture of M. S. Robertson, {\it Proc. Amer. Math. Soc.} {\bf 91} (1984), 108--110. 
	
	%	\bibitem{ma minda-1992-a} {\sc  W.C. Ma} and {\sc  D. Minda}, A unified treatment of some special classes of univalent functions, in {\it Proceedings of the Conference on Complex Analysis}(Tianjin, 1992), 157--169, Conf. Proc. Lecture Notes Anal., I, Int. Press, Cambridge.
	
	%	\bibitem{ma minda-1992-b} {\sc  W.C. Ma} and {\sc  D. Minda}, Uniformly convex functions, {\it Ann. Polon. Math.} {\bf 57} (1992), 165--175.
	
	%	\bibitem{ma minda-1993} {\sc  W.C. Ma} and {\sc  D. Minda}, Uniformly convex functions II, {\it Ann. Polon. Math.} {\bf 58} (1993), 275--285.
	
	%	\bibitem{masih-2020} {\sc V. S. Masih} and {\sc S. Kanas}, Subclasses of Starlike and Convex functions associated with associated with the Limaçon Domain, {\it Symmetry} {\bf 12} (2020), 942. 
	%	\bibitem{mendiratta-2015} {\sc R. Mendiratta}, {\sc S. Nagpal} and {\sc V. Ravichandran}, On a subclass of strongly starlike functions associated with exponential function, {\it Bull. Malays. Math. Sci. Soc.} {\bf 38} (2015), 365--386.
	
	%	\bibitem{Miller-Mocanu-diff-book} {\sc S. S. Miller} and {\sc P. T. Mocanu}, Differential Subordinations-Theory and Applications, {\it Marcel Dekker, Inc., New York}, 2000.
	
	%	\bibitem{pal-2014-a} {\sc S. Pal}, From Stinespring dilation to Sz.-Nagy dilation on the symmetrized bidisc and operator models, {\it New York J. Math.} {\bf 20} (2014), 545--564.
	
	%	\bibitem{pal-2014-b} {\sc S. Pal} and {\sc O. M. Shalit}, Spectral sets and distinguished varieties in the symmetrized bidisc, {\it J Funct. Anal.} {\bf 266} (2014), 5779-5800.
	
	%	\bibitem{pal-2015} {\sc S. Pal}, The failure of rational dilation on the tetrablock, {\it J. Funct. Anal.} {\bf 269} (2015), 1903-1924.
	
	%	\bibitem{pal-2018-a} {\sc S. Pal}, Canonical decomposition of operators associated with the symmetrized polydisc, {\it Complex Anal. Oper. Theory} {\bf 12} (2018), 931--943.
	
	%		\bibitem{pal-2018-b} {\sc S. Pal} and {\sc S. Roy}, A generalized Schwarz lemma for two domains related to $\mu$-synthesis, {\it Complex Manifolds} {\bf 5} (2018), 1--8.
	
	%	\bibitem{pal-2021-a} {\sc S. Pal} and {\sc S. Roy}, A Schwarz lemma for two families of domains and complex geometry, {\it Complex Var. Elliptic Equ.} {\bf 66} (2021), 756--782.
	
	%	\bibitem{pal-2021-b} {\sc S. Pal} and {\sc S. Roy}, Characterizations of the symmetrized polydisc via another family of domains, {\it Internt. J. Math.} {\bf 32} (2021), 2150036.
	
		\bibitem{paulsen-2002} {\sc V. I. Paulsen, G. Popescu}, and {\sc D. Singh}, On Bohr's inequality, {\it Proc. Lond. Math. Soc.} s3-85 (2002), 493--512.
	
			\bibitem{paulsen-2004} {\sc V. I. Paulsen} and {\sc D. Singh}, Bohr’s inequality for uniform algebras, {\it Proc. Amer. Math. Soc.} {\bf 132} (2004), 3577-3579.
	
			\bibitem{prengel-2005} {\sc C. Prengel}, {\it Domains of convergence in infinite dimensional holomorphy}, Ph.D. Thesis, University of Oldenburg, 2005.
	
	%	\bibitem{Ponnusamy-Vasu-Vuorinen-2009}	{\sc S. Ponnusamy,  A. Vasudevarao} and {\sc  M. Vuorinen}, Region of variability for spirallike functions with respect to a boundary point, 
	%	{\it Colloq. Math.} {\bf 116} (2009), 31--46.
	
	%	\bibitem{S.Ponnusamy variability regions-2013} {\sc S. Ponnusamy, H. Yamamoto} and {\sc H. Yanagihara}, Variability regions for certain families of harmonic univalent mappings,
	%	{\it  Complex Var. Elliptic Equ.} {\bf  58} (1) (2013), 23--34.
	
	%	\bibitem{Ponnusamy-Wirths-2020} {\sc S. Ponnusamy} and {\sc K-J. Wirths}, Bohr type inequalities for functions with a multiple zero at the origin, https://arxiv.org/pdf/2006.06441.pdf.
	
		\bibitem{popescu-2019} {\sc G. Popescu}, Bohr inequalities for free holomorphic functions on polyballs, {\it Adv. Math.} {\bf 347} (2019), 1002-1053.
	
	
	%\bibitem{raina-2015} {\sc R. K.Raina} and {\sc J. Sokol}, Some properties related to a certain class of starlike functions, {\it C. R. Math. Acad. Sci. Paris} {\bf 353} (2015), 973--978.
	
	%		\bibitem{Range-book} {\sc M. R. Range}, Holomorphic functions and integral representations in several complex variables, {\it Graduate Texts in Mathematics}, {\bf 108}, Springer-Verlag, New York, 1986.
	
	%\bibitem{Ravichandran-2004} {\sc V. Ravichandran}, Starlike and convex functions with respect to conjugte points, {\it Acta Math. Acad. Paedagog. Nyhazi.} {\bf 20} (2004), 31--37.
	
	%	\bibitem{robertson-1981} {M. S. Robertson}, Univalent functions with respect to a boundary point, {\it J. Math. Anal. Appl.} {\bf 81} (1981), 327--345.
	
		\bibitem{Rocchetta-Slice-Bohr-Math-Nachr} {\sc C. D. Rocchetta, G. Gentili}, and {\sc G. Sarfatti}, The Bohr theorem for slice regular functions, {\it Math. Nachr.} \textbf{285} (2012), 2093–-2105.
	
	%\bibitem{rogosinski-1930} {\sc W. Rogosinski}, Uber Bildschranken bei Potenzreihen und ihren Abschnitten, {\it Math. Z.} {\bf 17} (1923), 260-276.
	
	%		\bibitem{rong-2022} {\sc F. Rong} and {\sc S. Yang}, On Fridman invariants and generalized squeezing functions, {\it Chin. Ann. Math. Ser. B} {\bf 43} (2022), 161--174.
	
	%	\bibitem{ronning-1991} {\sc F. Ronning}, On starlike functions associated with parabolic regions, {\it Anna. Univ. Mariae Curie-Sklodowska Sect. A}, {\bf 45} (1991), 117--122.
	
	%\bibitem{ronning-1993} {\sc F. Ronning}, Uniformly convex functions and a corresponding class of starlike functions, {\it Proc. Amer. Math. Soc.} {\bf 118} (1993), 189--196.
	
	%	\bibitem{sharma-2016} {\sc K. Sharma}, {\sc N. K. Jain} and {\sc V. Ravichandran}, Starlike functions associated with cardiod, {\it Afr. Mat.} {\bf 27} (2016), 923--939.
	
	%	\bibitem{sidon-1927} {\sc S. Sidon}, Uber einen satz von Hernn Bohr, {\it Math. Zeit.}  {\bf 26} (1927),  731-732.
	
	%\bibitem{silverman-1990} {\sc H. Silverman} and {\sc E.M. Silvia}, Subclasses of univalent functions starlike with respect to a boundary point, {\it Houston J. Math.} {\bf 16} (1990), 289--299. 
	
	%	\bibitem{rogosinski-1943} {\sc W. Rogosinski}, On the coefficients of subordinate functions, {\it Proc. London Math. Soc.} {\bf 48} (1943), 48--82.
	
	%	\bibitem{Ruscheweyh-1985} {\sc St. Ruscheweyh}, Two remarks on bounded analytic functions, \emph{Serdica} \textbf{11}(1) (1985), 731--732.
	
	
	%	\bibitem{stempak-1994} {\sc K. Stempak}, Ces{\'a}ro averaging operators, {\it Proc. R. Soc. Edinb., Sect. A, Math.} {\bf 124} (1994), 121--126.
	
	%	\bibitem{Y. Sun} {\sc Y. Sun}, {\sc Y-P Jiyang} and {\sc A. Rasila}, On a certain subclass of close-to-convex harmonic mappings, {\it Complex Var. Elliptic Equ.} {\bf 61} (2016) 1627-1643.
	
	%	\bibitem{vasu-book} {\sc Derek K. Thomas}, {\sc Nikola Tuneski} and {\sc Allu Vasudevarao}, {\it Univalent functions}. A primer, De Gruyter Studies in Mathematics, {\bf 69}. De Gruyter, Berlin, 2018.
	
	%	\bibitem{tomic-1962} {\sc M. Tomic}, Sur un theoreme de H. Bohr,  {\it Math. Scand.}  {\bf 11} (1962), 103--106.
	
	\bibitem{Xu-quaternion-AMPA} {\sc Z. Xu}, Generalized Bohr radius for slice regular functions over quaternions,  {\it Ann. Mat. Pura
		Appl.}  {\bf 200} (2021), 1217--1229.
	
	\bibitem{Xu-quaternion-PRSESA} {\sc Z. Xu}, Bohr theorems for slice regular functions over octonions,  {\it Proc. Roy. Soc. Edinburgh Sect. A}  {\bf 151} (2021), 1595--1610.
	
	%	\bibitem{Zapalowski-2008} {\sc P. Zapa\l{}owski}, Geometry of symmetrized ellipsoids, {\it Univ. Iagel. Acta Math.}  {\bf 46} (2008), 105--116.
	
	%	\bibitem{Zwonek-Arch-Math-2013} {\sc W. Zwonek}, Geometric properties of the tetrablock, {\it Arch. Math. (Basel)} {\bf 2013}, 159-165.
\end{thebibliography}
\end{document}
